\documentclass[11pt]{article}

\usepackage[utf8]{inputenc}
\usepackage[T1]{fontenc}
\usepackage[USenglish]{babel}
\usepackage[a4paper,margin=32mm,headheight=14pt]{geometry}
\usepackage{lmodern}
\usepackage{microtype}
\usepackage{csquotes}

\usepackage{amsmath,amssymb,amscd}
\usepackage{amsthm}
\usepackage{mathtools}
\usepackage{mathrsfs}
\usepackage{enumitem}
\usepackage{aliascnt}
\usepackage{xcolor}
\usepackage{fancyhdr}
\usepackage[
  colorlinks=true,
  linkcolor=blue!45!black,
  citecolor=green!35!black,
  urlcolor=blue!55!black,
  pdfauthor={Alex Junior Gomez Saltachin},
  pdftitle={Derived categories and Brauer groups of log del Pezzo surfaces with 1/3(1,1) singularities},
  pdfsubject={Algebraic geometry},
  pdfkeywords={Log del Pezzo surfaces, canonical stacks, exceptional collections, semiorthogonal decompositions, quotient singularities, Brauer groups, special McKay correspondence}
]{hyperref}
\usepackage[capitalize,nameinlink]{cleveref}

% Theorem environments with a common counter.
\theoremstyle{plain}
\newtheorem{theorem}{Theorem}[section]
\newaliascnt{proposition}{theorem}
\newtheorem{proposition}[proposition]{Proposition}
\aliascntresetthe{proposition}
\newaliascnt{lemma}{theorem}
\newtheorem{lemma}[lemma]{Lemma}
\aliascntresetthe{lemma}
\newaliascnt{corollary}{theorem}
\newtheorem{corollary}[corollary]{Corollary}
\aliascntresetthe{corollary}

\theoremstyle{definition}
\newaliascnt{definition}{theorem}
\newtheorem{definition}[definition]{Definition}
\aliascntresetthe{definition}
\newaliascnt{example}{theorem}
\newtheorem{example}[example]{Example}
\aliascntresetthe{example}
\theoremstyle{remark}
\newaliascnt{remark}{theorem}
\newtheorem{remark}[remark]{Remark}
\aliascntresetthe{remark}

\crefname{theorem}{Theorem}{Theorems}
\crefname{proposition}{Proposition}{Propositions}
\crefname{lemma}{Lemma}{Lemmas}
\crefname{corollary}{Corollary}{Corollaries}
\crefname{definition}{Definition}{Definitions}
\crefname{example}{Example}{Examples}
\crefname{remark}{Remark}{Remarks}

% Macros
\newcommand{\A}{\mathbb A}
\newcommand{\Pp}{\mathbb P}
\newcommand{\F}{\mathbb F}
\newcommand{\C}{\mathbb C}
\newcommand{\Z}{\mathbb Z}
\newcommand{\Db}{D^b}
\newcommand{\coh}{\operatorname{coh}}
\newcommand{\Sing}{\operatorname{Sing}}
\newcommand{\Perf}{\operatorname{Perf}}
\newcommand{\Pic}{\operatorname{Pic}}
\newcommand{\rk}{\operatorname{rk}}
\newcommand{\Hom}{\operatorname{Hom}}
\newcommand{\Ext}{\operatorname{Ext}}
\newcommand{\Bl}{\operatorname{Bl}}
\newcommand{\GL}{\mathrm{GL}}
\newcommand{\Cl}{\mathrm{Cl}}
\newcommand{\Br}{\mathrm{Br}}
\newcommand{\Spec}{\mathrm{Spec}}
\newcommand{\Bs}{\operatorname{Bs}}
\newcommand{\RHom}{\operatorname{RHom}}
\newcommand{\XX}{\mathcal X}
\newcommand{\Xten}{X_{10}}
\newcommand{\XXten}{\mathcal X_{10}}
\newcommand{\Xtilden}{\widetilde X}
\newcommand{\Xtenten}{\widetilde X_{10}}
\newcommand{\muthree}{\boldsymbol{\mu}_3}
\newcommand{\OO}{\mathcal O}

\usepackage[
  backend=biber,
  style=numeric,
  sorting=nyt,
  giveninits=true,
  doi=true,
  url=false,
  isbn=false
]{biblatex}
\addbibresource{bibliography.bib}

% Page design for the arXiv version.
\pagestyle{fancy}
\fancyhf{}
\fancyhead[L]{\small A. J. Gomez Saltachin}
\fancyhead[R]{\small Derived categories and Brauer groups}
\fancyfoot[C]{\thepage}
\renewcommand{\headrulewidth}{0.4pt}
\renewcommand{\footrulewidth}{0pt}
\setlength{\parindent}{1.5em}
\setlength{\parskip}{0pt}
\emergencystretch=2em

\begin{document}

\title{Derived categories and Brauer groups of log del Pezzo surfaces with \texorpdfstring{$\frac13(1,1)$}{1/3(1,1)} singularities}
\author{Alex Junior Gomez Saltachin\\[0.5em]
\small Departamento de Matem\'atica, Pontif\'icia Universidade Cat\'olica do Rio de Janeiro\\
\small Rua Marqu\^es de S\~ao Vicente, 225, Rio de Janeiro, RJ 22451-900, Brazil\\
\small \href{mailto:alex.gomez@pucp.edu.pe}{alex.gomez@pucp.edu.pe}}
\date{}

\maketitle

\begin{abstract}
We study the derived categories of log del Pezzo surfaces with only $\tfrac13(1,1)$ singularities. For the Fano index one families classified by Corti and Heuberger, we construct explicit full exceptional collections on their canonical stacks using blow-up presentations of the minimal resolutions and cascades of smooth points. More generally, for any log del Pezzo surface $X$ with $k$ such singularities, we obtain the length formula $12-K_X^2+4k/3$ independent of the Fano index.
We show that the intersection lattice of the exceptional curves determines the Brauer group: it vanishes for $k\leq5$ and is isomorphic to $\mathbb Z/3$ for $k=6$. Applying the Karmazyn--Kuznetsov--Shinder framework, we obtain semiorthogonal decompositions of the bounded derived categories of coherent sheaves on the singular surfaces for the one- and three-point cascades, together with a nontrivial Brauer-twisted decomposition for the six-point cascade.
For quasismooth degree-$10$ hypersurfaces in $\mathbb P(1,2,3,5)$, we construct a thirteen-object full exceptional collection on the canonical stack, compute its Euler matrix, and obtain a semiorthogonal decomposition of the derived category of the singular surface. We also identify a general such hypersurface as the anticanonical model of the blow-up of $\mathbb P(1,1,3)$ at eight general smooth points, with the converse holding for general centers.
Finally, for rational surfaces with arbitrary baskets of $\tfrac1{n_i}(1,1)$ singularities, we give a general block formula for the Euler matrices of full exceptional collections on their canonical stacks.
\end{abstract}


\medskip
\noindent\textbf{2020 Mathematics Subject Classification.}
Primary 14F08; Secondary 14J45, 14A20, 14J17, 14E15, 14J26.

\smallskip
\noindent\textbf{Keywords.}
Log del Pezzo surfaces, canonical stacks, exceptional collections, semiorthogonal decompositions, quotient singularities, Brauer groups, special McKay correspondence.

\section{Introduction}

Del Pezzo surfaces offer a particularly explicit setting in which mirror
symmetry interacts with the birational classification of Fano varieties.
The Fukaya--Seidel formalism for Lefschetz fibrations \cite{Seidel2008}
gives the symplectic framework for homological mirror symmetry with
Landau--Ginzburg mirrors. An early stacky instance of this picture was
given by Auroux, Katzarkov, and Orlov
\cite{AurouxKatzarkovOrlov2008}, who proved homological mirror symmetry
for weighted projective planes, treating their generally singular coarse
spaces through smooth quotient stacks. Ueda \cite{Ueda2006} proved the
correspondence for smooth toric del Pezzo surfaces, and Auroux,
Katzarkov, and Orlov \cite{AurouxKatzarkovOrlov2006} established it for
del Pezzo surfaces obtained by blowing up $\mathbb P^2$. Ueda and Yamazaki
\cite{UedaYamazaki2013} subsequently extended this picture to certain
quotient stacks of toric del Pezzo surfaces.

A complementary mirror-symmetric perspective
relates deformation families, Fano polygons, and periods for the same
surfaces. In the smooth case, the ten deformation families correspond to the ten
mutation-equivalence classes of Fano polygons with empty residual basket,
and the classical periods of the corresponding maximally mutable Laurent
polynomials recover their regularized quantum periods
\cite{AkhtarEtAl2016,KasprzykNillPrince2017}. More generally, Akhtar
et~al.\ \cite{AkhtarEtAl2016} formulated a mirror-symmetric classification
program relating mutation classes of Fano polygons to
$\mathbb Q$-Gorenstein deformation classes of locally
$\mathbb Q$-Gorenstein rigid del Pezzo surfaces with cyclic quotient
singularities that admit $\mathbb Q$-Gorenstein toric degenerations.

The simplest non-smooth residual singularity in this program is
$\tfrac13(1,1)$. This $\mathbb Q$-Gorenstein rigid singularity coincides
with its residue in the sense of singularity content \cite{AkhtarEtAl2016,AkhtarKasprzyk2014}.
Surfaces with only these singularities provide a first testing
ground for extending the smooth theory: they exhibit new orbifold
phenomena while remaining accessible to explicit classification.
Corti and Heuberger classified them into $29$ $\mathbb Q$-Gorenstein
deformation families and constructed explicit models; precisely $26$
of these families admit $\mathbb Q$-Gorenstein degenerations to toric
surfaces \cite{Corti2016}.

For these families, Kasprzyk, Nill, and Prince
\cite{KasprzykNillPrince2017} classified the corresponding Fano polygons
up to mutation. Oneto and Petracci \cite{OnetoPetracci2018} computed many
of the quantum periods and compared them with classical periods of
associated Laurent polynomials, using conjectural extensions of Quantum
Lefschetz and the Abelian/non-Abelian Correspondence to the orbifold
setting.

In keeping with this stacky $B$-model perspective, the orbifold invariants
in this picture are naturally defined on the canonical smooth
Deligne--Mumford stack $\pi\colon\mathcal X\to X$ of a log del Pezzo
surface $X$ \cite{AkhtarEtAl2016,OnetoPetracci2018}.
This stack records the stabilizer data at the quotient singularities,
while its bounded derived category of coherent sheaves gives the
categorical $B$-model. We ask how the contribution of a $\tfrac13(1,1)$
point to this category relates to the exceptional curve on the minimal
resolution, and what categorical decompositions descend to $X$.

To address these questions, let $\nu\colon\widetilde X\to X$ be the
minimal resolution. We compare
\[
D^b(\operatorname{coh}\widetilde X),\qquad
D^b(\operatorname{coh}\mathcal X),\qquad \text{and}\qquad
D^b(\operatorname{coh}X)
\]
by constructing exceptional collections on $\widetilde X$ and
using them to obtain semiorthogonal decompositions on the canonical stack
$\mathcal X$ and, under suitable descent conditions, on the singular
surface $X$ itself. We carry out these
constructions along the Corti--Heuberger cascades.

Ishii--Ueda's special McKay correspondence~\cite{Ishii2015} gives a fully faithful functor $\Phi\colon D^b(\operatorname{coh}\widetilde X)\to D^b(\operatorname{coh}\mathcal X)$. For singularities of type $\tfrac1{n_i}(1,1)$, its semiorthogonal complement is generated by $n_i-2$ exceptional objects supported on the residual gerbe at each singular point. For any rational surface with quotient singularities, Orlov's blow-up formula yields a full exceptional collection on its minimal resolution. Ishii--Ueda's result then induces such a collection on the canonical stack, as observed by Gugiatti--Rota~\cite[Remark~4.8]{GugiattiRota2023} and recalled in \cref{prop:stack-existence}. Canonical stacks of log del Pezzo surfaces admit full exceptional collections, since their resolutions are rational (\cref{prop:log-del-pezzo-rational}).

To pass to the singular surface, Karmazyn--Kuznetsov--Shinder (KKS) introduced a framework to construct semiorthogonal decompositions of $D^b(\operatorname{coh}X)$ from suitable decompositions on a resolution \cite[Theorems~2.12 and~3.16]{KKS2020}. A decomposition on $\widetilde X$ that is compatible with the contraction descends to $X$. Under the stronger adherence condition, which prescribes the form of the components associated with contracted curves, the components corresponding to contracted exceptional curves become equivalent to bounded derived categories of modules over finite-dimensional local algebras.

For surfaces with only $\tfrac13(1,1)$ singularities, general results guarantee full exceptional collections on their canonical stacks and determine their Brauer group orders via the Corti--Heuberger defect \cite[Section~2.1, Proposition~18(b), Lemma~28]{Corti2016}. Blow-up presentations of the minimal resolutions of the six rigid bases, together with Orlov's formula, give explicit exceptional collections. We identify the Brauer groups through saturation of the exceptional lattice, including the alternating class for $Z_6$, and construct untwisted KKS decompositions for the one- and three-point cascades and a Brauer-twisted decomposition for the six-point cascade. For $X_{10}$, we relate general hypersurfaces to eight general points on $\mathbb P(1,1,3)$. The Euler-matrix formula extends to arbitrary baskets of $\tfrac1{n_i}(1,1)$ points.

The constructions below provide explicit $B$-model data---exceptional
collections, Euler pairings, and coarse-space components---for the
$\tfrac13(1,1)$ families appearing in the mirror-symmetric classification
program.

Every Fano index one family $X_{k,d}$, indexed by the number $k$ of its singular points and by the degree $d=K_X^2$, belongs to one of six cascades whose members arise from a rigid base $Z_k$ by blowing up $K_{Z_k}^2-d$ smooth points \cite[Theorems~4 and~6, Corollary~8]{Corti2016}. Here $Z_1=\mathbb P(1,1,3)$, while for $2\le k\le6$, $Z_k$ is the unique surface in its family; Table~\ref{tab:rigid-base-invariants} lists the invariants.

In a marked blow-up presentation, the marking records the ordered centers, their incidence data, and the curves contracted to singular points (\cref{def:marked-blowup-presentation}). The starting smooth surfaces $Y_0$ and collection lengths $N_k$ are listed in Table~\ref{tab:rigid-base-invariants}. Write $\mathcal A_k=\langle A^{(k)}_1,\dots,A^{(k)}_{N_k}\rangle$ for the exceptional collection on $\widetilde Z_k$ (\cref{prop:explicit-rigid-bases}, Appendix~\ref{app:explicit-bases}). The presentations of $\widetilde Z_k$ may involve infinitely near points, whereas the additional blow-up centers on $Z_k$ are distinct by ampleness, as explained after \cref{def:marked-blowup-presentation}.

The following theorem gives full exceptional collections on the canonical stacks of surfaces in the Fano index one families. Here the residual-gerbe objects are the exceptional objects supported at the stacky singular points.

\begin{theorem}[Explicit collections in the Corti--Heuberger cascades]
\label{thm:explicit-cascades-main}
Let $X\in X_{k,d}$ be a member of a Fano index one Corti--Heuberger
family, with singular points $p_1,\dots,p_k$, and let $X\to Z_k$ be a blow-up at
$s=K_{Z_k}^2-d$ distinct smooth points of the rigid base.
Let $\gamma\colon\widetilde X\to\widetilde Z_k$ be the induced map of
minimal resolutions, with exceptional curves $D_1,\dots,D_s$.
For the base collection
$\mathcal A_k=\langle A^{(k)}_1,\dots,A^{(k)}_{N_k}\rangle$,
the ordered collection
\[
\bigl\langle \mathcal O_{D_s}(-1),\dots,\mathcal O_{D_1}(-1),\,
L\gamma^*A^{(k)}_1,\dots,L\gamma^*A^{(k)}_{N_k}\bigr\rangle
\]
is a full exceptional collection in
$D^b(\operatorname{coh}\widetilde X)$.
Its image under the Ishii--Ueda embedding $\Phi$, preceded by the
residual-gerbe objects $\mathcal E_{p_1},\dots,\mathcal E_{p_k}$,
is a full exceptional collection in
$D^b(\operatorname{coh}\mathcal X)$ of length
\[
N_k+s+k=12-d+\frac{4k}{3}.
\]
\end{theorem}

Appendix~\ref{app:explicit-bases} constructs marked blow-up presentations of
the six base resolutions, and \cref{cor:explicit-cascade-collections}
extends their Orlov collections along the Corti--Heuberger cascades.

For these rational surfaces, Bright's exact sequence~\cite{Bright2013} yields
\[
\operatorname{Br}(X)\cong\operatorname{coker}
\bigl(\theta\colon\operatorname{Pic}\widetilde X\to\Lambda^\vee\bigr),
\]
where $\Lambda\subset\operatorname{Pic}\widetilde X$ is the sublattice generated by the curves contracted by $\nu$. The map $\theta$ sends a divisor class to its intersection pairing with these curves (\cref{cor:brauer-cokernel}). The saturation calculation in Appendix~\ref{app:brauer} then identifies this cokernel with the Pontryagin dual of the saturation quotient:
\[
\operatorname{Br}(X)\cong
\operatorname{Hom}\!\left(\Lambda^{\mathrm{sat}}/\Lambda,
\mathbb Q/\mathbb Z\right),
\qquad\text{where}\qquad
\Lambda^{\mathrm{sat}}=(\Lambda\otimes\mathbb Q)
\cap\operatorname{Pic}\widetilde X.
\]
Because $\Lambda^{\mathrm{sat}}/\Lambda\cong\operatorname{Cl}(X)_{\mathrm{tors}}$, the Brauer group is also the dual of the torsion subgroup of the divisor class group (\cref{lem:brauer-glue-duality}). Since an untwisted adherent KKS decomposition forces $\operatorname{Cl}(X)$ to be torsion-free~\cite[Corollary~4.8]{KKS2020}, a nonzero Brauer group immediately obstructs such a decomposition. Invariance of the Brauer group under blow-ups at smooth points (\cref{lem:brauer-cascade-invariance}) reduces its computation to the six rigid bases and gives the following theorem.

\begin{theorem}[Brauer obstruction in the Corti--Heuberger cascades]
\label{thm:brauer-cascades-intro}
Let $X$ belong to a Corti--Heuberger cascade with $k$ singular points, all of
type $\tfrac13(1,1)$. Then
\[
\operatorname{Br}(X)\cong
\begin{cases}
0, & 1\le k\le 5,\\
\mathbb Z/3, & k=6.
\end{cases}
\]
Among the six cascades, only the six-point case has a Brauer obstruction
to an untwisted adherent decomposition in the sense of KKS.
\end{theorem}

Appendix~\ref{app:brauer} computes these intersection maps via toric divisor relations and configurations of curves forming $k$-gons, explicitly exhibiting the alternating generator
\[
\tfrac13(E_1-E_2+E_3-E_4+E_5-E_6) \in \Lambda^{\mathrm{sat}}/\Lambda
\]
for $Z_6$. The bases $Z_k$ are toric precisely when $k\in\{1,3,6\}$~\cite[Tables~1, 2 and~4]{Corti2016}, \cite[Table~1]{OnetoPetracci2018}. On their minimal resolutions, we construct subcategories adherent to the curves contracted to singular points—incorporating a nontrivial Brauer twist in the six-point case. Because compatibility and adherence are preserved under additional blow-ups away from these curves (\cref{lem:kks-blowups}), we obtain semiorthogonal decompositions of $D^b(\operatorname{coh}X)$ for the one- and three-point cascades (Corollaries~\ref{cor:one-point-kks-sod} and~\ref{rem:three-point-toric-descent}), and of $D^b(\operatorname{coh}(X,\alpha))$ for the six-point cascade, with a nonzero class $\alpha\in\operatorname{Br}(X)$ (\cref{cor:six-point-twisted-descent}).

For $k\in\{2,4,5\}$, the Brauer group vanishes. Although necessary, this vanishing does not by itself produce compatible decompositions with simultaneous adherent subcategories. These non-toric bases lack the requisite torus action, and their toric degenerations acquire additional singularities (\cref{rem:toric-bases-versus-degenerations}).

Quasismooth degree-$10$ hypersurfaces $X_{10}\subset\mathbb P(1,2,3,5)$ belong to the family $X_{1,1/3}$. Their exceptional collections and derived category decompositions follow from our cascade results. We single out this family to compare its hypersurface parameter space with that of configurations of eight smooth points on $\mathbb P(1,1,3)$, proving that general hypersurfaces correspond to eight general smooth centers and vice versa (part~(iii) of the theorem below).

\begin{theorem}[The hypersurface \(X_{10}\)]\label{thm:X10-main}
Let $X_{10}\subset\mathbb P(1,2,3,5)$ be a quasismooth hypersurface of
degree $10$, with canonical stack $\mathcal X_{10}$ and minimal resolution
$\nu\colon\widetilde X_{10}\to X_{10}$. Then $X_{10}$ has a unique singular
point $p$, of type $\tfrac13(1,1)$, and $\widetilde X_{10}$ is the blow-up
$\sigma\colon\widetilde X_{10}\to\mathbb F_3$ at eight distinct points off
$C_0$, with pairwise disjoint exceptional curves $D_1,\dots,D_8$.
\begin{enumerate}[label=(\roman*)]
\item $D^b(\operatorname{coh}\mathcal X_{10})$ admits the full exceptional
collection
\[
\begin{aligned}
\bigl\langle &\mathcal E_p,\ \Phi(\mathcal O_{D_1}(-1)),\dots,\Phi(\mathcal O_{D_8}(-1)),\
\Phi(\sigma^*\mathcal O),\ \Phi(\sigma^*\mathcal O(f)),\\
&\Phi(\sigma^*\mathcal O(C_0+3f)),\ \Phi(\sigma^*\mathcal O(C_0+4f))\bigr\rangle
\end{aligned}
\]
of length $13$, whose Euler matrix is computed in \cref{prop:X10-euler-form}.
\item The singular surface admits a semiorthogonal decomposition
\[
\begin{gathered}
D^b(\operatorname{coh}X_{10})=\bigl\langle D^b(K(3,1)\text{-mod}),\,F_1,\dots,F_{10}\bigr\rangle,\\
K(3,1)\cong\mathbb C[z_1,z_2]/(z_1,z_2)^2,
\end{gathered}
\]
with exceptional objects $F_i$ made explicit in \cref{thm:X10-kks-sod}.
\item A general such hypersurface is the anticanonical model of the
blow-up of $\mathbb P(1,1,3)$ at eight general smooth points, and conversely
general centers give a general hypersurface
(\cref{prop:general-X10,prop:X10-blowup-dominance}).
\end{enumerate}
\end{theorem}

\Cref{prop:X10-quasismooth} establishes the unique singularity, while
Corti--Heuberger~\cite[Corollary~8(1)]{Corti2016} gives the specified blow-up
presentation at eight smooth points; their distinctness follows from ampleness,
as explained after \cref{def:marked-blowup-presentation}. Part~(i) specializes
\cref{thm:explicit-cascades-main} to $k=1$ and $s=8$. Part~(ii) follows from the KKS decomposition of $\mathbb P(1,1,3)$~\cite[Example~3.17]{KKS2020}, with compatibility and adherence preserved under the eight blow-ups (\cref{cor:cyclic-one-point-kks-sod}). Part~(iii) combines the Reid--Suzuki construction~\cite{ReidSuzuki2004} with the dominance of the classifying map for the hypersurfaces obtained from eight-point configurations, proved in Appendix~\ref{app:general-X10}.

The $X_{10}$ case fits into the broader investigation of log del Pezzo hypersurfaces and their derived categories. Johnson--Koll\'ar~\cite{JohnsonKollar2001} classified anticanonically embedded log del Pezzo hypersurfaces in weighted projective spaces. On the categorical side, Elagin~\cite{Elagin2007} and Gugiatti--Rota~\cite{GugiattiRota2023} constructed exceptional collections for series of log del Pezzo hypersurfaces different from the degree-ten family considered here.

Finally, we give a uniform block upper-triangular formula for the Euler
matrices of full exceptional collections on canonical stacks of rational
surfaces with arbitrary baskets of $\tfrac1{n_i}(1,1)$ singularities.
The diagonal blocks record the local singularity contributions and the
Euler matrix on the resolution; when the resolution collection consists
of line bundles, the off-diagonal blocks are determined by their
intersections with the exceptional curves (\cref{prop:mixed-basket-euler}).
For $X_{10}$, the matrix is a specialization of this formula. We also
obtain uniform collection-length formulas (\cref{cor:mixed-basket-length}),
with the $\tfrac13(1,1)$ case given in \cref{cor:length}.

\Cref{sec:rationality,sec:canonical-stack,sec:kks} recall exceptional
collections on resolutions, the Ishii--Ueda embedding, and KKS descent.
\Cref{sec:cascade} constructs the collections and decompositions for the
cascades and $X_{10}$, and \cref{sec:further-directions} computes the Euler
matrices. Appendix~\ref{app:general-X10} compares the parameter spaces of
eight-point configurations and degree-$10$ hypersurfaces, proving that
general centers give general hypersurfaces and conversely.
Appendix~\ref{app:brauer} computes the Brauer groups, and
Appendix~\ref{app:explicit-bases} gives explicit blow-up presentations of
the minimal resolutions of the six rigid bases.

\subsection*{Notation and conventions}

All varieties and stacks are defined over \(\mathbb C\).
We write $\mathbb F_n$ for the $n$-th Hirzebruch surface, $C_0$ for its
negative section, and $f$ for its fiber class. The notation $\Pp(E)$
denotes the Grothendieck projectivization parametrizing one-dimensional
quotients of $E$. We write $\nu:\widetilde X\to X$ for the minimal
resolution of the singular surface under consideration; other letters
are reserved for auxiliary blow-ups and contractions.
In the general theory, $r=\#\operatorname{Sing}(X)$ and $p_i$ is of type
$\tfrac1{n_i}(1,1)$; for one point we write $n$. In the Corti--Heuberger notation
$X_{k,d}$, $Z_k$, we have $n_i=3$ and $r=k$. We distinguish resolution
curves $E_i$ over the singular points, blow-up curves $D_i$ over smooth
centers, and local objects $\mathcal E_{p_i,j}$, defined in
\cref{def:local-cyclic-blocks}; for a $\tfrac13(1,1)$-point we write
$\mathcal E_{p_i}=\mathcal E_{p_i,2}$.


% ------------------------------------------------------------
\section{Resolutions and exceptional collections}
\label{sec:rationality}
\label{sec:orlov}
% ------------------------------------------------------------

To construct full exceptional collections on canonical stacks in
\cref{sec:canonical-stack}, we first produce them on the minimal
resolutions of log del Pezzo surfaces. As recalled in
\cref{prop:log-del-pezzo-rational} below, the minimal resolution
\(\widetilde X\) of a log del Pezzo surface \(X\) is rational.
Successively contracting \((-1)\)-curves yields a minimal rational surface,
hence \(\Pp^2\) or \(\F_n\) with \(n\neq1\) \cite[Theorem~V.10]{Beauville1996}.
Reversing these contractions exhibits \(\widetilde X\) as an iterated
blow-up of that model. Orlov's blow-up formula transports the standard
collection on the starting surface to the resolution and proves
\cref{cor:resolution-fec}.

We begin with the surface geometry, recalling the rationality argument
before turning to exceptional collections. The relevant background on
log del Pezzo surfaces can be found in \cite{AlexeevNikulin},
and the rationality statement is also recalled in
\cite[Remark~4.3]{GugiattiRota2023}. Our treatment of exceptional
collections on smooth rational surfaces follows \cite{HillePerling2011}.

\begin{definition}
A \emph{log del Pezzo surface} is a normal projective surface \(X\) with
Kawamata log terminal (klt) singularities such that \(-K_X\) is ample.
\end{definition}

For normal surfaces, klt singularities are precisely quotient singularities \cite[Proposition~4.18]{Kollar1998}. They are also rational \cite[Theorem~5.22]{Kollar1998}. If \(\nu:\widetilde X\to X\) is the minimal resolution, then \(\nu_*\OO_{\widetilde X}=\OO_X\) and \(R^i \nu_*\OO_{\widetilde X}=0\) for all \(i>0\). The Leray spectral sequence gives a canonical isomorphism
\[
H^1(\widetilde X,\OO_{\widetilde X})
\cong
H^1(X,\OO_X).
\]

We recall the rationality argument.

\begin{proposition}\label{prop:log-del-pezzo-rational}
Let \(X\) be a log del Pezzo surface, and let \(\nu:\widetilde X\to X\) be its minimal resolution. The surface \(\widetilde X\) is rational. In particular, \(X\) is rational.
\end{proposition}

\begin{proof}
Since \(X\) is klt and \(-K_X\) is ample, Kawamata--Viehweg vanishing for
klt varieties \cite[Theorem~2.70]{Kollar1998} gives
\(H^1(X,\OO_X)=0\). The rationality of the singularities and the Leray
isomorphism above then give
\(H^1(\widetilde X,\OO_{\widetilde X})=0\).

It remains to show that \(H^0(\widetilde X,\OO_{\widetilde X}(2K_{\widetilde X}))=0\). Suppose, to the contrary, that there exists an effective divisor \(D\sim 2K_{\widetilde X}\). The pullback \(-\nu^*K_X\) is nef because \(-K_X\) is ample, so \((-\nu^*K_X)\cdot D\geq 0\). On the other hand, write the discrepancy formula as
\[
K_{\widetilde X}=\nu^*K_X+\sum_i a_iE_i.
\]
The projection formula gives \(\nu^*K_X\cdot E_i=0\) for every
\(\nu\)-exceptional curve \(E_i\) and \((\nu^*K_X)^2=K_X^2\).
Intersecting the discrepancy formula with \(-\nu^*K_X\) therefore gives
\[
(-\nu^*K_X)\cdot D
=
(-\nu^*K_X)\cdot 2K_{\widetilde X}
=
-2K_X^2<0.
\]
This contradiction implies \(H^0(\widetilde X,\OO_{\widetilde X}(2K_{\widetilde X}))=0\). By Castelnuovo's rationality criterion, \(\widetilde X\) is rational. Since \(X\) is birational to \(\widetilde X\), the surface \(X\) is rational as well.
\end{proof}

\begin{example}
The weighted projective plane \(\Pp(1,1,3)\) has a unique singularity
of type \(\frac13(1,1)\). Its minimal resolution is
\(\F_3=\Pp_{\Pp^1}(\OO\oplus\OO(3))\), whose negative section
\(C_0\), with \(C_0^2=-3\), is the exceptional locus
\cite[Example~3.17]{KKS2020}.
\end{example}

We next recall the exceptional collections and blow-up formula used below.

\begin{definition}
Let \(\mathcal T\) be a triangulated category. An object \(G\in\mathcal T\) is \emph{exceptional} if \(\Hom(G,G)=\C\) and \(\Ext^k(G,G)=0\) for all \(k\neq 0\). An ordered sequence \(\langle G_1,\dots,G_n\rangle\) is an \emph{exceptional collection} if each \(G_i\) is exceptional and \(\Ext^\bullet(G_j,G_i)=0\) for all \(j>i\). The collection is \emph{full} if the smallest full triangulated subcategory of \(\mathcal T\) containing \(G_1,\ldots,G_n\) is \(\mathcal T\).
\end{definition}

The projective line has the standard full exceptional collection
\[
\Db(\coh \Pp^1)=\left\langle \OO,\OO(1)\right\rangle,
\]
and the projective plane has Beilinson's full exceptional collection \cite{Beilinson1978}
\[
\Db(\coh \Pp^2)=\left\langle \OO,\OO(1),\OO(2)\right\rangle.
\]

Orlov's projective bundle formula \cite[Corollary 2.7]{Orlov1993}, applied to
\[
\F_n=\Pp_{\Pp^1}(\OO\oplus\OO(n))
\]
gives a full exceptional collection on \(\F_n\). The intersection numbers are \(C_0^2=-n\), \(C_0\cdot f=1\), and \(f^2=0\), and the relative tautological line bundle is
\[
\OO_{\F_n}(1)\simeq \OO(C_0+nf).
\]
Orlov's projective-bundle formula gives
\[
\Db(\coh\F_n)
=
\left\langle
p^*\Db(\coh\Pp^1),
p^*\Db(\coh\Pp^1)\otimes\OO_{\F_n}(1)
\right\rangle,
\]
where \(p:\F_n\to\Pp^1\) is the ruling. For every \(s\in\Z\),
\(\langle\OO_{\Pp^1}(s-n),\OO_{\Pp^1}(s-n+1)\rangle\) is a full
exceptional collection of \(\Db(\coh\Pp^1)\).
Using the standard pair in the first component and this shifted pair in
the second, together with
\(\OO(C_0+sf)=\OO_{\F_n}(1)\otimes p^*\OO_{\Pp^1}(s-n)\), gives
\begin{equation}\label{eq:hirzebruch-family}
\left\langle\OO,\OO(f),\OO(C_0+sf),\OO(C_0+(s+1)f)\right\rangle,
\qquad s\in\Z,
\end{equation}
a full exceptional collection for every value of \(s\); the case \(s=n\) recovers the standard collection.

\begin{theorem}[Orlov's blow-up formula {\cite[Theorem~4.3 and Corollary~4.4]{Orlov1993}}]\label{thm:orlov-blow-up}
Let \(S\) be a smooth projective surface, and let \(\beta:\Bl_pS\to S\) be the blow-up of \(S\) at a point \(p\), with exceptional curve \(D\). There is a semiorthogonal decomposition
\[
\Db(\coh \Bl_pS)
=
\left\langle
\OO_D(-1),L\beta^*\Db(\coh S)
\right\rangle,
\]
where \(\OO_D(-1)\) is regarded as an object of \(\Db(\coh \Bl_pS)\) via pushforward along the inclusion \(D\hookrightarrow \Bl_pS\). In particular, if \(\Db(\coh S)\) admits a full exceptional collection, the category \(\Db(\coh \Bl_pS)\) also admits one.
\end{theorem}

\begin{corollary}[{\cite[Theorem~5.6]{HillePerling2011}}]\label{cor:rational-surface-fec}
Every smooth projective rational surface admits a full exceptional collection.
\end{corollary}

\begin{proof}
Express the surface as an iterated point blow-up of \(\Pp^2\) or a
Hirzebruch surface \(\F_n\) \cite[Theorem~V.10]{Beauville1996}. By successive applications
of \cref{thm:orlov-blow-up}, the standard full exceptional collection on
the starting surface extends to a full exceptional collection on the
given surface, with one exceptional object added at each step.
\end{proof}

Together, \cref{prop:log-del-pezzo-rational,cor:rational-surface-fec} yield the following consequence.

\begin{corollary}\label{cor:resolution-fec}
Let \(X\) be a log del Pezzo surface, and let \(\nu:\widetilde X\to X\) be its minimal resolution. The category \(\Db(\coh \widetilde X)\) admits a full exceptional collection.
\end{corollary}

% ------------------------------------------------------------
\section{Canonical stacks and exceptional collections}
\label{sec:canonical-stack}
% ------------------------------------------------------------

This section introduces canonical stacks and the local components in the
special McKay correspondence for singularities of type
\(\frac1{n_i}(1,1)\), with \(n_i\geq2\). The global decomposition yields
the existence statement of \cref{prop:stack-existence}, while the
discrepancy and rank computations give \cref{cor:mixed-basket-length}.
We then specialize to \(\frac13(1,1)\)-singularities and compute the local
adjoint needed for the Euler matrices.

\begin{definition}
Let \(X\) be a normal surface with quotient singularities. The \emph{canonical stack} of \(X\) is the smooth Deligne--Mumford stack \(\pi:\XX\to X\) whose coarse moduli space is \(X\), whose local models at quotient singularities are the corresponding quotient stacks, and whose structure morphism is an isomorphism over the smooth locus of \(X\).
Existence for quotient singularities follows from Vistoli's construction
\cite[Proposition~2.8 and its proof]{Vistoli1989}. Uniqueness up to
isomorphism over \(X\) follows from the universal property in
\cite[Theorem~4.6 and Remark~4.9]{FantechiMannNironi2010}.
\end{definition}

At a singular point of type \(\frac1n(1,1)\), with \(n\geq2\), the
local analytic model is \(\A^2/\boldsymbol\mu_n\), and the canonical
stack is locally \([\A^2/\boldsymbol\mu_n]\). We use the ring-action
convention \(\rho_i(\zeta)=\zeta^i\), with the coordinate functions
\(x,y\in\C[x,y]\) both of weight \(\rho_1\).
All character labels below use this ring-action convention.

For the cyclic quotient \(\A^2/\boldsymbol\mu_n\), let \(R=\C[x,y]\)
and let \(\tau:Y\to\Spec R^{\boldsymbol\mu_n}\) be the minimal resolution.
For a character \(\rho\), set
\(M_\rho:=(R\otimes\rho^\vee)^{\boldsymbol\mu_n}\), and let its associated
\emph{full sheaf} on \(Y\) be
\(\mathcal M_\rho:=\tau^*M_\rho/\text{torsion}\).
These constructions associate indecomposable full sheaves on \(Y\) with
characters of \(\boldsymbol\mu_n\) \cite[Section~2]{Ishii2015}.
Wunram calls such a sheaf \emph{special} when
\(H^1(Y,\mathcal M_\rho^\vee)=0\).
Equivalently, \(\rho\) is a \emph{special representation}; otherwise it is
called \emph{non-special}. This terminology follows Wunram
\cite{Wunram1988} and the exposition in \cite[Section~2]{Ishii2015}.

\begin{lemma}\label{lem:non-special-cyclic}
Let \(X\) be a normal surface with singular points \(p_1,\ldots,p_r\), where \(p_i\) has
type \(\frac1{n_i}(1,1)\), with \(n_i\geq2\). At \(p_i\), the
special representations of \(\boldsymbol\mu_{n_i}\) are precisely
\(\rho_0\) and \(\rho_1\). The non-special representations are
\(\rho_2,\ldots,\rho_{n_i-1}\), so there are \(n_i-2\) of them.
For \(n_i=2\) this list is empty; for \(n_i=3\) it consists of
\(\rho_2\) alone.
\end{lemma}

\begin{proof}
At each point, the Hirzebruch--Jung continued fraction
\(n_i/1=[n_i]\) gives the recursion
\(j_0=n_i\), \(j_1=1\), and \(j_2=j_0-n_i j_1=0\).
Wunram's cyclic special McKay correspondence identifies the entries
preceding the terminal zero with the special character labels, read
modulo \(n_i\) \cite[Theorems~2.3 and~3.2]{Ishii2015}.
These give \(\rho_{j_0}=\rho_{n_i}=\rho_0\) and
\(\rho_{j_1}=\rho_1\); all remaining characters are non-special
\cite[Example~2.6]{GugiattiRota2023}.
\end{proof}

\begin{definition}\label{def:local-cyclic-blocks}
Let \(X\) be a normal surface with singular points \(p_1,\ldots,p_r\) of types
\(\frac1{n_i}(1,1)\), with \(n_i\geq2\), and let \(\XX\) be its
canonical stack. For the residual gerbe
\(\iota_{p_i}:B\boldsymbol\mu_{n_i}\hookrightarrow\XX\),
define the \emph{local exceptional objects}
\[
\mathcal E_{p_i,j}:=
\iota_{p_i*}\bigl(\OO_{B\boldsymbol\mu_{n_i}}\otimes\rho_j\bigr),
\qquad 2\leq j\leq n_i-1,
\]
and their ordered exceptional collection
\[
\mathscr B_{p_i}:=
\langle\mathcal E_{p_i,2},\ldots,\mathcal E_{p_i,n_i-1}\rangle.
\]
The collection is empty when \(n_i=2\). Locally,
\(\mathcal E_{p_i,j}\) is represented by the equivariant module
\(\C[x,y]/(x,y)\otimes\rho_j\),
namely \(\OO_0\otimes\rho_j\) in the notation of Ishii--Ueda.
For a point \(p\) of type \(\frac13(1,1)\), we retain the abbreviation
\(\mathcal E_p:=\mathcal E_{p,2}\).
\end{definition}

\begin{theorem}[Ishii--Ueda {\cite[Theorems~1.4 and~3.1]{Ishii2015}}]\label{thm:ishii-ueda-global}
Let \(X\) be a surface with at worst quotient singularities, let \(\XX\) be its canonical stack, and let \(Y\to X\) be the minimal resolution. There is a fully faithful functor
\[
\Phi:\Db(\coh Y)\to \Db(\coh \XX)
\]
and a semiorthogonal decomposition
\[
\Db(\coh \XX)
=
\left\langle
\mathcal E_1,\dots,\mathcal E_q,\Phi(\Db(\coh Y))
\right\rangle,
\]
where \(\mathcal E_1,\dots,\mathcal E_q\) form an exceptional collection.

If the singularities are \(p_1,\ldots,p_r\) of types
\(\frac1{n_i}(1,1)\), with \(n_i\geq2\), the local description
\cite[Theorem~3.1]{Ishii2015} and
\cite[Example~2.6]{GugiattiRota2023} gives
\[
\Db(\coh\XX)=
\left\langle
\mathscr B_{p_1},\ldots,\mathscr B_{p_r},
\Phi\bigl(\Db(\coh Y)\bigr)
\right\rangle,
\]
with the collections of \cref{def:local-cyclic-blocks} and
\(q=\sum_{i=1}^r(n_i-2)\). Collections at distinct points are mutually
orthogonal and can be ordered arbitrarily; within each collection the
character indices increase from \(2\) to \(n_i-1\).
\end{theorem}

We fix the ordering of \cite[Theorem~1.4]{Ishii2015}, with the exceptional
objects preceding the image of \(\Phi\). Thus
\(\RHom(\Phi A,\mathcal E_i)=0\), which gives the lower-left zero blocks
in the Euler matrices.

The following existence statement was observed in
\cite[Remark~4.8]{GugiattiRota2023}.

\begin{proposition}[Existence on rational quotient surfaces]
\label{prop:stack-existence}
Let \(X\) be a normal projective rational surface with quotient
singularities, and let \(\XX\) be its canonical stack. Then \(\Db(\coh\XX)\) admits a full
exceptional collection. In particular, this holds for every log del
Pezzo surface.
\end{proposition}

\begin{proof}
The minimal resolution of \(X\) is smooth and rational, so
\cref{cor:rational-surface-fec} supplies a full exceptional collection.
Insert it into the Ishii--Ueda decomposition of
\cref{thm:ishii-ueda-global}. The log del Pezzo case follows from
\cref{prop:log-del-pezzo-rational}.
\end{proof}

For the one-curve cyclic singularities, the local collections also determine
an ordered collection and its length.

\begin{corollary}[Mixed-basket length and ordered collection]
\label{cor:mixed-basket-length}
Let \(X\) be a normal projective rational surface with singular
points \(p_1,\ldots,p_r\) of types \(\frac1{n_i}(1,1)\), with
\(n_i\geq2\). Let \(\XX\) be its canonical stack and
\(\nu:\widetilde X\to X\) its minimal resolution. With the local
collections \(\mathscr B_{p_i}\) of \cref{def:local-cyclic-blocks},
for any full exceptional collection
\(\langle G_1,\ldots,G_N\rangle\) on \(\widetilde X\), the
Ishii--Ueda embedding gives the ordered
full exceptional collection
\[
\left\langle
\mathscr B_{p_1},\ldots,\mathscr B_{p_r},
\Phi(G_1),\ldots,\Phi(G_N)
\right\rangle
\]
on \(\XX\). The lengths of the resolution and stack collections are,
respectively,
\[
\begin{aligned}
N&=\rk K_0(\widetilde X)
=12-K_X^2+\sum_{i=1}^r\frac{(n_i-2)^2}{n_i},\\
\ell&=N+\sum_{i=1}^r(n_i-2)
=12-K_X^2+\sum_{i=1}^r
\frac{2(n_i-1)(n_i-2)}{n_i}.
\end{aligned}
\]
\end{corollary}

\begin{proof}
The resolution has a full exceptional collection by
\cref{cor:rational-surface-fec}. Write \(E_i\subset\widetilde X\)
for the exceptional curve over \(p_i\). Thus \(E_i^2=-n_i\), and the
discrepancy formula is
\[
K_{\widetilde X}
=
\nu^*K_X-\sum_{i=1}^r\frac{n_i-2}{n_i}E_i.
\]
To verify the coefficient of \(E_i\), write it as \(a_i\) and
intersect with \(E_i\). Adjunction gives
\(K_{\widetilde X}\cdot E_i=-2-E_i^2=n_i-2\), whereas
\(\nu^*K_X\cdot E_i=0\). Hence
\(-n_i a_i=n_i-2\), giving \(a_i=-(n_i-2)/n_i\).
The curves \(E_i\) are disjoint, so
\[
K_{\widetilde X}^2
=K_X^2-\sum_{i=1}^r\frac{(n_i-2)^2}{n_i}.
\]
Noether's formula
\cite[Appendix~A, Example~4.1.2]{Hartshorne1977} gives
\(12\chi(\OO_{\widetilde X})=K_{\widetilde X}^2+c_2(\widetilde X)\).
For the smooth rational surface \(\widetilde X\)
(\cref{prop:log-del-pezzo-rational}), we have
\(\chi(\OO_{\widetilde X})=1\) and
\(c_2(\widetilde X)=2+\rho(\widetilde X)\), the latter being its
topological Euler characteristic, so
\(\rho(\widetilde X)=10-K_{\widetilde X}^2\). Hence
\[
N=\rk K_0(\widetilde X)=2+\rho(\widetilde X)
=12-K_{\widetilde X}^2
=12-K_X^2+\sum_{i=1}^r\frac{(n_i-2)^2}{n_i}.
\]
Inserting the resolution collection into
\cref{thm:ishii-ueda-global} gives the displayed collection, with
\(n_i-2\) additional objects at each point. Its length is
\[
\ell=N+\sum_{i=1}^r(n_i-2)
=12-K_X^2+\sum_{i=1}^r
\frac{2(n_i-1)(n_i-2)}{n_i}.\qedhere
\]
\end{proof}

For the singularities used in the Corti--Heuberger cascades, the local
collections consist of single exceptional objects. Setting \(n_i=3\) in \cref{cor:mixed-basket-length}
gives the following specialization.

\begin{corollary}\label{cor:length}
Let \(X\) be a log del Pezzo surface whose singularities are
\(p_1,\ldots,p_r\), all of type \(\frac13(1,1)\). Let \(\XX\) be
its canonical stack, and let
\(\widetilde X\to X\) be its minimal resolution. Then
\(\Db(\coh\XX)\) admits a full exceptional collection of length
\(\ell=12-K_X^2+\frac{4r}{3}\). The resolution collection has length
\(N=\rk K_0(\widetilde X)=12-K_X^2+\frac r3\), and each local collection in
\cref{cor:mixed-basket-length} is the singleton
\(\mathcal E_{p_i}=\mathcal E_{p_i,2}\).
\end{corollary}

The preceding existence and length results use only the Ishii--Ueda
semiorthogonal decomposition. For the Euler matrices we also need the opposite
pairings \(\RHom(\mathcal E_i,\Phi A)\). In the projective setting,
\(\Phi\) has both adjoints: the semiorthogonal ordering above is controlled
by the right adjoint, which kills the local objects, while the left adjoint
\(\Psi\) computes them through
\(\RHom(\mathcal E_i,\Phi A)\simeq\RHom(\Psi(\mathcal E_i),A)\).
We now fix the local conventions needed to compute this left adjoint.

With the convention \(\rho_i(\zeta)=\zeta^i\) and \(x,y\) of weight
\(\rho_1\), the equivariant Koszul resolution of
\(\OO_0\otimes\rho_i\) has free terms \(\OO\otimes\rho_i\),
\((\OO\otimes\rho_{i+1})^{\oplus2}\), and
\(\OO\otimes\rho_{i+2}\) in degrees \(0,-1,-2\), respectively.

On \(\mathcal U_3=[\A^2/\muthree]\), the form \(dx\wedge dy\) has
character \(\rho_2\), so
\(\omega_{\mathcal U_3}\simeq\OO_{\A^2}\otimes\rho_2\), whereas the
anticanonical bundle has character \(\rho_1\)
\cite[Section~3, equation~(17)]{GugiattiRota2023}. The local character of
\(\OO_{\mathcal X_{10}}(1)\), the anticanonical bundle in our hypersurface
example, is \(\rho_1\), not \(\rho_2\). This distinction fixes the character
convention in the calculation below and in the Euler pairings of Section~6.

\begin{lemma}[Local left adjoint for \(\frac1n(1,1)\)]
\label{lem:cyclic-left-adjoint}
Let \(n\geq3\), let
\(Y=\operatorname{Tot}_{\Pp^1}\OO_{\Pp^1}(-n)\), with projection
\(p:Y\to\Pp^1\) and zero section \(E\), and let \(\psi\) be the left
adjoint of the local Ishii--Ueda embedding \(\phi\). Then
\begin{equation}\label{eq:cyclic-left-adjoint}
\psi(\OO_0\otimes\rho_i)\simeq
\begin{cases}
0,&0\leq i\leq n-3,\\
\OO_E(-n)[1],&i=n-2,\\
\OO_E(-n+1),&i=n-1.
\end{cases}
\end{equation}
Here \(\OO_E(t)\) has degree \(t\) on \(E\simeq\Pp^1\).
\end{lemma}

\begin{proof}
With our character convention, the formula follows from
\cite[Example~3.5, equation~(26)]{GugiattiRota2023} after translating the
indices. We spell out this translation, together with the Koszul calculation,
to make the character residues and cohomological shifts explicit. In the
notation of \cite[Theorem~3.1]{GugiattiRota2023}, one has
\(\rho_{\mathrm{nat}}=\rho_1^{\oplus2}\),
\(\rho_{\det}=\rho_2\), and \(\rho_1\) is the nontrivial special character.
Thus \(\rho_{\det}^{\vee}=\rho_{n-2}\), and the condition
\(\rho_j\otimes\rho_{\det}=\rho_1\) gives \(j=n-1\). The dual appearing in
\(M_\rho=(\C[x,y]\otimes\rho^\vee)^{\mu_n}\) acts on the representation
used to form the module; it does not change the label of the equivariant
skyscraper \(\OO_0\otimes\rho\). Passing to the opposite geometric indexing
convention negates every character index modulo \(n\). Although only the
characters \(2\leq j\leq n-1\) occur in the exceptional complement, we carry
out the computation for all characters to display the convention completely.

Put \(P_a=p^*\OO_{\Pp^1}(a)\) for every \(a\in\Z\), and
\(\mathcal L_a=\OO_{\A^2}\otimes\rho_a\). The \(\mathcal L_a\) are
periodic modulo \(n\), whereas the line bundles \(P_a\) are not. In
particular, \(P_{a+n}\) cannot be identified with \(P_a\), so character
residues and line-bundle degrees must be kept separate. The weight-one
quotient of \(\C x\oplus\C y\) over \(\Pp^1\) is
\(\OO_{\Pp^1}(1)\), and its powers yield the tautological character
summands \(P_a\) for \(0\leq a<n\). The description of the tautological
bundles associated with the universal \(\mu_n\)-cluster over \(Y\) in
\cite[proof of Theorem~3.1, equations~(20)--(22)]{GugiattiRota2023} therefore
gives \(\psi(\mathcal L_a)=P_{-a}\) in exactly this range.

We now apply \(\psi\) to the equivariant Koszul resolution
\[
0\longrightarrow\mathcal L_{j+2}
\xrightarrow{(-y,x)}\mathcal L_{j+1}^{\oplus2}
\xrightarrow{(x,y)}\mathcal L_j
\longrightarrow\OO_0\otimes\rho_j\longrightarrow0.
\]
Let \(s_0,s_1\) be the homogeneous coordinates on \(\Pp^1\),
pulled back to \(Y\), and let
\(t\in H^0(Y,P_{-n})\) be the section cutting out \(E\).
Write \(u=(-s_1,s_0)^{\mathsf T}\) and \(v=(s_0,s_1)\).
The character indices must first be reduced modulo \(n\), and only then may
\(\mathcal L_a\) be replaced by \(P_{-a}\). With this order of operations,
the resulting complexes in degrees \(-2,-1,0\) are
\[
\begin{array}{ll}
0\leq j\leq n-3:
 & [P_{-j-2}\xrightarrow{u}P_{-j-1}^{\oplus2}
       \xrightarrow{v}P_{-j}],\\[3pt]
j=n-2:
 & [P_0\xrightarrow{tu}P_{1-n}^{\oplus2}
       \xrightarrow{v}P_{2-n}],\\[3pt]
j=n-1:
 & [P_{-1}\xrightarrow{u}P_0^{\oplus2}
       \xrightarrow{tv}P_{1-n}].
\end{array}
\]
Here the factor \(t\) records multiplication across the character boundary
from \(n-1\) to \(0\); crucially, it has degree \(-n\), not degree zero.
The first complex is the exact Euler sequence and hence has no cohomology.
For the second complex, \(\ker v=P_{-n}\), and its only cohomology is the
quotient \(P_{-n}/tP_0=\OO_E(-n)\) in degree \(-1\). For the third,
\(\operatorname{coker}u=P_1\), while the induced map
\(P_1\xrightarrow{t}P_{1-n}\) is injective with cokernel
\(\OO_E(1-n)\) in degree zero. These three cases give
\eqref{eq:cyclic-left-adjoint}, including the stated shifts.
\end{proof}

For several singular points the calculation applies componentwise.
The restriction \(n\geq3\) here concerns only the non-special objects in the
semiorthogonal complement of \(\phi(\Db(\coh Y))\);
the KKS descent theorem used in \cref{sec:kks} also allows \(n=2\).

% ------------------------------------------------------------
\section{The singular surface and the KKS approach}
\label{sec:kks}
% ------------------------------------------------------------

We now study \(\Db(\coh X)\) for the singular surface \(X\).
For a single \(\frac1n(1,1)\)-point, a decomposition on \(\F_n\)
compatible with contraction of its negative section descends by the
Karmazyn--Kuznetsov--Shinder (KKS) theorem
\cite[Theorems~2.12 and~3.16]{KKS2020}; adherence identifies the component
corresponding to the contracted curve.
\Cref{lem:kks-blowups} proves that compatibility and adherence are preserved
under successive point blow-ups whose centers avoid the strict transforms
of the contracted curves, including infinitely near centers. The Corti--Heuberger application in
\cref{sec:cascade} uses \(n=3\).

KKS descent requires
a decomposition on the resolution compatible with the contraction.
The theorem assumes rational singularities, a hypothesis satisfied here
by the klt property, as recalled in \cref{sec:rationality}.

We begin with the compatibility condition that permits descent from a
resolution.

\begin{definition}[{KKS compatibility, \cite[Definition~2.7]{KKS2020}}]\label{def:kks-compatibility}
Let \(\tau:\widetilde X\to X\) be a resolution of a normal projective surface
with rational singularities,
viewed as a birational contraction, and let
\[
\Db(\coh\widetilde X)
=
\left\langle
\widetilde{\mathcal A}_1,\dots,\widetilde{\mathcal A}_t
\right\rangle
\]
be a semiorthogonal decomposition. The decomposition is \emph{compatible
with the contraction \(\tau\)} if, for every irreducible curve \(E\)
contracted by \(\tau\), the object \(\OO_E(-1)\) belongs to one of the
components \(\widetilde{\mathcal A}_i\).
\end{definition}

\begin{theorem}[{KKS descent, \cite[Theorem~2.12(i)]{KKS2020}}]
\label{thm:kks-descent}
In the setting of \cref{def:kks-compatibility}, a compatible decomposition
descends uniquely to
\[
\Db(\coh X)=\langle\mathcal A_1,\ldots,\mathcal A_t\rangle,
\qquad
\mathcal A_i=R\tau_*(\widetilde{\mathcal A}_i).
\]
Let \(\mathcal D_i\) be the union of the exceptional curves assigned to
\(\widetilde{\mathcal A}_i\). The restriction of \(R\tau_*\) induces
an equivalence with the Verdier quotient
\[
\mathcal A_i\simeq
\widetilde{\mathcal A}_i/
\langle\OO_E(-1)\mid E\subset\mathcal D_i\rangle.
\]
Here \(E\) runs over the irreducible components of \(\mathcal D_i\).
In particular, pushforward is an equivalence on a component with no assigned curves.
The descended projection functors have finite cohomological amplitude.
Both assertions are contained in KKS Theorem~2.12(i).
\end{theorem}

To identify the descended components, we need the stronger condition of
adherence. In the one-curve case used here, its untwisted form has the
following concrete description.

\begin{definition}[{One-curve case of \cite[Definition~3.6]{KKS2020}}]
Let \(S\) be a smooth projective surface such
that \(H^\bullet(S,\OO_S)\cong\C\), and let \(E\subset S\) be a smooth
rational curve with \(E^2=-n\), where \(n\geq2\). A triangulated subcategory
\(\mathcal A\subset\Db(\coh S)\) is \emph{untwisted adherent} to \(E\) if
there is a line bundle \(L_0\) such that
\[
\mathcal A
=
\left\langle
L_0,\OO_E(-1)
\right\rangle,
\qquad
L_0\cdot E=n-1.
\]
\end{definition}

By \cite[Lemma~3.5, Definition~3.6, and Example~3.7]{KKS2020}, this is
equivalent to a presentation
\[
\mathcal A=\langle L_0,L_1\rangle,
\qquad
L_1=L_0(E),
\qquad
L_0\cdot E=n-1.
\]
More generally, \(b\)-twisted adherence replaces \(\OO_E(-1)\) by
\(\OO_E(-1+b)\) and the intersection condition by
\(L_0\cdot E=n+b-1\). The one-point decompositions below use \(b=0\);
nonzero twists are needed for the Brauer-twisted six-point construction
in Appendix~\ref{app:brauer-six}.

The following lemma proves preservation of compatibility and adherence
under blow-up for all the contracted curves simultaneously.

\begin{lemma}
\label{lem:kks-blowups}
Let \(S\) be a smooth projective surface and let
\(C_1,\ldots,C_s\subset S\) be pairwise disjoint smooth rational
curves with \(C_i^2=-n_i\), where \(n_i\geq2\). Let
\(\beta:S'\to S\) be the blow-up of a point outside
\(\bigcup_i C_i\), and write \(C_i'\) for the strict transforms.
If a semiorthogonal decomposition of \(\Db(\coh S)\) contains
components \(\mathcal A_i\) such that
\[
\OO_{C_i}(-1)\in\mathcal A_i,
\]
then \(\OO_{C_i'}(-1)\in L\beta^*\mathcal A_i\) for every \(i\).
More generally, if \(H^\bullet(S,\OO_S)\cong\C\) and subcategories
\(\mathcal A_i\) have \(b_i\)-twisted adherent presentations
\[
\mathcal A_i=\langle L_i,L_i(C_i)\rangle,
\qquad L_i\cdot C_i=n_i+b_i-1,
\]
then every \(L\beta^*\mathcal A_i\) is \(b_i\)-twisted adherent to
\(C_i'\) and contains \(\OO_{C_i'}(b_i-1)\).
In particular, untwisted adherence is preserved when all \(b_i=0\).
These conclusions hold after any finite sequence of point blow-ups
whose centers avoid the successive strict transforms of all the
\(C_i\), even when a center is a point infinitely near to an earlier
blow-up center.
\end{lemma}

\begin{proof}
The calculation for a single curve can be carried out independently for
each of the pairwise disjoint curves \(C_i\). Since \(\beta\) is an
isomorphism near every \(C_i\), for any integer \(b_i\) we have
\[
L\beta^*\OO_{C_i}(b_i-1)\simeq\OO_{C_i'}(b_i-1),
\qquad
\beta^*C_i=C_i'.
\]
For the line bundles,
\[
\beta^*(L_i(C_i))=(\beta^*L_i)(C_i'),
\qquad
(\beta^*L_i)\cdot C_i'=L_i\cdot C_i=n_i+b_i-1.
\]
A point blow-up preserves the cohomology of the structure sheaf, and
derived pullback is fully faithful by Orlov's formula. Each pullback
component therefore remains adherent, with the same integer \(b_i\).
Induction gives the conclusions for any finite sequence. The same argument
applies to infinitely near centers, since each blow-up remains an
isomorphism near the successive strict transforms of the \(C_i\).
The first pullback identity with \(b_i=0\) proves compatibility
independently of adherence.
\end{proof}

For one cyclic point, the following proposition
states the KKS decomposition of \(\Pp(1,1,n)\) from
\cite[Example~3.17]{KKS2020} in our notation.

\begin{proposition}[KKS weighted-projective-plane decomposition]
\label{prop:kks-weighted-plane}
Let \(n\geq2\), and let \(\tau:\F_n\to\Pp(1,1,n)\) contract the
negative section \(C_0\), with \(C_0^2=-n\). Writing \(f\) for the
fiber class, the full exceptional collection
\[
\langle\OO(-C_0-f),\OO(-f),\OO,\OO(C_0+nf)\rangle
\]
begins with generators of an untwisted adherent subcategory
\(\widetilde{\mathcal A}_0=
\langle\OO(-C_0-f),\OO(-f)\rangle\) over \(C_0\).
It descends to
\[
\Db(\coh\Pp(1,1,n))=
\left\langle\Db(K(n,1)\text{-mod}),\OO,\OO(n)\right\rangle,
\]
where the Kalck--Karmazyn algebra is
\[
K(n,1)\cong\C[z_1,\ldots,z_{n-1}]/(z_1,\ldots,z_{n-1})^2.
\]
\end{proposition}

\begin{proof}
The collection is the projective-bundle collection used in
\cite[Example~3.17]{KKS2020}, with their exceptional section and fiber
class denoted here by \(C_0\) and \(f\). Fullness follows directly from
\eqref{eq:hirzebruch-family} with \(s=n+1\), which gives
\[
\langle\OO,\OO(f),\OO(C_0+(n+1)f),\OO(C_0+(n+2)f)\rangle.
\]
Tensoring by \(\OO(-f)\) yields
\[
\langle\OO(-f),\OO,\OO(C_0+nf),\OO(C_0+(n+1)f)\rangle.
\]
Since
\[
\omega_{\F_n}=\OO(-2C_0-(n+2)f),
\]
the Serre functor satisfies
\[
S\bigl(\OO(C_0+(n+1)f)\bigr)
=\OO(-C_0-f)[2].
\]
The Serre transform of the last object in a full exceptional collection
may be moved cyclically to the front, giving
\[
\langle\OO(-C_0-f)[2],\OO(-f),\OO,\OO(C_0+nf)\rangle.
\]
Shifting an individual exceptional object preserves semiorthogonality and
does not change the triangulated subcategory it generates, so the first
object may be replaced by \(\OO(-C_0-f)\). The divisor sequence
\[
0\longrightarrow\OO(-C_0-f)\longrightarrow\OO(-f)
\longrightarrow\OO_{C_0}(-1)\longrightarrow0
\]
puts \(\OO_{C_0}(-1)\) in \(\widetilde{\mathcal A}_0\), proving
compatibility. The equality
\((-C_0-f)\cdot C_0=n-1\) proves untwisted adherence.
The surface is rational with rational singularities, so KKS
\cite[Theorem~2.12(i)]{KKS2020} descends this compatible decomposition.
By untwisted adherence, \cite[Theorem~3.16]{KKS2020} identifies the first
descended component with \(\Db(K(n,1)\text{-mod})\). The remaining components descend
equivalently to the line bundles \(\OO\) and \(\OO(n)\), since
\(\tau^*\OO(n)=\OO(C_0+nf)\).
\cite[Example~3.14(2)]{KKS2020} gives the displayed presentation of
\(K(n,1)\).
\end{proof}

The KKS decomposition of \(\Pp(1,1,n)\) above comes from an adherent
subcategory on \(\F_n\). Its derived pullback remains adherent under
successive blow-ups away from the strict transforms of \(C_0\), by
\cref{lem:kks-blowups}, giving the following decomposition.

\begin{corollary}[One-point descent under blow-ups]
\label{cor:cyclic-one-point-kks-sod}
Let \(X\) be a normal projective surface with a unique singularity
of type \(\frac1n(1,1)\), where \(n\geq2\). Suppose its minimal
resolution \(\nu:S\to X\) is obtained from \(\F_n\) by \(m\)
point blow-ups whose centers avoid the negative section \(C_0\)
and its successive strict transforms, and that the exceptional locus
of \(\nu\) is the strict transform of \(C_0\).
Then
\[
\Db(\coh X)=\left\langle
\Db(K(n,1)\text{-mod}),F_1,\ldots,F_{m+2}
\right\rangle,
\]
where all \(F_i\) are exceptional. The centers may be infinitely near,
provided they lie over the smooth locus and avoid the contracted curve.
\end{corollary}

\begin{proof}
We pull back the three components of the decomposition on \(\F_n\)
in \cref{prop:kks-weighted-plane} through the successive blow-ups.
At each blow-up \(\beta\) with exceptional curve \(D\),
Orlov's decomposition, written with the exceptional object on the right, is
\[
\Db(\coh S')=\langle L\beta^*\Db(\coh S),\OO_D\rangle.
\]
By \cref{lem:kks-blowups}, the first component remains untwisted adherent
to the strict transform of \(C_0\). The other two initial components
and the \(m\) new components are generated by exceptional objects. The
new exceptional objects are understood via derived pullback from their
creation stages, as in \cref{lem:kks-blowups}; full faithfulness preserves
exceptionality. The decomposition
is compatible with \(\nu\), whose only exceptional curve is assigned
to the first component. By \cref{thm:kks-descent}, all other components descend
equivalently, and \cite[Theorem~3.16]{KKS2020} identifies the first
with \(\Db(K(n,1)\text{-mod})\).
\end{proof}

For the one-point Cavey--Prince cascades from \(\Pp(1,1,n)\)
\cite{CAVEY2020}, the corollary gives decompositions whenever the minimal
resolutions admit blow-up presentations satisfying its hypotheses.
For \(n=3\), the KKS decomposition of \(\Db(\coh\Pp(1,1,3))\) is
\(\langle\Db(K(3,1)\text{-mod}),\OO,\OO(3)\rangle\), as used in
\cref{sec:cascade}.

For several singularities, constructing simultaneous subcategories
adherent to the exceptional curves is a separate problem. KKS gives a
global obstruction in terms of the divisor class group.

\begin{theorem}[{KKS obstruction, \cite[Corollary~4.8]{KKS2020}}]
\label{thm:kks-obstruction}
Let \(X\) be a normal projective rational surface with cyclic quotient
singularities, and let \(\tau:Y\to X\) be its minimal resolution with
exceptional divisor \(\mathcal D\). Suppose there is a semiorthogonal decomposition
\[
\Db(\coh Y)=\langle\widetilde{\mathcal A}_1,\ldots,
\widetilde{\mathcal A}_t\rangle,
\qquad \mathcal D=\bigsqcup_{i=1}^t\mathcal D_i,
\]
such that every nonempty \(\mathcal D_i\) is a connected component of \(\mathcal D\)
and every \(\widetilde{\mathcal A}_i\) is untwisted adherent to
\(\mathcal D_i\) in the sense of \cite[Definition~3.6]{KKS2020}.
For \(\mathcal D_i=\varnothing\), adherence means that
\(\widetilde{\mathcal A}_i\) is generated by one line bundle,
using the convention of \cite[Remark~3.8]{KKS2020}.
Then \(\Cl(X)\) is torsion-free.
\end{theorem}

\begin{proof}
KKS Corollary~3.18 gives a decomposition into derived categories of
local finite-dimensional algebras and explicitly allows empty chains;
for an empty chain the algebra is \(\C\). Additivity makes \(G_0(X)\) free,
and \(G_0(X)_{\mathrm{tors}}\cong\Cl(X)_{\mathrm{tors}}\).
This is the proof of \cite[Corollary~4.8]{KKS2020}, using
\cite[Corollary~3.18, Lemma~4.1, and Proposition~4.4]{KKS2020}.
\end{proof}

Bright's exact sequence translates this obstruction into a Brauer-group
calculation. We use the cohomological Brauer group
\(\Br(X)=H^2_{\mathrm{\acute et}}(X,\mathbb G_m)_{\mathrm{tors}}\). In
the present projective setting, its torsion classes are represented by Azumaya
algebras by Gabber's theorem, as used in
\cite[proof of Theorem~6]{Bright2013}; thus no distinction between the two
Brauer groups arises for the finite groups computed here. This is also the
Brauer group entering the obstruction in \cite[Corollary~4.8]{KKS2020}.

\begin{lemma}[Bright's exact sequence {\cite[Proposition~1 and Section~3]{Bright2013}}]\label{rem:brauer-obstruction}
Let \(X\) be a normal projective surface with rational singularities and
minimal resolution \(\tau:Y\to X\), where \(Y\) is a rational surface. Let
\(\Lambda_\tau\subset \Pic(Y)\) denote the subgroup generated by the exceptional
curves. The intersection pairing induces a homomorphism
\[
\theta:\Pic(Y)\longrightarrow \Lambda_\tau^\vee,
\qquad
D\longmapsto \bigl(E\mapsto D\cdot E\bigr),
\]
and there is an exact sequence
\[
0
\longrightarrow
\Pic(X)
\longrightarrow
\Pic(Y)
\xrightarrow{\theta}
\Lambda_\tau^\vee
\longrightarrow
\Br(X)
\longrightarrow
\Br(Y).
\]
Here \(\Br(Y)=0\).
\end{lemma}

\begin{corollary}\label{cor:brauer-cokernel}
Let \(X\) be a normal projective surface with quotient singularities and
minimal resolution \(\tau:Y\to X\), where \(Y\) is a rational surface. For the intersection map of
\cref{rem:brauer-obstruction}, there is an isomorphism
\[
\Br(X)\cong\operatorname{coker}
\bigl(\theta:\Pic(Y)\longrightarrow\Lambda_\tau^\vee\bigr).
\]
In particular, \(\Br(X)=0\) if and only if there are divisor classes
\(L_i\in\Pic(Y)\) with \(L_i\cdot E_j=\delta_{ij}\) for all exceptional
curves \(E_j\).
\end{corollary}

\begin{proof}
Quotient singularities are rational, so \cref{rem:brauer-obstruction}
applies, and \(\Br(Y)=0\). Hence
\(\Br(X)\cong\operatorname{coker}(\theta)\). Therefore \(\Br(X)=0\)
exactly when \(\theta\) is surjective, equivalently when every dual basis
vector is represented by a divisor class \(L_i\), that is,
\(L_i\cdot E_j=\delta_{ij}\).
\end{proof}

All log del Pezzo surfaces in this paper satisfy the rationality hypothesis
by \cref{prop:log-del-pezzo-rational}. In this setting, the vanishing in
\cref{cor:brauer-cokernel} is equivalent to the KKS \emph{torsion-free}
condition: the Weil divisor class group \(\Cl(X)\) has no torsion
\cite[Definition~4.7]{KKS2020}. It is necessary for an untwisted adherent
decomposition by \cite[Corollary~4.8]{KKS2020}.

The rigid base of the one-point Corti--Heuberger cascade is
\(\Pp(1,1,3)\). The exceptional lattice of
\(\F_3\to\Pp(1,1,3)\) is \(\Z C_0\). The relation
\(f\cdot C_0=1\) makes the intersection map surjective, so
\cref{cor:brauer-cokernel} gives
\[
\Br(\Pp(1,1,3))=0.
\]

The corresponding KKS decomposition of \(\Pp(1,1,3)\) is given by
\cref{prop:kks-weighted-plane} with \(n=3\).
Appendix~\ref{app:brauer} computes the Brauer obstruction throughout
the six Corti--Heuberger cascades.

% ------------------------------------------------------------
\section{Corti--Heuberger cascades and the surface \texorpdfstring{$X_{10}$}{X10}}
\label{sec:cascade}
% ------------------------------------------------------------

We now specialize the preceding constructions to log del Pezzo surfaces
with only \(\frac13(1,1)\)-singularities. The Corti--Heuberger
classification organizes the Fano index one families into six cascades
from rigid bases; after recalling the deformation terminology, we construct
the corresponding exceptional collections and then treat \(X_{10}\)
explicitly.

\begin{definition}[\(\mathbb Q\)-Gorenstein deformations and rigidity]
\label{def:qg-deformations}
Let \(X\) be a normal projective surface with quotient singularities.
A \emph{\(\mathbb Q\)-Gorenstein deformation} of \(X\) is a flat deformation
that locally lifts to equivariant
deformations of the index-one, or canonical, covers of the quotient
singularities, namely the local cyclic covers on which the canonical
divisor becomes Cartier.
For projective surfaces we use proper families, and \(\mathbb Q\)-Gorenstein deformation
equivalence means being connected by a chain of such families over
connected bases \cite[Basic Concepts, Definition~2 and proof of Lemma~6]{AkhtarEtAl2016}.
A surface is \emph{\(\mathbb Q\)-Gorenstein rigid} if every sufficiently small
\(\mathbb Q\)-Gorenstein deformation is trivial.
\end{definition}

Rigidity of the surface is a global condition, distinct from rigidity of
its singularities. The singularity \(\frac13(1,1)\) is
locally \(\mathbb Q\)-Gorenstein rigid, so small \(\mathbb Q\)-Gorenstein
deformations preserve its analytic
type; the degree is also locally constant \cite[Section~1.1.1]{Corti2016}.

Corti--Heuberger classify the non-smooth del Pezzo surfaces with only
\(\frac13(1,1)\)-points into \(29\) families, grouped into six cascades
\cite[Theorems~3--4]{Corti2016}. The operations relating these families
are described by the following terminology.

\begin{definition}[Floating \((-1)\)-curves and cascades]
\label{def:floating-curves-cascades}
Let \(X\) be a log del Pezzo surface with only \(\frac13(1,1)\)-points.
A \emph{floating \((-1)\)-curve} on \(X\) is a smooth rational curve
contained in its smooth locus with self-intersection \(-1\)
\cite[Definition~5]{Corti2016}. A \emph{cascade} connects the deformation
families by blow-ups at smooth points, or equivalently by reversing
contractions of floating \((-1)\)-curves.
\end{definition}

Contracting a floating \((-1)\)-curve is the inverse of a blow-up at a
smooth point and increases \(K^2\) by one. These operations preserve
the singular points and their analytic types.

\begin{definition}[Marked blow-up presentations]
\label{def:marked-blowup-presentation}
A \emph{marked blow-up presentation} of a smooth surface \(Y\) is a
sequence of point blow-ups
\[
Y_0\xleftarrow{\beta_1}Y_1\xleftarrow{\beta_2}\cdots
\xleftarrow{\beta_m}Y_m\simeq Y,
\]
with fixed identifications of the starting and final surfaces and, at
each stage, a specified point \(q_i\in Y_{i-1}\), its incidences with
the strict transforms of previously introduced exceptional curves, and
a label for the new exceptional divisor \(D_i\subset Y_i\).
For a marked blow-up presentation, the centers may be infinitely
near points of \(Y_0\). When \(Y\) is a
resolution of a surface with quotient singularities, the marking also
specifies the curves contracted to those singularities. A \emph{marked
cascade presentation} specifies distinct smooth points on the fixed
surface \(Z_k\), their order, and the identification of their blow-up
with the member \(X\).
\end{definition}

The distinctness in a cascade follows from ampleness. In any sequence
of blow-ups at smooth points ending in a log del Pezzo surface \(X\),
an exceptional curve has anticanonical degree one when created. If
\(h\geq1\) later centers lie on its successive strict transforms,
its final strict transform \(D\), still in the smooth locus, satisfies
\begin{equation}\label{eq:cascade-distinct}
(-K_X)\cdot D=1-h\leq0,
\end{equation}
contrary to ampleness. Thus the cascade centers are distinct on the
base, and their exceptional curves are disjoint. This restriction does
not apply to the blow-ups forming the minimal resolutions in Appendix~\ref{app:explicit-bases}.

We follow \cite[Theorem~4]{Corti2016}: \(X_{k,d}\) denotes the Fano index
one \(\mathbb Q\)-Gorenstein deformation family with \(k\) singular points and \(d=K_X^2\).
By \cite[Theorem~6(1)--(6) and Remark~7]{Corti2016}, the six rigid bases
are unique in their respective base families, apart from the two auxiliary
families discussed below; we denote them by \(Z_k\).
Table~\ref{tab:rigid-base-invariants}
collects their invariants and the numerical data of the marked blow-up
presentations constructed in Appendix~\ref{app:explicit-bases}. Here \(Y_0\) is
the starting smooth surface, \(N_k=\rk K_0(\widetilde Z_k)\), and
\(\ell_k=N_k+k\) is the length of the resulting full exceptional collection on the canonical stack.
We use \(S_{1,25/3}\) for \(\Pp(1,1,3)\), of Fano index five, to
distinguish it from the Fano index one notation \(X_{k,d}\).

\begin{table}[htbp]
\centering
\small
\renewcommand{\arraystretch}{1.2}
\begin{tabular}{@{}c c c c c c c@{}}
\(k\)&\(Z_k\)&\(K_{Z_k}^2\)&\(Y_0\)&\shortstack{Number of\\blow-ups}&\(N_k\)&\(\ell_k\)\\ \hline
1&\(S_{1,25/3}=\Pp(1,1,3)\)&\(25/3\)&\(\F_3\)&0&4&5\\
2&\(X_{2,17/3}\)&\(17/3\)&\(\F_3\)&3&7&9\\
3&\(X_{3,5}\)&5&\(\F_3\)&4&8&11\\
4&\(X_{4,7/3}\)&\(7/3\)&\(\F_2\)&7&11&15\\
5&\(X_{5,5/3}\)&\(5/3\)&\(\Pp^2\)&9&12&17\\
6&\(X_{6,2}\)&2&\(\F_0\)&8&12&18
\end{tabular}
\caption{Rigid cascade bases, marked blow-up presentations, and collection lengths.}
\label{tab:rigid-base-invariants}
\end{table}

\begin{proposition}\label{prop:explicit-rigid-bases}
Each minimal resolution \(\widetilde Z_k\) has a marked
blow-up presentation by successive point blow-ups of a smooth surface \(Y_0\),
with infinitely near centers allowed. The starting surfaces and numbers of
blow-ups are specified in Table~\ref{tab:rigid-base-invariants}.
Appendix~\ref{app:explicit-bases} records the incidence data, lists the
centers, and constructs the resulting Orlov full exceptional collections.
Applying the Ishii--Ueda embedding and adjoining the residual-gerbe objects
gives full exceptional collections on the canonical stacks of the bases,
of lengths \(\ell_k\) as recorded in Table~\ref{tab:rigid-base-invariants}
and given by \cref{cor:length}.
\end{proposition}

By \cite[Corollary~8(1)--(6)]{Corti2016}, every surface in \(X_{k,d}\)
is obtained from \(Z_k\) by \(K_{Z_k}^2-d\) blow-ups at smooth points,
distinct by \eqref{eq:cascade-distinct}. Orlov's formula extends the base
collection through these blow-ups, and the Ishii--Ueda embedding, together
with the residual-gerbe objects, gives the collection on the canonical
stack in the order of \cref{cor:explicit-cascade-collections}. Invariance
of the Brauer group under smooth blow-ups
(\cref{lem:brauer-cascade-invariance}) reduces its computation throughout
each cascade to the six bases treated in Appendix~\ref{app:brauer}.

The two auxiliary Fano-index-two families \(B_{1,16/3}\) and
\(B_{2,8/3}\) are treated separately in
Appendix~\ref{app:brauer-classification}, using the alternative contraction
sequences of \cite[Remark~9]{Corti2016}.

The collection length is given by \cref{cor:length}. We single out \(X_{10}\)
because its hypersurface description can be compared directly with its
presentation as an eight-point blow-up of \(\Pp(1,1,3)\); this comparison
underlies the parameter-space result of Appendix~\ref{app:general-X10} and
makes the exceptional collection and KKS decomposition explicit.
Consider a degree-\(10\) hypersurface
\(X_{10}\subset\Pp(1,2,3,5)\), writing \(x_0,x_1,x_2,x_3\) for
coordinates of weights \(1,2,3,5\). The singularities of a general member
are standard \cite{JohnsonKollar2001,IanoFletcher2000}; we verify them for
an arbitrary quasismooth member.

\begin{lemma}\label{prop:X10-quasismooth}
A general degree-\(10\) hypersurface in \(\Pp(1,2,3,5)\) is
quasismooth. Every quasismooth member avoids the coordinate
points of weights \(2\) and \(5\) and has a unique singular point
\(p=[0:0:1:0]\), of type \(\frac13(1,1)\).
\end{lemma}

\begin{proof}
For any quasismooth member, the coefficients of
\(x_1^5\) and \(x_3^2\) are nonzero: these are the only degree-ten
monomials whose first derivatives can be nonzero on the respective
punctured coordinate axes. Otherwise the corresponding axis would lie in
the cone and be singular. Thus the hypersurface avoids the ambient
singular points of weights \(2\) and \(5\).
At \(p\), the quotient chart has weights \((1,2,5)\equiv(1,2,2)\pmod3\).
Quasismoothness forces a nonzero coefficient of \(x_0x_2^3\), the only
monomial giving a linear term there. Eliminating \(x_0\) leaves tangent
weights \((2,2)\), equivalent to \((1,1)\). These are all the ambient
singular points, so quasismoothness implies smoothness elsewhere.
\end{proof}

For every member satisfying \cref{prop:X10-quasismooth}, weighted
adjunction and the intersection formula
of \cite{Dolgachev1982} give
\[
\omega_{X_{10}}^\vee=\OO_{X_{10}}(1),
\qquad
K_{X_{10}}^2=\frac13.
\]
The anticanonical class is ample, so \(X_{10}\) is a log del Pezzo surface.
The higher Fano index families in \cite[Theorem~3]{Corti2016} have different
numerical invariants, so the Fano index is one;
\cite[Theorem~4 and Table~2]{Corti2016} identifies the family as
\(X_{1,1/3}\).

By \cite[Corollary~8(1)]{Corti2016}, every such \(X_{10}\) is
the blow-up of \(\Pp(1,1,3)\) at eight smooth points, distinct by
\eqref{eq:cascade-distinct}. Its minimal resolution is therefore
the blow-up of \(\F_3\) at eight distinct points off \(C_0\).
The following lemma and proposition relate general eight-point
configurations to general hypersurfaces. The ring presentation is due to
Reid--Suzuki \cite[Section~1 and Table~1.1]{ReidSuzuki2004}.

\begin{lemma}[Reid--Suzuki]
\label{prop:reid-suzuki-X10}
For eight distinct general smooth points
\(P_1,\dots,P_8\in\Pp(1,1,3)\), the blow-up
\(S=\Bl_{P_1,\ldots,P_8}\Pp(1,1,3)\) has ample anticanonical divisor and
\[
R(S,-K_S)\cong
\frac{\C[x,y,z,w]}{(f_{10})},
\qquad
\deg(x,y,z,w)=(1,2,3,5),
\]
where \(f_{10}\) is weighted homogeneous of degree \(10\). Consequently,
\(S\) is isomorphic to a quasismooth hypersurface
\(X_{10}\subset\Pp(1,2,3,5)\).
\end{lemma}

\begin{proof}
By \cite[Corollary~8(1)]{Corti2016}, an ample member is obtained by blowing
up eight distinct smooth points of \(\Pp(1,1,3)\). The parameter space of
ordered distinct smooth eight-tuples is irreducible, and ampleness of
\(-3K\) is open in the corresponding family of blow-ups
\cite[Tag~0D2S]{StacksProject}. Thus \(-K_S\) is ample for general
centers. Reid and Suzuki compute its anticanonical ring in the displayed
form \cite[Section~1 and Table~1.1]{ReidSuzuki2004}; ampleness then gives
\(S\simeq\operatorname{Proj}R(S,-K_S)\).

The resulting hypersurface is quasismooth. Its punctured anticanonical
cone is locally a \(\mathbb G_m\)-bundle over the smooth locus of \(S\).
Over its \(\frac13(1,1)\)-point it is locally the quotient of
\(\A^2\times\mathbb G_m\) by \(\muthree\), whose faithful action on the
second factor is free. Hence the punctured cone is smooth.
\end{proof}

For a general hypersurface, the following proposition strengthens the
eight-point presentation by allowing the centers to be general.

\begin{proposition}[The general hypersurface]
\label{prop:general-X10}
A general quasismooth
\(X_{10}\subset\Pp(1,2,3,5)\) is isomorphic to the anticanonical model of
the blow-up of \(\Pp(1,1,3)\) at eight distinct general smooth points.
Its minimal resolution is therefore the blow-up of \(\F_3\) at eight
distinct points away from the negative section \(C_0\).
\end{proposition}

Appendix~\ref{app:general-X10} compares the parameter spaces of eight-point
configurations and degree-$10$ hypersurfaces. Parameter counts and
finiteness of blow-up presentations show that the resulting hypersurfaces
form a dominant family.

The exceptional collection and its Euler matrix can be computed from
the eight-point blow-up presentation. For any \(X_{10}\) satisfying the
hypotheses of \cref{thm:X10-main}, let
\[
\sigma:\widetilde X_{10}\longrightarrow \F_3
\]
be a marked blow-up presentation at eight distinct points off \(C_0\). Write
\(D_1,\dots,D_8\) for its pairwise disjoint exceptional curves and
\(f\) for the fiber class on \(\F_3\). Let
\[
\nu:\widetilde X_{10}\longrightarrow X_{10}
\]
be the minimal resolution, and let \(E\) be the strict transform of \(C_0\).
The computation at the end of \cref{sec:kks}, together with invariance
under smooth blow-ups (\cref{lem:brauer-cascade-invariance}), gives
\(\Br(X_{10})=0\).

Orlov's formula extends the standard collection on \(\F_3\) through the
eight blow-ups. Applying the Ishii--Ueda embedding and adjoining the
residual-gerbe object gives the collection on the canonical stack, of
length \(13\) by \cref{cor:length}, displayed in
\cref{thm:X10-main}. The curves \(D_i\) are pairwise disjoint, so
their exceptional objects may be listed in any order.

\begin{corollary}\label{cor:X10-stack-fec} 
Let \(\pi:\mathcal X_{10}\to X_{10}\) be the canonical stack. The category \(\Db(\coh\mathcal X_{10})\) has a full exceptional collection of length \(13\).
\end{corollary}

More generally, rename the parameter in the Hille--Perling family
\eqref{eq:hirzebruch-family} as \(t\) and set \(t=s+4\)
\cite[Proposition~5.2]{HillePerling2011}. This gives the following
collections; \(s=-1\) is the standard one.

\begin{proposition}[The Euler matrix for \(\mathcal X_{10}\)]
\label{prop:X10-euler-form}
Let \(s\in\Z\) and let
\[
\begin{aligned}
\mathcal C(s)=
\big\langle
&\mathcal E_p,\,
\Phi(\OO_{D_1}(-1)),\ldots,\Phi(\OO_{D_8}(-1)),\,
\Phi(\sigma^*\OO),\,\Phi(\sigma^*\OO(f)),\\
&\Phi(\sigma^*\OO(C_0+(s+4)f)),\,
\Phi(\sigma^*\OO(C_0+(s+5)f))
\big\rangle
\end{aligned}
\]
be an exceptional collection of \(D^b(\coh \mathcal X_{10})\) associated with the Hille--Perling toric system for \(\F_3\). The Euler matrix \(S_{ij}(s)=\chi(\mathcal C_i(s),\mathcal C_j(s))\)
is
\[
\begingroup
\setlength{\arraycolsep}{6pt}
\renewcommand{\arraystretch}{1.15}
S_{\mathcal X_{10}}(s)
=
\left(
\begin{array}{c|c|c}
1
&
0_{1\times 8}
&
\begin{matrix}
0&-1&-(s+1)&-(s+2)
\end{matrix}
\\ \hline
0_{8\times 1}
&
I_8
&
-\mathbf 1_{8\times 4}
\\ \hline
0_{4\times 1}
&
0_{4\times 8}
&
\begin{matrix}
1&2&2s+7&2s+9\\
0&1&2s+5&2s+7\\
0&0&1&2\\
0&0&0&1
\end{matrix}
\end{array}
\right).
\endgroup
\]

\end{proposition}

\begin{proof}
This is the specialization \(n=3\), \(m=8\) of
\cref{prop:general-one-point-euler}. The curves \(D_i\) are disjoint,
so their objects may be ordered arbitrarily. The four line bundles have degrees
\(0,1,s+1,s+2\) on \(E\), giving the displayed first row.
\end{proof}


For the singular surface \(X_{10}\), we extend the KKS decomposition of
\(\Pp(1,1,3)\) recalled in \cref{sec:kks} through the eight blow-ups;
\cref{lem:kks-blowups} preserves the adherent component.

\begin{theorem}
\label{thm:X10-kks-sod}
Every \(X_{10}\) satisfying the hypotheses of \cref{thm:X10-main}
admits a semiorthogonal decomposition
\[
\Db(\coh X_{10})
=
\left\langle
\Db(K(3,1)\text{-mod}),F_1,\dots,F_{10}
\right\rangle,
\]
where \(\eta:X_{10}\to\Pp(1,1,3)\) is the blow-down of the chosen
presentation, \(\nu\) is the minimal resolution, and
\[
\begin{aligned}
F_1&=\OO_{X_{10}},&F_2&=\eta^*\OO_{\Pp(1,1,3)}(3),\\
F_{i+2}&=\OO_{\nu(D_i)}&& (1\leq i\leq8)
\end{aligned}
\]
are exceptional objects.
\end{theorem}

\begin{proof}
\Cref{cor:cyclic-one-point-kks-sod}, with \(n=3\) and \(m=8\), gives
this decomposition using the eight-point presentation valid for every
quasismooth \(X_{10}\). Since the centers are distinct
and \(\nu\) is an isomorphism near every \(D_i\), the objects
\(\OO_{D_i}\) in Orlov's decomposition descend to \(\OO_{\nu(D_i)}\). If
\(\tau:\F_3\to\Pp(1,1,3)\) contracts \(C_0\), then
\(\tau\sigma=\eta\nu\) and
\(\tau^*\OO_{\Pp(1,1,3)}(3)=\OO_{\F_3}(C_0+3f)\).
The projection formula and \(R\nu_*\OO=\OO_{X_{10}}\) therefore
identify the two descended line bundles with \(F_1,F_2\).
\end{proof}


\begin{corollary}\label{cor:one-point-kks-sod}
Let \(X\) belong to a Fano index one family \(X_{1,d}\). Choose a marked cascade
presentation by blowing up \(m\) distinct smooth points of \(\Pp(1,1,3)\), so that
\(m=25/3-d\). Then
\[
\Db(\coh X)
=
\left\langle
\Db(K(3,1)\text{-mod}),F_1,\dots,F_{m+2}
\right\rangle,
\]
where \(F_1,\dots,F_{m+2}\) are exceptional objects.
\end{corollary}

\begin{proof}
The marked cascade presentation induces the marked blow-up presentation required
by \cref{cor:cyclic-one-point-kks-sod}, with \(n=3\); its centers
are distinct on \(\F_3\) and avoid the negative section by
\eqref{eq:cascade-distinct}. Apply that corollary.
\end{proof}

% ------------------------------------------------------------
\section{Euler matrices for exceptional collections on canonical stacks}
\label{sec:further-directions}
% ------------------------------------------------------------

\Cref{cor:mixed-basket-length} supplies the full exceptional collections
on the canonical stacks considered here. We give a block formula for their
Euler matrices, separating the local contributions from the Euler matrix
of the resolution collection and the pairings between them. For a
line-bundle collection, these pairings reduce to intersection numbers with
the exceptional curves. We then specialize to a single cyclic point
resolved by blow-ups of a Hirzebruch surface and compute every block
explicitly.

\begin{proposition}[Mixed-basket Euler block formula]
\label{prop:mixed-basket-euler}
Let \(X\) be a normal projective rational surface with singular points \(p_1,\ldots,p_r\) of
types \(\frac1{n_i}(1,1)\), where \(n_i\geq2\), and let
\(E_i\subset\widetilde X\) be the exceptional curve over \(p_i\).
Fix a full exceptional collection \(\langle G_1,\ldots,G_N\rangle\)
on the minimal resolution. For the associated full exceptional collection on the canonical stack of
\cref{cor:mixed-basket-length}, the Euler matrix is
\[
S=
\begin{pmatrix}
T_{n_1}&0&\cdots&0&B_1\\
0&T_{n_2}&\cdots&0&B_2\\
\vdots&\vdots&\ddots&\vdots&\vdots\\
0&0&\cdots&T_{n_r}&B_r\\
0&0&\cdots&0&H
\end{pmatrix},
\qquad H_{ab}=\chi(G_a,G_b).
\]
For an arbitrary exceptional collection on the resolution, the off-diagonal
blocks \(B_i\) are defined by \((B_i)_{j,a}=\chi(\mathcal E_{p_i,j},\Phi(G_a))\),
where \(2\leq j\leq n_i-1\). The local block is
\begin{equation}\label{eq:Tn}
T_n=(I_{n-2}-J_{n-2})^2,
\qquad (J_{n-2})_{ab}=\delta_{b,a+1}\quad(1\leq a,b\leq n-2),
\end{equation}
where \(J_{n-2}\) is the nilpotent shift with ones immediately above
the diagonal. For \(n=2\), the block is empty.

When the resolution collection consists of line bundles
\(G_a=\OO(A_a)\), these pairings are determined by intersection numbers
alone. Put
\(d_{ia}=A_a\cdot E_i\). For \(n_i\geq4\),
\[
(B_i)_{j,a}=
\begin{cases}
0,&2\leq j\leq n_i-3,\\
d_{ia}+1,&j=n_i-2,\\
-d_{ia},&j=n_i-1.
\end{cases}
\]
For \(n_i=3\), the sole row is \((B_i)_{2,a}=-d_{ia}\); for
\(n_i=2\), \(B_i\) has no rows. Thus the resolution matrix \(H\)
and the intersection numbers \(A_a\cdot E_i\) determine the entire
Euler matrix for a line-bundle collection on the resolution.
\end{proposition}

\begin{proof}
The equivariant Koszul resolution on
\([\A^2/\boldsymbol\mu_n]\) gives
\[
\chi(\OO_0\otimes\rho_u,\OO_0\otimes\rho_v)
=\delta_{uv}-2\delta_{v,u+1}+\delta_{v,u+2}.
\]
Within \(2\leq u,v\leq n-1\), no cyclic wraparound term occurs,
so this is \(T_n\). Distinct local collections have disjoint support,
which gives the off-diagonal zeros among them. Full faithfulness of
\(\Phi\) gives
\[
\chi(\Phi(G_a),\Phi(G_b))=\chi(G_a,G_b),
\]
so the lower-right block is exactly \(H\); \cref{thm:ishii-ueda-global}
supplies the lower-left zero blocks.

Let \(\Psi\) be the left adjoint of \(\Phi\). Left adjunction gives
\[
\RHom_{\XX}(\mathcal E_{p_i,j},\Phi(G_a))
\simeq
\RHom_{\widetilde X}(\Psi(\mathcal E_{p_i,j}),G_a).
\]
\Cref{lem:cyclic-left-adjoint}, equivalently
\eqref{eq:cyclic-left-adjoint}, computes
\(\Psi(\mathcal E_{p_i,j})\) locally near \(E_i\). For the indices
belonging to the non-special collection, it gives
\[
\Psi(\mathcal E_{p_i,n_i-2})=\OO_{E_i}(-n_i)[1],
\qquad
\Psi(\mathcal E_{p_i,n_i-1})=\OO_{E_i}(-n_i+1),
\]
with all earlier non-special characters sent to zero. The first
displayed source belongs to the local collection only when \(n_i\geq4\),
and the second only when \(n_i\geq3\).
For the divisor inclusion \(E_i\hookrightarrow\widetilde X\),
duality and \(E_i^2=-n_i\) give
\[
\chi_{\widetilde X}(\OO_{E_i}(t),\OO(A_a))
=-\chi_{\Pp^1}\bigl(\OO(d_{ia}-n_i-t)\bigr).
\]
Together with \(\chi(\OO_{\Pp^1}(q))=q+1\), this gives \(d_{ia}+1\) for the
shifted penultimate image and \(-d_{ia}\) for the final image.
\end{proof}

\begin{remark}[Local checks and orthogonality]
\label{rem:local-euler-check}
For a vector bundle \(V\) on \([\A^2/\mu_n]\), gerbe adjunction
and the ring-action convention of \cref{sec:canonical-stack} give
\[
\RHom(\OO_0\otimes\rho_j,V)
\simeq (V_0\otimes\rho_{-j-2})^{\mu_n}[-2].
\]
Its Euler characteristic is the multiplicity of \(\rho_{j+2}\) in
the fiber \(V_0\). In particular,
\(\chi(\mathcal E_j,\mathcal L_a)=\delta_{a,j+2}\), with
indices read modulo \(n\).

We also use the images of the special local generators. For the special
characters \(j=0,1\), the corresponding free sheaves belong to the image of
\(\phi\) \cite[Proposition~1.1]{Ishii2015}, so
\(\phi(P_0)=\mathcal L_0\) and \(\phi(P_{-1})=\mathcal L_1\).
Applying \(\phi\) to the Euler sequence
\(0\to P_{-1}\to P_0^{\oplus2}\to P_1\to0\) yields
\[
\phi(P_1)\simeq
\operatorname{coker}\bigl(\mathcal L_1
\xrightarrow{(-y,x)}\mathcal L_0^{\oplus2}\bigr)
\simeq (x,y)\otimes\rho_{n-1}.
\]
The coordinate maps are those of the universal cluster. Thus, for
\(n\geq3\),
\[
0\longrightarrow\phi(P_1)\longrightarrow\mathcal L_{n-1}
\longrightarrow\mathcal E_{n-1}\longrightarrow0.
\]
Hence \(\chi(\mathcal E_{n-1},\phi(P_1))=-1\).
This verifies the Euler pairing with the image of \(\OO(f)\), including the
\(-1\) in the \(X_{10}\) Euler matrix. The fiber-multiplicity rule cannot
be applied to \(\phi(P_1)\) itself: it is torsion-free but not locally free
at the origin.

For \(n\geq5\), the objects \(\mathcal E_j\) with
\(2\leq j\leq n-3\) are orthogonal to the resolution image on both
sides: one direction is semiorthogonality, and the other follows from
\(\psi(\mathcal E_j)=0\). This does not split off an orthogonal
summand of the whole category, since
\(\chi(\mathcal E_j,\mathcal E_{j+1})=-2\) and
\(\chi(\mathcal E_j,\mathcal E_{j+2})=1\).
\end{remark}

The local block depends only on the singularity type, while the resolution
contribution depends on the chosen collection and, for line bundles, only
on its intersection numbers with the exceptional curves. We now
specialize \cref{prop:mixed-basket-euler} to a single
\(\frac1n(1,1)\)-point whose resolution is obtained from \(\F_n\) by
\(m\) point blow-ups away from the successive strict transforms of
\(C_0\).

To describe the collection and its matrix blocks, let \(n\geq2\) and
\(m\geq0\). Suppose that a normal projective surface \(X\) has
a minimal resolution fitting into a diagram
\[
\widetilde X\xrightarrow{\ \sigma\ }\F_n,
\qquad
\widetilde X\xrightarrow{\ \nu\ }X,
\]
where \(\sigma\) is a sequence of \(m\) point blow-ups, possibly
infinitely near, whose centers avoid the successive strict transforms
of the negative section \(C_0\subset\F_n\), and \(\nu\)
contracts the strict transform \(E\) of \(C_0\) to the unique singular point
of \(X\), of type \(\frac1n(1,1)\). Let \(D_a\) be the exceptional
curve on the surface produced by the \(a\)-th blow-up and \(G_a\) the derived
pullback of \(\OO_{D_a}(-1)\) from that surface to \(\widetilde X\). We use the
intersection conventions
\[
C_0^2=-n,
\qquad
C_0\cdot f=1,
\qquad
f^2=0.
\]
Write \(\mathcal E_i=\mathcal E_{p,i}\) for the objects of
\cref{def:local-cyclic-blocks} at the unique singular point, and
retain the notation \(\XX,\Phi\) of \cref{prop:mixed-basket-euler}.
For \(s\in\Z\), the projective-bundle collection
\eqref{eq:hirzebruch-family}, with its parameter renamed \(t\) and set to
\(t=s+n+1\), together with Orlov's blow-up formula and the Ishii--Ueda
decomposition, gives the ordered full exceptional collection
\[
\begin{aligned}
\mathcal C_{n,m}(s)=\big\langle
&\mathcal E_2,\dots,\mathcal E_{n-1},
\Phi(G_m),\dots,\Phi(G_1),\\
&\Phi(\sigma^*\OO),\Phi(\sigma^*\OO(f)),
\Phi(\sigma^*\OO(C_0+(s+n+1)f)),\\
&\Phi(\sigma^*\OO(C_0+(s+n+2)f))
\big\rangle .
\end{aligned}
\]

The local block is \(T_n\) from \eqref{eq:Tn}.
For \(n\geq4\), the local-to-Hirzebruch block
\(B_n(s)\in\operatorname{Mat}_{(n-2)\times4}(\Z)\) is
\begin{equation}\label{eq:Bn}
B_n(s)
=
\begin{pmatrix}
0_{(n-4)\times4}\\[2pt]
\begin{matrix}
1&2&s+2&s+3\\
0&-1&-(s+1)&-(s+2)
\end{matrix}
\end{pmatrix},
\end{equation}
where the zero block is absent for \(n=4\). The penultimate and final rows
are indexed by \(\mathcal E_{n-2}\) and \(\mathcal E_{n-1}\), respectively.
For \(n=3\), the local collection consists only of \(\mathcal E_2\), so
\(T_3=(1)\). The representation \(\rho_{n-2}=\rho_1\) is special and does
not belong to the Ishii--Ueda complement; consequently, only the final row
of \(B_n(s)\) occurs:
\[
B_3(s)
=
\begin{pmatrix}
0&-1&-(s+1)&-(s+2)
\end{pmatrix}.
\]
For \(n=2\), both \(T_2\) and \(B_2(s)\) are empty, and the
collection consists only of the resolution objects.

The Euler matrix \(H_n(s)\in M_4(\Z)\) of the four line bundles on \(\F_n\) is
\[
H_n(s)
=
\begin{pmatrix}
1&2&2s+n+4&2s+n+6\\
0&1&2s+n+2&2s+n+4\\
0&0&1&2\\
0&0&0&1
\end{pmatrix}.
\]

\begin{proposition}[The one-point specialization for a Hirzebruch surface]
\label{prop:general-one-point-euler}
For the collection \(\mathcal C_{n,m}(s)\) and the blocks defined above,
the \((n+m+2)\times(n+m+2)\)
matrix of the Euler form is
\[
\boxed{
S_{n,m}(s)
=
\begin{pmatrix}
T_n&0&B_n(s)\\
0&I_m&-\mathbf1_{m\times4}\\
0&0&H_n(s)
\end{pmatrix}
}.
\]
Here \(\mathbf1_{m\times4}\) denotes the matrix all of whose entries are
equal to \(1\).
\end{proposition}

\begin{proof}
\Cref{prop:mixed-basket-euler} gives \(T_n\) and the local pairings
with the line bundles. Since \(\sigma\) is an isomorphism near \(E\),
\((\sigma^*M)\cdot E=M\cdot C_0\) for a divisor class \(M\) on
\(\F_n\). For the four line bundles in
\(\mathcal C_{n,m}(s)\), the intersection numbers are
\[
M\cdot C_0=0,\ 1,\ s+1,\ s+2.
\]
These values give \eqref{eq:Bn}, including the stated truncation when \(n=3\).
The adjoint images in \cref{lem:cyclic-left-adjoint} are supported on
\(E\). Since every \(G_a\) is supported away from \(E\), the remaining local pairings
satisfy
\[
\chi\bigl(\mathcal E_i,\Phi(G_a)\bigr)=0
\]
for every \(i\) and \(a\).

For a blow-up \(\beta:S'\to S\) at \(p\), with exceptional
curve \(D\), divisor duality gives, for any perfect complex \(A\),
an isomorphism of complexes of vector spaces and the resulting value of
the Euler form:
\[
\RHom_{S'}(\OO_D(-1),L\beta^*A)
\simeq Li_p^*A[-1],
\qquad
\chi(\OO_D(-1),L\beta^*A)=-\rk(A).
\]
Earlier Orlov objects have rank zero, and line bundles have rank one.
Induction, full faithfulness, and Orlov semiorthogonality
\cite[Theorem~4.3 and Corollary~4.4]{Orlov1993} therefore give
\[
\chi(G_a,G_b)=\delta_{ab},
\qquad
\chi(G_a,\sigma^*L)=-1,
\qquad
\chi(\sigma^*L,G_a)=0
\]
for every line bundle \(L\) on \(\F_n\). These pairings yield \(I_m\),
\(-\mathbf1_{m\times4}\), and the opposite zero block, even for
infinitely near centers. In that case the off-diagonal Ext groups
need not vanish, but their Euler characteristics still do.

The block \(H_n(s)\) follows from Riemann--Roch on \(\F_n\). Using
\[
K_{\F_n}=-2C_0-(n+2)f
\]
together with
\[
\chi(\OO(A),\OO(B))
=
1+\frac12(B-A)\cdot(B-A-K_{\F_n}),
\]
gives \(H_n(s)\) for the four Hille--Perling line bundles and
completes the specialization of \cref{prop:mixed-basket-euler}.
\end{proof}

For \(n=3\) and \(m=8\), one has \(T_3=(1)\) and
\(B_3(s)=(0,-1,-(s+1),-(s+2))\), while \(H_3(s)\) is the
\(4\times4\) block in \cref{prop:X10-euler-form}. Together with \(I_8\) and
\(-\mathbf1_{8\times4}\), the formula for \(S_{n,m}(s)\) recovers
\(S_{\mathcal X_{10}}(s)\).

\Cref{prop:general-one-point-euler} also applies to the one-point
Cavey--Prince cascades starting from \(\Pp(1,1,n)\) whenever their
marked blow-up presentations satisfy the stated hypotheses
\cite{CAVEY2020}, including infinitely near centers. The Orlov objects
are derived pullbacks from the stages at which the corresponding exceptional
divisors are created.

% ------------------------------------------------------------

\section*{Concluding remarks}

Semiorthogonal decompositions of \(\Db(\coh X)\) are obtained for the one-
and three-point cascades by
\cref{cor:one-point-kks-sod,rem:three-point-toric-descent}.
For the six-point cascade, \cref{cor:six-point-twisted-descent}
gives a decomposition of \(\Db(\coh(X,\alpha))\) for a nonzero
\(\alpha\in\Br(X)\); by \cref{thm:kks-obstruction}, no untwisted
adherent decomposition of the KKS type considered here exists.
For \(k=2,4,5\), the bases \(Z_k\) are non-toric
(\cref{rem:toric-bases-versus-degenerations}), so the toric construction
used above is unavailable; despite Brauer vanishing, constructing
simultaneous untwisted adherent subcategories remains open.
Appendix~\ref{app:general-X10} establishes, for \(X_{10}\), the
correspondence between general degree-\(10\) hypersurfaces and general
eight-point configurations on \(\Pp(1,1,3)\); analogous comparisons for
other families require separate classifying maps and fiber analyses
(\cref{rem:presentations-versus-dominance}).

For other mixed baskets, \cref{prop:mixed-basket-euler} reduces the
Euler calculation for an exceptional collection of line bundles on the resolution to its
resolution matrix and intersection data. Two natural testing grounds are
the Cavey--Prince models
\(X_{n_1+n_2}\subset\Pp(1,1,n_1,n_2)\)
\cite[Corollary~7.1]{CAVEY2020} and Miura's classification
\cite{Miura2019} of del Pezzo surfaces with \(\tfrac13(1,1)\)- and
\(\tfrac14(1,1)\)-singularities and no floating \((-1)\)-curves.
A further question, in the spirit of the mutation-equivalence perspective
recalled in the Introduction, is whether exceptional collections obtained
from different marked blow-up presentations of the same surface are
related by mutations. More broadly, extending these categorical
constructions to wider singularity baskets would provide further
\(B\)-model data for the mirror-symmetric classification program beyond
the \(\tfrac13(1,1)\) case.

% ------------------------------------------------------------

\appendix
\section{The general degree-ten hypersurface as a blow-up}
\label{app:general-X10}
% ------------------------------------------------------------

Section~\ref{sec:cascade} proves the blow-up presentation for every individual
hypersurface under consideration. Here we prove \cref{prop:general-X10}
by comparing the parameter space of ordered eight-point configurations on
\(\Pp(1,1,3)\) with the parameter space of quasismooth degree-\(10\)
hypersurfaces in \(\Pp(1,2,3,5)\). The construction of
\cref{prop:reid-suzuki-X10} associates such a hypersurface to eight general
smooth points.

Throughout this appendix, put
\[
\begin{gathered}
U\subset\Pp H^0(\Pp(1,2,3,5),\OO(10)),
\qquad G=\operatorname{Aut}\Pp(1,2,3,5),\\
Q=\Pp(1,1,3),\qquad H=\operatorname{Aut}Q,
\end{gathered}
\]
where \(U\) is the nonempty open set of quasismooth hypersurfaces.
Let \(W\) be a nonempty \(H\)-invariant open subset of the space of
ordered distinct smooth eight-tuples on \(Q\) for which
\cref{prop:reid-suzuki-X10} holds. For \(P=(P_1,\ldots,P_8)\in W\),
write \(S_P=\Bl_{P_1,\ldots,P_8}Q\). Thus \(-K_{S_P}\) is ample,
\(K_{S_P}^2=1/3\), and \(S_P\) is its own anticanonical model.

We first record the elementary dimensions of these parameter spaces. In
coordinates of weights \(1,2,3,5\), a degree-\(10\) monomial
\(x^a y^b z^c w^e\) satisfies \(a+2b+3c+5e=10\). For fixed \(e,c\),
there are \(\lfloor(10-5e-3c)/2\rfloor+1\) choices of \(b\), with \(a\)
then determined. The contributions for \(e=0,1,2\) are \(14,5,1\), so
there are \(20\) monomials and \(U\) is a nonempty open subset of
\(\Pp^{19}\). Thus \(\dim U=19\). The dimensions of the graded pieces in
degrees \(1,2,3,5\) are \(1,2,3,6\), respectively. Since invertibility is
an open condition and graded coordinate changes are taken modulo the
one-dimensional subgroup of weighted scalars
\cite[Theorem~4.2]{Cox1995},
\[
\dim G=1+2+3+6-1=11.
\]
The space of ordered distinct smooth eight-tuples is a nonempty open subset
of \((Q_{\mathrm{sm}})^8\). Since \(Q_{\mathrm{sm}}\) is an irreducible
surface, \(W\) is irreducible of dimension \(16\). Finally, a graded
automorphism of \(Q\) has the form
\[
(u,v)\longmapsto A(u,v),\qquad
z\longmapsto az+q_3(u,v),
\]
where \(A\in\GL_2\), \(a\ne0\), and \(q_3\) is a binary cubic, modulo
weighted scalars \cite[Theorem~4.2]{Cox1995}. Hence
\[
\dim H=4+1+4-1=8.
\]

The strategy is a dimension count. Ordered eight-point configurations have
dimension \(16\), while the space \(U\) of hypersurface equations has
dimension \(19\). Allowing graded coordinate changes contributes
\(\dim G=11\), whereas configurations giving the same surface occur, up to
finite ambiguity, in \(H\)-orbits of dimension \(8\). Thus the expected
image dimension is
\[
16+11-8=19=\dim U.
\]
The lemmas below justify the orbit and finiteness statements needed to turn
this count into a proof.

\begin{lemma}[Lifting graded hypersurface isomorphisms]
\label{lem:graded-ring-lift}
Let \(S=\C[x_0,\ldots,x_n]\) be a polynomial ring with positive grading
\(\deg x_i=a_i>0\). Let \(f,g\in S\) be nonzero homogeneous elements of
the same degree \(d\), and assume that \(d>\max_i a_i\). Then every graded
isomorphism
\[
S/(f)\simeq S/(g)
\]
lifts uniquely to a graded automorphism \(\Phi:S\xrightarrow{\sim}S\).
Moreover, \(\Phi(f)=cg\) for some \(c\in\C^\times\).
\end{lemma}

\begin{proof}
Let \(\phi:S/(f)\to S/(g)\) be a graded isomorphism. Since \(f\) and
\(g\) have degree \(d\), their ideals have no nonzero components in the
generator degrees \(a_i<d\). Thus the image of each \(x_i\) under \(\phi\)
has a unique homogeneous representative in \(S_{a_i}\). Sending \(x_i\) to
this representative defines a graded endomorphism \(\Phi:S\to S\) inducing
\(\phi\).

Let \(S_+\) be the positive-degree ideal. Since \(d>\max_i a_i\), every
degree-\(d\) monomial has at least two positive-degree factors, and hence
\(f,g\in S_+^2\). Therefore the indecomposable space of each quotient
ring, namely its positive-degree ideal modulo its square, identifies with
\(S_+/S_+^2\). Consequently, \(\phi\) induces an automorphism of
\(S_+/S_+^2\). Let \(b_1<\cdots<b_r\) be the distinct weights and let
\(X_j\) denote the vector of variables of weight \(b_j\). If a monomial of
weighted degree \(b_j\) contains a variable of weight \(b_j\), then it is
that variable alone, since all weights are positive. Thus
\[
\Phi(X_j)=A_jX_j+P_j(X_1,\ldots,X_{j-1}),
\]
where \(A_j\) represents the weight-\(b_j\) component of the induced
automorphism of \(S_+/S_+^2\), and is therefore invertible. This is block
triangular in increasing weight: after recovering \(X_1,\ldots,X_{j-1}\),
one recovers \(X_j\) as
\(A_j^{-1}(\Phi(X_j)-P_j(X_1,\ldots,X_{j-1}))\). Hence \(\Phi\) is a
graded automorphism.

Since \(\Phi\) is invertible and induces \(\phi\), it carries \((f)\) onto
\((g)\). Comparing the degree-\(d\) components gives \(\Phi(f)=cg\) for
some \(c\in\C^\times\). Finally, any lift of \(\phi\) must send each
generator to its unique homogeneous representative in degree \(a_i<d\),
so \(\Phi\) is unique.
\end{proof}

The following two lemmas apply this result to weighted hypersurfaces of
Fano index one.

\begin{lemma}[Index-one hypersurface orbits]
\label{lem:X10-hypersurface-parameters}
Let \(\Pp=\Pp(a_0,\ldots,a_n)\), where \(n\geq2\), and let
\(d=\sum_{i=0}^n a_i-1\). Two well-formed quasismooth hypersurfaces of
degree \(d\) in \(\Pp\) are isomorphic if and only if they lie in the
same \(\operatorname{Aut}(\Pp)\)-orbit.
\end{lemma}

\begin{proof}
Write \(S=\C[x_0,\ldots,x_n]\), with \(\deg x_i=a_i\), and let
\(X=V(f)\). Weighted adjunction \cite{Dolgachev1982} gives
\(\OO_X(1)\simeq\OO_X(-K_X)\) as reflexive sheaves. Hence the homogeneous
coordinate ring is the intrinsic anticanonical section ring:
\[
R(X,-K_X):=\bigoplus_{m\geq0}H^0(X,\OO_X(-mK_X))
\simeq S/(f).
\]

Let \(X'=V(g)\) be another such hypersurface. An isomorphism
\(X\simeq X'\) induces an isomorphism of their intrinsic anticanonical
rings, hence a graded isomorphism \(S/(f)\simeq S/(g)\). Since
\(d>\max_i a_i\), \cref{lem:graded-ring-lift} applies and
the isomorphism lifts to a graded automorphism \(\Phi\) of
\(S\) satisfying \(\Phi(f)=cg\) for some \(c\in\C^\times\). The induced
automorphism of \(\Pp\) therefore maps \(X\) onto \(X'\). The converse
follows by restricting an ambient automorphism.
\end{proof}

The same anticanonical-ring description identifies ambient stabilizers
with intrinsic automorphism groups.

\begin{lemma}[Index-one hypersurface stabilizers]
\label{lem:X10-hypersurface-stabilizers}
Let \(\Pp\) and \(d\) be as in \cref{lem:X10-hypersurface-parameters},
and put \(\Gamma=\operatorname{Aut}(\Pp)\). For every well-formed
quasismooth hypersurface \(X\subset\Pp\) of degree \(d\), restriction
induces an isomorphism
\[
\operatorname{Stab}_\Gamma(X)\xrightarrow{\sim}\operatorname{Aut}(X).
\]
\end{lemma}

\begin{proof}
Write \(S=\C[x_0,\ldots,x_n]\), with \(\deg x_i=a_i\), and fix
\(X=V(f)\). Restriction defines
\[
\operatorname{res}\colon
\operatorname{Stab}_\Gamma(X)\longrightarrow\operatorname{Aut}(X).
\]
We construct an inverse \(\iota\). Choose
\(\theta:\OO_X(1)\xrightarrow{\sim}\omega_X^\vee\), which, as in the
proof of \cref{lem:X10-hypersurface-parameters}, induces a graded-ring
identification \(j_\theta:S/(f)\xrightarrow{\sim}R(X,-K_X)\). We will
show that the resulting class in \(\Gamma\) is independent of this choice.

Let \(\alpha\in\operatorname{Aut}(X)\). For each \(m\geq0\), pullback
along \(\alpha^{-1}\), followed by the natural identification of the
pullback of the anticanonical sheaf, gives an automorphism
\[
\alpha_m:H^0(X,\OO_X(-mK_X))\xrightarrow{\sim}
H^0(X,\OO_X(-mK_X)).
\]
Functoriality of pullback gives
\(\alpha_{m+m'}(ss')=\alpha_m(s)\alpha_{m'}(s')\), so these maps define
a graded-ring automorphism
\[
\alpha_R:R(X,-K_X)\xrightarrow{\sim}R(X,-K_X).
\]
By \cref{lem:graded-ring-lift}, the transported automorphism
\(j_\theta^{-1}\alpha_Rj_\theta\) of \(S/(f)\) has a unique lift
\(\Phi_\alpha:S\xrightarrow{\sim}S\) preserving \((f)\). Set
\(\iota(\alpha):=[\Phi_\alpha]\in\operatorname{Stab}_\Gamma(X)\).

To verify independence of \(\theta\), let \(\theta'\) be another choice.
The automorphism \(\theta^{-1}\theta'\) of the rank-one reflexive sheaf
\(\OO_X(1)\) is multiplication by a global unit. By
\cite[Chapter~I, Theorem~3.4]{Hartshorne1977}, every global regular
function on the projective variety \(X\) is constant. Thus
\(\theta'=\lambda\theta\) for some \(\lambda\in\C^\times\), and
\(j_{\theta'}\) differs from \(j_\theta\) in degree \(m\) by
multiplication by \(\lambda^m\). This scaling commutes with \(\alpha_R\),
so
\(j_{\theta'}^{-1}\alpha_Rj_{\theta'}=
j_\theta^{-1}\alpha_Rj_\theta\). Equivalently, the scaling on \(S\) is
\[
\rho_\lambda:(x_0,\ldots,x_n)\longmapsto
(\lambda^{a_0}x_0,\ldots,\lambda^{a_n}x_n).
\]
As \(\rho_\lambda\) belongs to the \(\mathbb G_m\)-action defining
\(\Pp\), it represents the identity in \(\Gamma\). Hence
\([\Phi_\alpha]\) is independent of \(\theta\), so \(\iota\) is well
defined.

Finally, \(\Phi_\alpha\) induces \(\alpha_R\) on the anticanonical ring.
Since \(X=\operatorname{Proj}R(X,-K_X)\), its restriction to \(X\) is
\(\alpha\), and hence
\(\operatorname{res}\circ\iota=\operatorname{id}_{\operatorname{Aut}(X)}\).
Conversely, let \(\gamma\in\operatorname{Stab}_\Gamma(X)\), choose a graded
representative \(\Psi\in\operatorname{Aut}_{\mathrm{gr}}(S)\) preserving
\((f)\), and put \(\alpha=\operatorname{res}(\gamma)\). The quotient
automorphisms induced by \(\Psi\) and \(j_\theta^{-1}\alpha_Rj_\theta\)
define the same automorphism of \(\operatorname{Proj}(S/(f))\). Their
linearizations of \(\OO_X(1)\) may differ by a scalar, so for some
\(\lambda\in\C^\times\) they differ in degree \(m\) by multiplication by
\(\lambda^m\). Replacing \(\Psi\) by
\(\rho_\lambda^{-1}\circ\Psi\) makes the quotient automorphisms equal.
Uniqueness in \cref{lem:graded-ring-lift} then gives
\(\rho_\lambda^{-1}\circ\Psi=\Phi_\alpha\). Since
\([\rho_\lambda]=1\) in \(\Gamma\), we have \([\Psi]=[\Phi_\alpha]\), and
\(\iota\circ\operatorname{res}
=\operatorname{id}_{\operatorname{Stab}_\Gamma(X)}\). Therefore
\(\operatorname{res}\) is the asserted isomorphism.
\end{proof}

\begin{lemma}[Generic configuration stabilizers]
\label{lem:X10-configuration-parameters}
There is a nonempty open subset of \(W\) on which the stabilizer in \(H\)
of every ordered tuple is trivial.
\end{lemma}

\begin{proof}
Let \(\mathcal C\) be the space of ordered distinct smooth eight-tuples on
\(Q\), which is irreducible by the discussion above. On the chart
\(u\ne0\), write \(t=v/u\) and \(\zeta=z/u^3\), and associate to an
ordered tuple \(P=((t_j,\zeta_j))_{j=1}^8\) the matrix
\[
M(P)=
\begin{pmatrix}
\zeta_1&1&t_1&t_1^2&t_1^3\\
\vdots&\vdots&\vdots&\vdots&\vdots\\
\zeta_8&1&t_8&t_8^2&t_8^3
\end{pmatrix}.
\]
Let \(\mathcal C^\circ\subset\mathcal C\) be the locus where all eight
points lie in this chart and \(\operatorname{rank}M(P)=5\). This is open,
since its complement in the chart locus is the common zero locus of the
\(5\times5\) minors. If \(\operatorname{rank}M(P)=5\), then the last four
columns \(1,t,t^2,t^3\) have rank \(4\): otherwise their rank would be at
most \(3\), and adjoining the \(\zeta\)-column could give total rank at
most \(4\). The evaluation matrix
\[
\begin{pmatrix}
1&t_1&t_1^2&t_1^3\\
\vdots&\vdots&\vdots&\vdots\\
1&t_8&t_8^2&t_8^3
\end{pmatrix}
\]
has rank \(4\) only if at least four of the \(t_j\) are distinct.

On \(Q\setminus\{\text{singular point}\}\), the projection
\([u:v:z]\mapsto[u:v]\) to \(\Pp^1\) is compatible with the \(H\)-action:
an automorphism
\[
(u,v)\longmapsto A(u,v),\qquad z\longmapsto az+q_3(u,v)
\]
induces the projective transformation \([A]\in\operatorname{PGL}_2\) on
the base. Therefore an automorphism fixing the ordered tuple fixes all the
corresponding projected points. Since at least four of them are distinct,
the induced M\"obius transformation is the identity. Thus
\(A=\lambda I\) for some \(\lambda\in\C^\times\). Composing with the
weighted scalar
\[
(u,v,z)\longmapsto
(\lambda^{-1}u,\lambda^{-1}v,\lambda^{-3}z)
\]
normalizes the representative so that \(A=I\). After renaming the
coefficient of \(z\) and the binary cubic, the residual action is
\[
(t,\zeta)\longmapsto(t,a\zeta+q(t)),\qquad \deg q\leq3.
\]
Write \(q(t)=c_0+c_1t+c_2t^2+c_3t^3\). The condition that every
\((t_j,\zeta_j)\) be fixed is
\[
(a-1)\zeta_j+c_0+c_1t_j+c_2t_j^2+c_3t_j^3=0,
\qquad 1\leq j\leq 8.
\]
This is the homogeneous linear system with coefficient matrix \(M(P)\).
Since \(\operatorname{rank}M(P)=5\), its kernel is zero, so \(a=1\) and
\(q=0\). Thus the normalized representative is the identity.

It remains to check that \(\mathcal C^\circ\) is nonempty. Choose five
distinct values \(t_j\) and set \(\zeta_j=t_j^4\). The corresponding
minor is, up to sign, the Vandermonde determinant
\[
\det(t_j^{\ell-1})_{1\leq j,\ell\leq5}
=\prod_{1\leq i<j\leq5}(t_j-t_i)\ne0.
\]
These five points can be completed by three further distinct points in the
chart to an ordered smooth eight-tuple. Therefore \(\mathcal C^\circ\) is a
nonempty open subset of the irreducible configuration space \(\mathcal C\).
Since \(W\) is also nonempty open, \(W\cap\mathcal C^\circ\) is nonempty
open, and the stabilizer in \(H\) of every tuple in this intersection is
trivial.
\end{proof}

The next lemma shows that a fixed surface arises from finitely many
\(H\)-orbits of ordered configurations. This finiteness determines the
fiber dimension of the morphism \(F\) defined below and hence the
dimension of its image in \(U\).

\begin{lemma}[Finiteness of blow-up presentations]
\label{lem:X10-finite-presentations}
For \(P\in W\), the surface \(S_P\) has finitely many floating
\((-1)\)-curves. A fixed isomorphism class of surfaces \(S_P\)
arises from only finitely many \(H\)-orbits in \(W\).
If the stabilizer of \(P\) in \(H\) is trivial, then
\(\operatorname{Aut}(S_P)\) is finite.
\end{lemma}

\begin{proof}
Let \(\nu:Y\to S_P\) be the minimal resolution and put
\(A=\nu^*(-K_{S_P})\). Then \(A^2=1/3\). A floating
\((-1)\)-curve is disjoint from the singular locus, so its strict transform
\(C\) on \(Y\) is disjoint from the exceptional locus and is a smooth
rational \((-1)\)-curve. Adjunction gives \(K_Y\cdot C=-1\). The discrepancy
divisor is exceptional and disjoint from \(C\), hence
\[
\nu^*K_{S_P}\cdot C=K_Y\cdot C=-1,
\qquad A\cdot C=1.
\]
The coefficient \(3\) in \(v=C-3A\) is determined by
\[
\frac{A\cdot C}{A^2}=\frac{1}{1/3}=3,
\]
so \(v\) is the component of \(C\) orthogonal to \(A\). Thus
\[
A\cdot v=0,\qquad
v^2=(C-3A)^2=C^2-6A\cdot C+9A^2=-4.
\]
Because \(S_P\) has Cartier index \(3\), the divisor \(-3K_{S_P}\) is
Cartier. Consequently \(3A=\nu^*(-3K_{S_P})\) is an integral divisor class
on \(Y\), and hence
\[
v=C-3A\in\operatorname{Pic}(Y)\cap A^\perp.
\]
The Hodge index theorem makes \(A^\perp\) negative definite
\cite[Chapter~V, Theorem~1.9]{Hartshorne1977}. A negative-definite lattice
has only finitely many integral vectors of a fixed norm, so there are only
finitely many possibilities for \(v\), hence for \(C\). A numerical class
of square \(-1\) contains at most one irreducible curve: two distinct curves
in that class would have intersection \(-1\), whereas distinct irreducible
curves on a smooth surface have nonnegative intersection. This proves
finiteness of the floating curves.

The eight exceptional curves of \(S_P\to Q\) are pairwise disjoint floating
\((-1)\)-curves because the eight centers are smooth points. Thus a
presentation as the blow-up of eight ordered distinct smooth points of
\(Q\) determines an ordered set of eight disjoint floating curves. Since
there are only finitely many floating curves, there are finitely many such
ordered sets. After contracting a fixed ordered set, an identification of
the resulting surface with \(Q\) determines the ordered centers, and two
such identifications differ by \(H=\operatorname{Aut}(Q)\). Hence a fixed
isomorphism class gives only finitely many \(H\)-orbits. This proves the
second assertion.

Finally, \(\operatorname{Aut}(S_P)\) permutes the finite set of floating
curves. A finite-index subgroup fixes the eight exceptional curves and
therefore descends injectively to the stabilizer of the ordered tuple \(P\)
in \(H\). Indeed, the contraction is an isomorphism away from the eight
exceptional curves, so an automorphism inducing the identity after contraction
is the identity on their complementary dense open subset, and hence on
\(S_P\). If the tuple stabilizer is trivial, this finite-index subgroup is
trivial, so \(\operatorname{Aut}(S_P)\) is finite.
\end{proof}

We prove that the morphism \(F\) assigning hypersurface equations to
configurations and graded coordinate changes is dominant by showing that
its nonempty fibers have dimension \(8\). These fibers account for the
action of \(H\), without requiring a global moduli space of surfaces.

\begin{proposition}[Dominance of the blow-up construction]
\label{prop:X10-blowup-dominance}
The locus in \(U\) of hypersurfaces isomorphic to \(S_P\) for some
\(P\in W\) contains a dense open subset of \(U\).
The same conclusion holds after replacing \(W\) by any nonempty open
subset of it.
\end{proposition}

\begin{proof}
Choose a nonempty open subset \(W_0\subset W\) on which the tuple
stabilizers are trivial, as in
\cref{lem:X10-configuration-parameters}. We first explain how to choose
hypersurface equations algebraically on a further nonempty open subset
of \(W_0\). Blowing up the eight tautological sections of
\(Q\times W_0\) gives an algebraic family
\(\pi:\mathscr S\to W_0\) with fiber \(S_P\) over \(P\). For \(m\geq0\),
put
\[
\mathcal L_m:=\OO_{\mathscr S}(-mK_{\mathscr S/W_0})^{\vee\vee}.
\]
The eight centers lie in the smooth locus of \(Q\), so near the fixed
quotient singularity the family is unchanged from \(Q\times W_0\); away
from that locus, the blow-up family is smooth over \(W_0\). Thus the
sheaves \(\mathcal L_m\) are flat over \(W_0\) and restrict to
\(\OO_{S_P}(-mK_{S_P})\) on each fiber.
For every \(m\geq0\), Kawamata--Viehweg vanishing for klt surfaces
\cite[Theorem~2.70]{Kollar1998}, applied to the reflexive sheaf
\(\OO_{S_P}(-mK_{S_P})\), gives
\[
H^1(S_P,\OO_{S_P}(-mK_{S_P}))=0,
\]
since \(-(m+1)K_{S_P}\) is ample. By cohomology and base change
\cite[Chapter~III, Theorem~12.11]{Hartshorne1977}, after shrinking \(W_0\)
the sheaves
\(\mathcal E_m:=\pi_*\mathcal L_m\), \(0\leq m\leq10\), are locally free
and their formation commutes with base change. Multiplication in the relative
anticanonical algebra gives morphisms
\(\mathcal E_a\otimes\mathcal E_b\to\mathcal E_{a+b}\), for
\(a+b\leq10\), compatible with multiplication in every fiber.

By \cref{prop:reid-suzuki-X10}, each fiber has one new anticanonical-ring
generator in each of the degrees \(1,2,3,5\), and no additional generators
in the relevant lower degrees. For \(m\in\{1,2,3,5\}\), let
\(\mathcal D_m\subset\mathcal E_m\) be the image of the multiplication maps
generated by lower-degree sections. The fiber dimension of
\(\mathcal E_m/\mathcal D_m\) is one; after shrinking \(W_0\), these
quotients are line bundles. Locally choosing generators of these four line
bundles and lifts to \(\mathcal E_m\) induces a degree-\(10\) evaluation map
\[
\OO_{W_0}[x,y,z,w]_{10}\longrightarrow\mathcal E_{10}.
\]
On every fiber its kernel has dimension one, since
\cref{prop:reid-suzuki-X10} gives a single relation of weighted degree
\(10\). After shrinking, the constant-rank property makes this kernel a
line bundle. We now replace \(W_0\) once and for all by a sufficiently small
nonempty affine open subset on which the four quotient line bundles and the
degree-\(10\) kernel line bundle are trivial. Shrinking further if necessary,
the quotient maps admit chosen global lifts. Choose global sections
\(x,y,z,w\) representing the four generators, and choose a global generator
of the degree-\(10\) kernel. This gives a relation
\(f_P(x,y,z,w)=0\) whose coefficients are regular functions on \(W_0\).
Its projective class therefore defines an honest morphism
\[
h:W_0\longrightarrow U,\qquad V(h(P))\simeq S_P,
\]
where \(V(h(P))\) denotes the hypersurface defined by the projective
form \(h(P)\). The subsequent saturation by the \(G\)-action makes the
argument independent of the chosen graded coordinates.

Allowing all graded coordinate changes gives the morphism
\[
F:G\times W_0\longrightarrow U,\qquad (g,P)\longmapsto g\cdot h(P).
\]
Its image consists precisely of the equations of hypersurfaces
isomorphic to some \(S_P\), \(P\in W_0\), by
\cref{lem:X10-hypersurface-parameters}.

Two pairs \((g,P)\) and \((g',P')\) can map to the same equation only when
\(S_P\simeq S_{P'}\), and \cref{lem:X10-finite-presentations} shows that
configurations with a fixed isomorphism class occur in finitely many
\(H\)-orbits. We now compute the dimension of each nonempty fiber of \(F\).

Fix \(u\in F(G\times W_0)\). Projection of \(F^{-1}(u)\) to
\(W_0\) has image
\[
\{P\in W_0:S_P\simeq V(u)\}.
\]
By \cref{lem:X10-finite-presentations}, this is a finite union of
intersections of \(H\)-orbits with \(W_0\). Since the tuple stabilizer is
trivial and \(\dim H=8\), the orbit-dimension formula
\cite[Proposition~I.1.8]{Borel1991} gives dimension \(8\) for each orbit.
Every nonempty intersection with the open subset \(W_0\) therefore also has
dimension \(8\).
For each tuple in this image, the choices of \(g\) satisfying
\(g\cdot h(P)=u\) form a coset of the stabilizer of \(h(P)\).
This stabilizer is finite by
\cref{lem:X10-hypersurface-stabilizers,lem:X10-finite-presentations}.
The projection therefore has finite fibers, and
\(\dim F^{-1}(u)=8\).

The fiber-dimension theorem now gives
\[
\dim F(G\times W_0)
=\dim G+\dim W_0-8
=\dim U.
\]
Thus \(F\) is dominant. By Chevalley's theorem
\cite[Tag~054K]{StacksProject}, its image is constructible; since it is
dense in the irreducible variety \(U\), it contains a dense open subset
\cite[Tag~005K]{StacksProject}. For any nonempty open subset
\(W'\subset W\), replace \(W_0\) by the nonempty open intersection
\(W_0\cap W'\), which is nonempty because \(W\) is irreducible; the same
dimension and fiber argument proves the final assertion.
\end{proof}

\begin{proof}[Proof of \cref{prop:general-X10}]
By \cref{prop:X10-blowup-dominance}, a general hypersurface in \(U\)
is isomorphic to \(S_P\) with eight distinct general smooth centers.
Since \(-K_{S_P}\) is ample, \(S_P\) is its anticanonical model.
The minimal resolution \(\tau:\F_3\to Q\) is an isomorphism over
the smooth locus of \(Q\). The eight centers therefore lift uniquely
to eight distinct points on \(\F_3\) away from \(C_0\). Blowing up
these points resolves \(S_P\); its exceptional locus over the
singular point is still the single \((-3)\)-curve \(C_0\), so this
resolution is minimal.
\end{proof}

The final assertion of \cref{prop:X10-blowup-dominance} gives the converse
used in \cref{thm:X10-main}(iii): restricting to any nonempty open locus of
centers still produces a dense open subset of \(U\), so general centers give
a general hypersurface.

% ------------------------------------------------------------

\section{Brauer groups and obstructions in the Corti--Heuberger cascades}
\label{app:brauer}
% ------------------------------------------------------------

The proof of \cref{thm:brauer-cascades-intro} combines the cascade description of
\cref{sec:cascade}, the cokernel criterion of \cref{cor:brauer-cokernel},
and invariance under blow-ups at smooth points (\cref{lem:brauer-cascade-invariance}).

We first describe the lattice mechanism for a normal projective rational surface
\(X\) with a finite set of singularities \(p_i\) of types
\(\frac1{n_i}(1,1)\), where \(n_i\geq2\). Its minimal resolution
has pairwise disjoint exceptional curves \(E_1,\ldots,E_r\) with
\(E_i^2=-n_i\). Put \(\Lambda=\bigoplus_i\Z E_i\). In the notation of
\cref{rem:brauer-obstruction}, the intersection map is
\[
\theta:\Pic(\widetilde X)\longrightarrow\Lambda^\vee,
\qquad
D\longmapsto(D\cdot E_1,\dots,D\cdot E_r).
\]
Here \(\Lambda^\vee=\Hom(\Lambda,\Z)\) has basis \(e_1,\dots,e_r\)
dual to the \(E_i\); in particular, \(\theta(E_i)=-n_i e_i\).
By \cref{cor:brauer-cokernel}, \(\Br(X)\cong\operatorname{coker}(\theta)\).

The four-point calculation below instead uses a surface in \(X_{4,4/3}\),
obtained from \(Z_4=X_{4,7/3}\) by one blow-up at a smooth point.

Bright's exact sequence applies to arbitrary rational surface singularities
\cite[Proposition~1]{Bright2013}, whereas the classification carried out in
\cite[Section~3]{Bright2013} concerns the classical singular del Pezzo case
obtained by contracting \((-2)\)-curves. In the Corti--Heuberger cascades the
exceptional curves over the singular points have self-intersection \(-3\).

\subsection{Saturation of the exceptional lattice}
\label{app:brauer-glue}

The intersection lattice \(P=\Pic(\widetilde X)\) is free and unimodular.
Indeed, the rational-surface description used in
\cref{cor:rational-surface-fec} reduces this to \(\Pp^2\) and the
Hirzebruch surfaces, whose intersection matrices have determinant \(\pm1\).
Each point blow-up adds an orthogonal rank-one summand of square \(-1\).
In particular, intersection identifies \(P\) with \(P^\vee\).

\begin{definition}\label{def:brauer-glue}
The \emph{saturation} of the exceptional lattice in \(P\) is
\[
\Lambda^{\mathrm{sat}}
=
\bigl(\Lambda\otimes\mathbb Q\bigr)\cap\Pic(\widetilde X).
\]
The lattice is \emph{primitive} if \(P/\Lambda\) is torsion-free,
equivalently if \(\Lambda^{\mathrm{sat}}=\Lambda\). The finite quotient
\(\Lambda^{\mathrm{sat}}/\Lambda\) is its \emph{saturation quotient}.
Using the intersection form to identify \(\Lambda^\vee\) with a lattice
in \(\Lambda\otimes\mathbb Q\), the \emph{discriminant group}
\[
A_\Lambda:=\Lambda^\vee/\Lambda\cong\bigoplus_{i=1}^r\Z/n_i
\]
carries the induced discriminant bilinear form
\[
b_\Lambda:A_\Lambda\times A_\Lambda\longrightarrow\mathbb Q/\Z,
\qquad
b_\Lambda(\bar x,\bar y)=x\cdot y\pmod\Z.
\]
\end{definition}

An \emph{integral overlattice} of \(\Lambda\) is a finite-index extension
inside \(\Lambda\otimes\mathbb Q\) whose pairings are integral. It lies
in \(\Lambda^\vee\), and its quotient by \(\Lambda\) is totally
isotropic for \(b_\Lambda\), meaning that every pairing between elements
of the quotient is zero. Conversely, the inverse image of such a subgroup
is an integral overlattice. This is the bilinear version of the
overlattice correspondence discussed in \cite[Section~1.4]{Nikulin1980}.
In particular, \(H_\Lambda:=\Lambda^{\mathrm{sat}}/\Lambda\subset A_\Lambda\) is totally isotropic,
since its representatives belong to the integral lattice \(P\).

The next lemma identifies the Brauer group canonically with the dual of
the saturation quotient. Nontrivial \(\Br(X)\), equivalently torsion in
\(\Cl(X)\), is the Brauer obstruction to an untwisted adherent KKS
decomposition: by \cref{thm:kks-obstruction}, such a decomposition requires
\(\Cl(X)\) to be torsion-free.

\begin{lemma}[The Brauer group and the saturation quotient]\label{lem:brauer-glue-duality}
There are natural isomorphisms
\[
\begin{aligned}
\Br(X)&\cong\operatorname{coker}(\theta)
\cong\Hom\bigl(\Lambda^{\mathrm{sat}}/\Lambda,\mathbb Q/\Z\bigr)\\
&\cong\Hom\bigl(\Cl(X)_{\mathrm{tors}},\mathbb Q/\Z\bigr).
\end{aligned}
\]
Thus \(\Br(X)=0\) if and only if \(\Lambda\) is primitive in \(P\).
As abstract finite abelian groups, \(\Br(X)\) and
\(\Lambda^{\mathrm{sat}}/\Lambda\) are isomorphic, but this last
identification is not canonical.
\end{lemma}

\begin{proof}
Applying \(\Hom(-,\Z)\) to \(0\to\Lambda\to P\to P/\Lambda\to0\)
gives an exact sequence
\[
P^\vee\longrightarrow\Lambda^\vee
\longrightarrow\Ext^1(P/\Lambda,\Z)\longrightarrow0,
\]
because \(P\) is free. Under the unimodular identification
\(P\cong P^\vee\), the first map is \(\theta\). For any finitely
generated abelian group \(G\), the sequence
\(0\to\Z\to\mathbb Q\to\mathbb Q/\Z\to0\) gives
\(\Ext^1(G,\Z)\cong\Hom(G_{\mathrm{tors}},\mathbb Q/\Z)\).
Pushforward of divisors identifies \(P/\Lambda\) with \(\Cl(X)\):
every prime divisor on \(X\) has a strict transform on the smooth
surface \(\widetilde X\), and the kernel is generated by the contracted
curves. Hence
\((P/\Lambda)_{\mathrm{tors}}=\Lambda^{\mathrm{sat}}/\Lambda
\cong\Cl(X)_{\mathrm{tors}}\).
The result follows from \cref{cor:brauer-cokernel}; a finite abelian
group is noncanonically isomorphic to its character group.
\end{proof}

The Gram matrix \(\operatorname{diag}(-n_1,\ldots,-n_r)\) gives
\(\Lambda^\vee=\bigoplus_i\Z(E_i/n_i)\). Every class in the saturation quotient thus
has a representative
\[
L=\sum_{i=1}^r\frac{a_i}{n_i}E_i\in\Pic(\widetilde X),
\qquad \bar a_i\in\Z/n_i,
\]
where the \(a_i\) are integer lifts of the residue classes. Since
the intersection form is integral and adjunction gives
\(K_{\widetilde X}\cdot E_i=n_i-2\), every such representative
satisfies
\begin{equation}\label{eq:mixed-glue-integrality}
L^2=-\sum_i\frac{a_i^2}{n_i}\in\Z,
\qquad
K_{\widetilde X}\cdot L=
\sum_i\frac{a_i(n_i-2)}{n_i}\in\Z.
\end{equation}
These are necessary numerical constraints on classes in the saturation quotient, not a
sufficiency criterion for their realization in \(\Pic(\widetilde X)\).

For the remainder of this appendix, specialize to \(k\) points all
of type \(\frac13(1,1)\), so \(r=k\), all \(n_i=3\), and
\(A_\Lambda\cong(\Z/3)^k\). Corti--Heuberger's defect is
\(\delta_{\mathrm{CH}}=\dim_{\mathbb F_3}(\Lambda^{\mathrm{sat}}/\Lambda)\).
They determine it in \cite[Section~2.1, Proposition~18(b), and
Lemma~28]{Corti2016}: \(\delta_{\mathrm{CH}}=1\) for \(X_{6,2},X_{6,1}\)
and \(\delta_{\mathrm{CH}}=0\) otherwise. Together with
\cref{lem:brauer-glue-duality}, they already give the group orders
in \cref{thm:brauer-cascades-intro}. The calculations below exhibit
the intersection maps and classes in the saturation quotient realizing these values.
The next lemma specializes \eqref{eq:mixed-glue-integrality} and
proves primitivity for one or two exceptional curves without choosing a blow-up presentation.

\begin{lemma}\label{lem:brauer-glue-vectors}
Every class in
\(\Lambda^{\mathrm{sat}}/\Lambda\) can be represented in the form
\[
L=\frac13\sum_{i=1}^k a_iE_i,
\qquad
a_i\in\{0,\pm1\}.
\]
Such a representative satisfies the congruences
\[
\sum_{i=1}^k a_i^2\equiv0\pmod3,
\qquad
\sum_{i=1}^k a_i\equiv0\pmod3.
\]
Consequently:
\begin{enumerate}[label=\textup{(\roman*)}]
\item if \(k=1\) or \(k=2\), the lattice \(\Lambda\) is primitive;
\item if \(3\leq k\leq5\), every nonzero class in the saturation quotient is represented by
\[
\pm\frac{E_a+E_b+E_c}{3}
\]
for a triple \(\{a,b,c\}\), and the isotropic subgroup
\(\Lambda^{\mathrm{sat}}/\Lambda\) has order at most \(3\);
\item if \(k=6\), classes supported on three or on six exceptional curves are
numerically possible. On \(Z_6\), the saturation quotient is generated by the
alternating support-six class, so no support-three class is realized; see
\cref{app:brauer-six}.
\end{enumerate}
\end{lemma}

\begin{proof}
For \(n_i=3\), choose the residue representatives
\(a_i\in\{0,\pm1\}\). Formula \eqref{eq:mixed-glue-integrality}
becomes
\[
-\frac13\sum_i a_i^2\in\Z,
\qquad \frac13\sum_i a_i\in\Z,
\]
which gives the two stated congruences.

The first congruence says that the number of nonzero coefficients is divisible
by \(3\). This proves (i). If \(3\leq k\leq5\), a nonzero class has support
of size \(3\), and the second congruence forces its three coefficients to
have the same sign. By the cited overlattice correspondence, the isotropic subgroup
is totally isotropic for \(b_\Lambda\). After changing their signs,
two support-three classes with distinct supports \(A\) and \(B\)
have pairing
\[
-\frac{\#(A\cap B)}{3}\pmod{\Z}.
\]
For two distinct triples in a set of at most five elements, the intersection
has size \(1\) or \(2\), so the two classes cannot be orthogonal. Hence all
nonzero classes in the saturation quotient generate the same cyclic subgroup, proving (ii). For
\(k=6\), supports of size \(3\) and \(6\) satisfy the first numerical
condition, subject also to the stated congruence on the signs. In particular,
the alternating coefficient vector satisfies both congruences.
Thus the numerical conditions permit the coefficient pattern
\(\frac13(E_1-E_2+E_3-E_4+E_5-E_6)\), as asserted in (iii).
\end{proof}

\begin{corollary}\label{cor:brauer-one-two}
Let \(X\) be a log del Pezzo surface with at most two singular points,
all of type \(\frac13(1,1)\). Its Brauer group is trivial. In particular,
\(\Br(Z_1)=\Br(\Pp(1,1,3))=0\) and \(\Br(Z_2)=0\).
\end{corollary}

\begin{proof}
Combine \cref{lem:brauer-glue-duality,lem:brauer-glue-vectors}
(for \(k=0\), this is \(\Br(Y)=0\), as in
\cref{rem:brauer-obstruction}).
\end{proof}

The bound of two singularities in \cref{cor:brauer-one-two} uses the
\((-3)\)-lattice congruences. For \(k=1\), this recovers the direct
fiber-class computation following \cref{cor:brauer-cokernel}. For
\(3\leq k\leq5\), the same lemmas
leave only the possibilities \(0\) and \(\Z/3\). The intersection maps
computed below decide between them.

\subsection{Behavior under cascade blow-ups}
\label{app:brauer-invariance}

The morphisms in a Corti--Heuberger cascade are blow-ups at smooth points by
\cite[Corollary~8]{Corti2016}. Their effect on the intersection map
is elementary.

\begin{lemma}\label{lem:brauer-cascade-invariance}
Let \(X\) be a normal projective surface with quotient singularities and
minimal resolution \(Y\to X\), where \(Y\) is a rational surface, and let \(X'\to X\) be the blow-up of a
smooth point. This blow-up preserves the Brauer group:
\[
\Br(X')\cong\Br(X).
\]
The same conclusion holds for a finite sequence of blow-ups at smooth points,
including infinitely near centers.
\end{lemma}

\begin{proof}
On minimal resolutions, the induced morphism
\[
\beta:\widetilde X'\longrightarrow\widetilde X
\]
is the blow-up of a point disjoint from the exceptional curves
\(E_1,\dots,E_k\) of \(\widetilde X\to X\), since the resolution is an
isomorphism over the smooth locus. The strict
transforms of these curves give natural identifications
\[
\Lambda_{X'}\cong\Lambda_X,
\qquad
\Lambda_{X'}^\vee\cong\Lambda_X^\vee.
\]
If \(D\) is the new exceptional curve, the Picard group decomposes as
\[
\Pic(\widetilde X')
=
\beta^*\Pic(\widetilde X)\oplus\Z D,
\qquad
D\cdot E_i=0.
\]
This decomposition implies that the intersection numbers satisfy
\[
(\beta^*L+nD)\cdot E_i=L\cdot E_i
\]
for every \(L\in\Pic(\widetilde X)\) and every \(i\). Under the displayed identification of the dual exceptional
lattices, the image and cokernel of \(\theta\) are therefore unchanged.
The Brauer groups are these cokernels by \cref{cor:brauer-cokernel};
naturality of Bright's sequence identifies this isomorphism with pullback.
The assertion for a sequence follows step by step, even when a center
lies on a curve introduced by an earlier blow-up: that curve lies over
the smooth locus, not over a singular point of the current surface.
\end{proof}

Combining this lemma with the marked cascade presentations supplied by
\cite[Corollary~8]{Corti2016} gives, for every \(X\in X_{k,d}\),
\begin{equation}\label{eq:brauer-cascade}
\Br(X)\cong\Br(Z_k).
\end{equation}
It is therefore enough to compute the cokernel on the rigid base
\(Z_k\), or on a surface connected to it by blow-ups at smooth points.

\subsection{\texorpdfstring{The cases \(k=3,4,5\)}{The cases k=3,4,5}}
\label{app:brauer-unobstructed}

% TODO (human writing/digest): Review the passage from the explicit toric
% and cyclic models to their base families, especially the separate
% one-blow-up comparison needed for the four-point model.

\paragraph{Three points: a toric surface.}
The three-point base admits a direct toric calculation.
\phantomsection\label{app:brauer-three}
Corti--Heuberger exhibit polygon \(20\) as a toric \(\mathbb Q\)-Gorenstein degeneration in the
family \(X_{3,5}\)
\cite[Section~7, Table~4 and Figure~11]{Corti2016}. We check its
singularities directly before using uniqueness of the base. Let
\(N\cong\Z^2\), and let
\(v_1,\dots,v_5\in N\) be the primitive ray generators of its fan, in
counterclockwise order, with \(V_i\) the corresponding torus-invariant
divisors:
\[
v_1=(1,0),\quad
v_2=(-1,3),\quad
v_3=(-2,3),\quad
v_4=(-1,0),\quad
v_5=(0,-1).
\]
The first three two-dimensional cones have determinant \(3\), while the
remaining two are smooth. The primitive rays inserted to resolve the three
singular cones are
\[
u_{12}=(0,1),\qquad
u_{23}=(-1,2),\qquad
u_{34}=(-1,1).
\]
Let \(E_1,E_2,E_3\) be the corresponding exceptional torus-invariant
curves. Each inserted ray is one third of the sum of its original
neighbors, and the determinants of all adjacent rays in the subdivided
fan equal one. Thus the subdivision is the minimal toric resolution,
obtained by the ray-subdivision procedure of
\cite[Chapters~10--11]{CoxLittleSchenck2011}. Computing the boundary
self-intersections from the fan relations as in
\cite[Theorem~10.4.4]{CoxLittleSchenck2011} gives the following table;
the last column records \(w_{i-1}+w_{i+1}=b_iw_i\), with square
\(-b_i\) for the corresponding curve:
\[
\begin{array}{c|r|l}
\text{curve}&\text{square}&\text{ray relation}\\ \hline
V_1&0&v_5+u_{12}=0\\
E_1&-3&v_1+v_2=3u_{12}\\
V_2&-1&u_{12}+u_{23}=v_2\\
E_2&-3&v_2+v_3=3u_{23}\\
V_3&-1&u_{23}+u_{34}=v_3\\
E_3&-3&v_3+v_4=3u_{34}\\
V_4&-1&u_{34}+v_5=v_4\\
V_5&0&v_4+v_1=0
\end{array}
\]
The inserted curves are therefore precisely the exceptional \((-3)\)-curves,
so the surface has exactly three \(\frac13(1,1)\)-points and no other
singularities.
After the unimodular change of coordinates
\((x,y)\mapsto(x,x+y)\), these are the rays of the polygon cited above.
It belongs to \(X_{3,5}\) (no higher Fano index family has three singular
points) and hence is isomorphic to \(Z_3\) by the
uniqueness statement recalled in \cref{sec:cascade}.

The boundary intersections give
\[
\theta(V_1)=e_1,\qquad
\theta(V_2)=e_1+e_2,\qquad
\theta(V_4)=e_3.
\]
Consequently, the dual basis vectors satisfy
\[
e_1=\theta(V_1),\qquad
e_2=\theta(V_2)-\theta(V_1),\qquad
e_3=\theta(V_4),
\]
so \(\theta\) is surjective and the Brauer group vanishes:
\[
\Br(Z_3)=\Br(X_{3,5})=0.
\]
Equivalently, the divisor classes
\[
\OO(-V_1),\qquad
\OO(V_1-V_2),\qquad
\OO(-V_4)
\]
have intersection \(-\delta_{ij}\) with \(E_j\).

\begin{corollary}[KKS decomposition in the three-point cascade]
\label{rem:three-point-toric-descent}
Let \(X\in X_{3,d}\) and \(s=5-d\). Then
\[
\Db(\coh X)=\left\langle
\Db(K(3,1)\text{-mod}),
\Db(K(3,1)\text{-mod}),
\Db(K(3,1)\text{-mod}),F_1,\ldots,F_{2+s}
\right\rangle,
\]
where every \(F_i\) is exceptional.
\end{corollary}

\begin{proof}
The toric base \(Z_3\) has torsion-free class group. Of its five
fixed points, the first three in the displayed fan are quotient
singularities and the last two are smooth. KKS
\cite[Proposition~5.6, Remark~5.8, and Corollary~5.10]{KKS2020}
therefore gives three untwisted adherent subcategories on the resolution
and two components generated by single exceptional objects. Here \(M\) in KKS Remark~5.8 is the
correcting line bundle on \(\widetilde Z_3\) used in tensoring with
\(M\otimes\OO(K_{\widetilde Z_3})\):
its degree is zero on every exceptional curve except the final curve
over the last fixed point, where it is \(-1\). Since that fixed point
is smooth, the exception is absent and one takes \(M=\OO\).
By \cite[Corollary~8(3)]{Corti2016}, \(X\) is obtained by
\(s\) blow-ups at smooth points of \(Z_3\), at distinct centers by
\eqref{eq:cascade-distinct}. On resolutions these avoid all three
\((-3)\)-curves. Lemma~\ref{lem:kks-blowups} preserves the three
adherent subcategories under pullback, and Orlov's decomposition with
\(\OO_D\) on the right adds \(s\) exceptional objects. KKS descent and
the identification of the descended adherent components with
\(\Db(K(3,1)\text{-mod})\)
\cite[Theorems~2.12(i) and~3.16]{KKS2020} give the stated decomposition.
\end{proof}

\begin{remark}[Toric surfaces and toric degenerations]
\label{rem:toric-bases-versus-degenerations}
Among nonsmooth del Pezzo surfaces with only \(\frac13(1,1)\)-points,
the toric surfaces are precisely those in
\[
S_{1,25/3},\quad X_{1,22/3},\quad X_{1,19/3},\quad
X_{3,5},\quad X_{3,4},\quad X_{6,2}.
\]
These are the toric cases in the Corti--Heuberger classification
\cite[Tables~1, 2, and~4]{Corti2016}; see also
\cite[Table~1]{OnetoPetracci2018}. In particular, \(Z_k\) is toric
precisely for \(k\in\{1,3,6\}\).

The toric degenerations of \(Z_2,Z_4,Z_5\) in
\cite[Table~4, polygons~23, 11, and~7]{Corti2016} have additional
\(T\)-singularities: respectively an \(A_1\)-point, an \(A_2\)-point,
and two \(\frac14(1,1)\)-points. Thus these toric special fibers
are not the rigid bases themselves. The semiorthogonal decompositions for
toric surfaces in \cite[Section~5]{KKS2020} describe the derived categories
of these special fibers.
\end{remark}

\paragraph{The \(k\)-gon construction.}
For the four- and five-point calculations we use the \(k\)-gon construction.
It starts with a smooth del Pezzo surface of degree \(k\) containing a cyclic
configuration \(C_1+\cdots+C_k\) of \((-1)\)-curves, blows up the vertices
\(C_i\cap C_{i+1}\), and contracts the strict transforms of the \(C_i\).
Corti--Heuberger show that the resulting surface has \(k\) singularities of
type \(\frac13(1,1)\), degree \(k/3\), and belongs to the family
\(X_{k,k/3}\) \cite[Section~3.5, ``\(k\)-gons'']{Corti2016}. Thus the
surfaces used below are respectively members of \(X_{4,4/3}\) and
\(X_{5,5/3}\). Let \(D_i\) be the exceptional curve over the
\(i\)-th vertex and \(E_i\) the strict transform of \(C_i\), with cyclic
indices. Direct intersection calculations give
\(E_i^2=-3\), since each \(C_i\) passes through two vertices.
The curves \(E_i\) are pairwise disjoint, and each \(D_i\) meets just
\(E_i\) and \(E_{i+1}\), transversally once. Consequently,
\[
\theta(D_i)=e_i+e_{i+1},
\qquad
\theta(E_i)=-3e_i.
\]
The degree also follows from the discrepancy computation in
the proof of \cref{cor:mixed-basket-length}: the \(k\) blow-ups lower
\(K^2\) from \(k\) to zero, and contraction of the \(k\) disjoint
\((-3)\)-curves raises it to \(k/3\).

\paragraph{Five points.}
For the five-point base, consider the degree-five del Pezzo surface
\(\Bl_{p_1,\dots,p_4}\Pp^2\), with exceptional classes \(Q_i\) and line
class \(H\). The curves
\[
Q_1,\quad H-Q_1-Q_2,\quad Q_2,\quad H-Q_2-Q_3,\quad H-Q_1-Q_4
\]
have adjacent intersections one and nonadjacent intersections zero, so
they form a cyclic pentagon of \((-1)\)-curves. Blowing up its vertices
and contracting the strict transforms of its sides gives a member of \(X_{5,5/3}\), which is the rigid base \(Z_5\) by
uniqueness. The five vectors
\(e_i+e_{i+1}\) are the columns of
\[
A_5=I+P_5,
\]
where \(P_5\) is the cyclic permutation matrix. A direct calculation gives
\[
\det(A_5)=\prod_{\zeta^5=1}(1+\zeta)=2.
\]
The quotient \(\Z^5/A_5\Z^5\) has order \(2\). Since the image of \(\theta\) also
contains \(3\Z^5\), and multiplication by \(3\) is invertible on this
order-two quotient, we have
\[
A_5\Z^5+3\Z^5=\Z^5.
\]
Hence \(\theta\) is surjective and the Brauer group vanishes:
\[
\Br(Z_5)=\Br(X_{5,5/3})=0.
\]

\paragraph{Four points.}
For the four-point cascade, it is convenient to use a member of \(X_{4,4/3}\), obtained from the
rigid base \(Z_4=X_{4,7/3}\) by one blow-up at a smooth point
\cite[Corollary~8]{Corti2016}. By \cref{lem:brauer-cascade-invariance},
it has the same cokernel as \(Z_4\). The vectors \(e_i+e_{i+1}\) span a
rank-three subgroup, so one additional divisor class is needed. To avoid
confusing the exceptional classes on a degree-four del Pezzo surface with the
\((-3)\)-curves above, write
\[
S=\Bl_{p_1,\dots,p_5}\Pp^2
\]
with basis \(H,Q_1,\dots,Q_5\), where \(H^2=1\), \(Q_i^2=-1\), and all
other pairings between basis elements vanish. Consider
\[
\begin{aligned}
C_1&=Q_1,\\
C_2&=H-Q_1-Q_2,\\
C_3&=Q_2,\\
C_4&=2H-Q_1-Q_2-Q_3-Q_4-Q_5.
\end{aligned}
\]
Their self-intersections are all \(-1\). The adjacent intersections are
\[
C_1\cdot C_2=C_2\cdot C_3=C_3\cdot C_4=C_4\cdot C_1=1,
\]
and the nonadjacent intersections are
\[
C_1\cdot C_3=0,
\qquad
C_2\cdot C_4=0.
\]
Their intersection graph is therefore a \(4\)-cycle. In addition,
\[
Q_3\cdot C_1=Q_3\cdot C_2=Q_3\cdot C_3=0,
\qquad
Q_3\cdot C_4=1.
\]
For general \(p_1,\dots,p_5\), these intersections are transverse; in
particular, \(Q_3\) avoids the four vertices of the cycle. If \(\pi\) is their
blow-up, the projection formula gives
\[
\theta(\pi^*Q_3)=e_4.
\]
Using this as the starting dual basis vector, the remaining basis vectors are
obtained successively:
\[
\begin{aligned}
e_1&=\theta(D_4)-e_4,\\
e_2&=\theta(D_1)-e_1,\\
e_3&=\theta(D_2)-e_2.
\end{aligned}
\]
These identities make \(\theta\) surjective. Since the \(4\)-gon representative
is a one-point blow-up of \(Z_4\),
\cref{lem:brauer-cascade-invariance} gives
\[
\Br(Z_4)=\Br(X_{4,7/3})\cong\Br(X_{4,4/3})=0.
\]

\subsection{The six-point Brauer obstruction}
\label{app:brauer-six}

Unlike the five smaller cascades, the six-point cascade carries a
nonzero Brauer obstruction. We first exhibit a relation among the toric
boundary divisors that produces an order-three class in the saturation quotient
\(\Lambda^{\mathrm{sat}}/\Lambda\subset A_\Lambda\). This proves
that \(\Br(Z_6)\ne0\). We then compute
the full image of the intersection map: a single alternating-sum
congruence modulo three describes it, showing that the Brauer group is
exactly \(\Z/3\).

For the rigid base \(Z_6=X_{6,2}\) of \cref{sec:cascade},
Corti--Heuberger describe its fan in
\cite[Section~2.2.5]{Corti2016}. We compute the intersection map by first
describing the minimal toric resolution and then determining the lattice
generated by the images of its boundary divisors. Let \(N\cong\Z^2\), and let
\(v_1,\dots,v_6\in N\) be the primitive ray generators of its fan, with
\(V_i\) the corresponding torus-invariant divisors:
\[
\begin{aligned}
v_1&=(1,0),&
v_2&=(-1,3),&
v_3&=(-2,3),\\
v_4&=(-1,0),&
v_5&=(1,-3),&
v_6&=(2,-3),
\end{aligned}
\]
in counterclockwise order. The unimodular change
\((x,y)\mapsto(x,x+y)\) sends these vertices to those of
\cite[Figure~1]{Corti2016}. Indices are cyclic modulo \(6\) throughout. Every
adjacent cone has determinant \(3\), and hence defines a cyclic quotient
singularity of order \(3\). The primitive rays inserted in the
minimal toric resolution are
\[
u_i=\frac{v_i+v_{i+1}}3,
\qquad
i\in\Z/6\Z.
\]
The inserted rays are explicitly
\[
\begin{aligned}
u_1&=(0,1),&
u_2&=(-1,2),&
u_3&=(-1,1),\\
u_4&=(0,-1),&
u_5&=(1,-2),&
u_6&=(1,-1).
\end{aligned}
\]
Let \(E_i\) be the exceptional torus-invariant curve corresponding to
\(u_i\). Since
\[
\det(v_i,u_i)=\det(u_i,v_{i+1})=1,
\]
the subdivision is smooth and gives the minimal toric resolution
\cite[Chapters~10--11]{CoxLittleSchenck2011}. The resolved boundary is the
alternating cycle
\[
V_1,E_1,V_2,E_2,\dots,V_6,E_6,
\]
and the ray relations are
\[
u_{i-1}+u_i=v_i,
\qquad
v_i+v_{i+1}=3u_i.
\]
Therefore, \cite[Theorem~10.4.4]{CoxLittleSchenck2011} gives
\[
V_i^2=-1,\qquad E_i^2=-3.
\]
Consequently, the minimal resolution of each singularity consists of one
\((-3)\)-curve, so all six singularities are of type \(\frac13(1,1)\).

\paragraph{A nonzero class in the saturation quotient.}
An order-three class in the saturation quotient comes from a divisibility relation. The principal-divisor
formula
\(\operatorname{div}(\chi^m)=\sum_\rho\langle m,v_\rho\rangle D_\rho\)
of \cite[Theorem~4.1.3]{CoxLittleSchenck2011}, applied to the character
\(m=(0,1)\), gives
\[
\begin{aligned}
0={}&E_1+3V_2+2E_2+3V_3+E_3-E_4\\
&\hspace{3em}-3V_5-2E_5-3V_6-E_6.
\end{aligned}
\]
Rearranging gives
\begin{equation}\label{eq:brauer-six-glue}
E_1-E_2+E_3-E_4+E_5-E_6=3L,
\end{equation}
where the divisor \(L\) is
\[
L=-V_2-V_3+V_5+V_6-E_2+E_5
\in\Pic(\widetilde Z_6).
\]
Thus \(L\) lies in \(\Lambda^{\mathrm{sat}}\) but not in \(\Lambda\):
the \(E_i\) are linearly independent, and its coefficients in this basis
are \(\pm1/3\). Its class has order three and realizes the alternating
pattern allowed by \cref{lem:brauer-glue-vectors}.
Moreover, \(L^2=-2\), confirming the isotropy
of its class for the discriminant pairing in
\cref{def:brauer-glue}.
In particular, \cref{lem:brauer-glue-duality} already implies that
\(\Br(Z_6)\ne0\). We next determine the whole group.

\paragraph{The cokernel.}
For a complete toric variety, the torus-invariant prime divisors generate
\(\Cl(\widetilde Z_6)\), and smoothness gives
\(\Cl(\widetilde Z_6)=\Pic(\widetilde Z_6)\)
\cite[Theorem~4.1.3]{CoxLittleSchenck2011}. Hence the toric boundary generates
\(\Pic(\widetilde Z_6)\), and its image under \(\theta\) is generated by
\[
\theta(E_i)=-3e_i,
\qquad
\theta(V_j)=e_{j-1}+e_j,
\]
with cyclic indices. Put \(c_j=e_{j-1}+e_j\), where \(e_0=e_6\), and let
\(A\subset\Z^6\) be the subgroup generated by \(c_2,\dots,c_6\).
These five vectors, together with \(e_6\), form a basis of \(\Z^6\):
the matrix of their columns is lower triangular with diagonal entries one.
Thus \(A\) is primitive of rank five. Moreover,
\[
c_1=c_2-c_3+c_4-c_5+c_6,
\]
so \(A\) is also the span of all six images \(\theta(V_j)\). Consider the
surjective homomorphism
\[
\lambda:\Z^6\longrightarrow\Z,
\qquad
(a_1,\dots,a_6)\longmapsto
a_1-a_2+a_3-a_4+a_5-a_6.
\]
Every \(c_j\) lies in \(\ker(\lambda)\). Since both
\(A\) and \(\ker(\lambda)\) have rank five and \(A\) is primitive,
this inclusion is an equality. The image of the full toric boundary is
therefore
\[
\operatorname{im}(\theta)
=
A+3\Z^6
=
\ker(\lambda)+3\Z^6.
\]
This sum is the kernel of the reduction of \(\lambda\) modulo \(3\). Indeed,
if \(\lambda(a)=3q\), then \(a-3q e_1\in\ker(\lambda)\), while the reverse
inclusion is immediate. Consequently,
\[
\operatorname{im}(\theta)
=
\left\{
(a_1,\dots,a_6)\in\Z^6:
a_1-a_2+a_3-a_4+a_5-a_6\equiv0\pmod3
\right\}.
\]
The homomorphism \(\lambda\bmod3:\Z^6\to\Z/3\) is surjective, so the
cokernel and Brauer group are
\[
\operatorname{coker}(\theta)\cong\Z/3,
\qquad
\Br(Z_6)=\Br(X_{6,2})\cong\Z/3.
\]
By \cref{lem:brauer-glue-duality}, the saturation quotient also has order three.
It is therefore generated by the class of \(L\) in
\eqref{eq:brauer-six-glue}; no support-three class in the saturation quotient occurs on \(Z_6\).
Indeed, its two nonzero elements are \([L]\) and \(-[L]\), whose
coefficient vectors modulo three are
\(\pm(1,-1,1,-1,1,-1)\); both have all six entries nonzero.
Here the global toric divisor relations select the alternating
support-six class, whereas the computations for \(k=3,4,5\) exclude
even the numerically permitted support-three classes.

For the adherence convention we use classes with
\(L_i\cdot E_j=-\delta_{ij}\); this differs from the dual-basis convention
of \cref{cor:brauer-cokernel} only by replacing each \(L_i\) by \(-L_i\).
We keep this sign convention in the twisted six-point construction below.

\paragraph{The obstruction to untwisted adherence.}
In particular, there cannot exist divisor classes \(L_i\) satisfying
\[
L_i\cdot E_j=-\delta_{ij},
\]
since such classes would make \(\theta\) surjective.

By \cref{thm:kks-obstruction}, an untwisted adherent decomposition would
force \(\Cl(X_{6,2})\) to be torsion-free, equivalently
\(\Br(X_{6,2})=0\). The computed nonzero class rules this out.
The KKS result for toric surfaces instead gives a semiorthogonal decomposition
of a Brauer-twisted category on \(Z_6\). Under the cascade blow-ups, the
Brauer class is pulled back and the corresponding twisted adherent
decomposition is preserved, as proved below.

\begin{corollary}[KKS decomposition of the twisted category in the six-point cascade]
\label{cor:six-point-twisted-descent}
Let \(X\in X_{6,d}\), let \(\eta:X\to Z_6\) be a marked cascade
presentation, and put \(s=2-d\). There is a generator
\(\alpha\in\Br(Z_6)\) such that, with
\(\alpha_X=\eta^*\alpha\),
\[
\Db(\coh(X,\alpha_X))=
\left\langle
\underbrace{\Db(K(3,1)\text{-mod}),\ldots,
\Db(K(3,1)\text{-mod})}_{6\ \text{components}},
F_1,\ldots,F_s
\right\rangle.
\]
Here \(\coh(X,\alpha_X)\) uses the KKS convention of coherent
right modules over an Azumaya algebra representing \(\alpha_X\),
and the \(F_i\) are
exceptional objects of its derived category; their list is empty
when \(s=0\). The class \(\alpha_X\) is nonzero.
\end{corollary}

\begin{proof}
The proof has three steps: we first identify the twisted adherence data on the
resolution of \(Z_6\), then pull these data through the cascade blow-ups using
\cref{lem:kks-blowups}, and finally identify the descended twisted components
via the KKS twisted pushforward.

Write \(\pi:\widetilde Z_6\to Z_6\) for the resolution and use
the counterclockwise boundary order \(V_1,E_1,\ldots,V_6,E_6\)
above. In KKS notation, \(E_{i,0}=V_i\), \(E_{i,1}=E_i\),
\(m_i=1\), and \(d_{i,1}=-E_i^2=3\); the last fixed point
corresponds to the cone between \(v_6\) and \(v_1\).
\cite[Proposition~5.6, equations~(5.4)--(5.6), p.~510]{KKS2020}
gives the component
\(\langle\mathcal L_{i,0},\mathcal L_{i,0}(E_i)\rangle\), where
\[
\mathcal L_{i,0}=\OO\Bigl(\sum_{j<i}(V_j+E_j)+V_i\Bigr),
\qquad
\deg(\mathcal L_{i,0}|_{E_i})=1+\delta_{i6}.
\]
Indeed, \(V_i\) meets \(E_i\) once, and the sum contributes only
\(V_1\cdot E_6=1\) when \(i=6\). Our adherence convention
\(\deg(\mathcal L_{i,0}|_{E_i})=3+b_i-1\) therefore gives
\(b_i=-1+\delta_{i6}\), exactly
\cite[Proposition~5.6, equation~(5.7)]{KKS2020}.
Adjunction gives \(K_{\widetilde Z_6}\cdot E_i=1\), so tensoring
with \(\OO(K_{\widetilde Z_6})\) adds one to every adherence parameter:
\[
b=(-1,-1,-1,-1,-1,0),\qquad
b+\theta(K_{\widetilde Z_6})=e_6.
\]
We also fix the Brauer sign: KKS's map \(B\) sends the degree vector
of a splitting bundle \(\mathcal V\), with
\(\pi^*\mathcal R\cong\operatorname{End}(\mathcal V)\), to
\([\mathcal R]\), not \([\mathcal R^{\mathrm{op}}]\).
The twisted pushforward is
\(F\mapsto R\pi_*(F\otimes\mathcal V^\vee)\), a right
\(\mathcal R\)-module, as in
\cite[Proposition~4.9 and equations~(4.13)--(4.14), pp.~498--500]{KKS2020}.
Since \(B\circ\theta=0\), the displayed equality gives
\(\alpha=B(b)=B(e_6)\), agreeing with
\cite[equation~(5.8) and Theorem~5.9]{KKS2020}.
Finally \(\lambda(e_6)=-1\), so this class generates \(\Br(Z_6)\).

There are no smooth fixed points, so the decomposition of
\(\Db(\coh(Z_6,\alpha))\) has six \(K(3,1)\)-components and no components
generated by single exceptional objects. Under the blow-ups
\(\widetilde X\to\widetilde Z_6\),
Lemma~\ref{lem:kks-blowups} preserves the adherent components and their
adherence parameters. Orlov's decomposition with \(\OO_D\) on the right adds
\(s\) components generated by single exceptional objects. The intersection-map description in
\cref{lem:brauer-cascade-invariance} identifies the resulting
Brauer class with \(\eta^*\alpha\), still a generator.
More explicitly, a vector bundle on \(\widetilde Z_6\) realizing
the adherence parameters, as in \cite[Corollary~4.10(1)]{KKS2020}, pulls back to
one with the same restrictions on the six exceptional curves.

Applying the twisted pushforward of \cite[equation~(4.14)]{KKS2020}
and then \cite[Theorem~4.19]{KKS2020} to each adherent component
identifies its image with \(\Db(K(3,1)\text{-mod})\).
For the additional Orlov objects, the restrictions to the six
exceptional curves are zero, so \cite[Lemma~4.14]{KKS2020}
identifies each with the twisted pullback of its pushforward.
Full faithfulness of twisted pullback preserves their exceptionality;
adjunction preserves their semiorthogonality to the preceding components.
Fullness follows from \cite[Lemma~4.13]{KKS2020} and the full
resolution decomposition. Thus the components generated by single exceptional objects require no
line-bundle adherence hypothesis. This is twisted descent, not
ordinary pushforward under \cref{thm:kks-descent}.
\end{proof}

\subsection{\texorpdfstring{Proof of \cref{thm:brauer-cascades-intro}}{Proof of Theorem 1.2}}
\label{app:brauer-classification}

% TODO (human writing/digest): Review the final assembly of the proof,
% keeping base uniqueness, blow-up at a smooth point invariance, and the
% coverage of the auxiliary families logically separate.

We first collect the computations on the bases. For \(k=1,2\),
\cref{cor:brauer-one-two} gives vanishing without choosing a blow-up presentation.
For \(k=3,5\), \cref{app:brauer-unobstructed} computes the Brauer group
on a toric surface and a pentagon surface in the respective base families.
Each is the corresponding \(Z_k\), by uniqueness of the base in
\cref{sec:cascade}. For \(k=4\), the same subsection uses a quadrilateral
surface in \(X_{4,4/3}\), not in the base family; its contraction
to a smooth point of \(Z_4\) and \cref{lem:brauer-cascade-invariance} give
\(\Br(Z_4)=0\). Finally, the toric computation in
\cref{app:brauer-six} gives \(\Br(Z_6)\cong\Z/3\). Altogether,
\[
\Br(Z_k)
\cong
\begin{cases}
0, & 1\leq k\leq5,\\
\Z/3, & k=6.
\end{cases}
\]

For every surface \(X\) in a Fano index one family \(X_{k,d}\),
\eqref{eq:brauer-cascade} now gives \(\Br(X)\cong\Br(Z_k)\).

The auxiliary families \(B_{1,16/3}\) and \(B_{2,8/3}\) are covered
directly by \cref{cor:brauer-one-two}, since their members have one and
two singular points, respectively. Their Brauer groups can also be
computed by invariance under blow-ups at smooth points: the alternative
sequences of blow-ups and floating contractions in \cite[Remark~9]{Corti2016}, recalled in
\cref{sec:cascade}, give a blow-up at a smooth point \(X'\to B\) and four
floating contractions from \(X'\) to \(Z_k\), for \(k=1,2\) respectively.
Reversing those contractions and applying
\cref{lem:brauer-cascade-invariance} gives
\[
\Br(B)\cong\Br(X')\cong\Br(Z_k)=0.
\]
Together with the already treated base \(Z_1\), these account for all
higher Fano index families.

We have now covered every family in the six cascades. The final assertion
follows from \cite[Corollary~4.8]{KKS2020}: the nonzero Brauer group in the
six-point cascade obstructs an untwisted adherent decomposition, whereas this
Brauer obstruction is absent in the other five cascades. This proves
\cref{thm:brauer-cascades-intro}. \qed

% ------------------------------------------------------------

\section{Explicit exceptional collections on the rigid cascade bases}
\label{app:explicit-bases}

This appendix constructs marked blow-up presentations of the minimal
resolutions \(\widetilde Z_k\) and uses Orlov's formula to write the
corresponding exceptional collections, proving
\cref{prop:explicit-rigid-bases}.
The descriptions of the singular surfaces come from Corti--Heuberger and
Appendix~\ref{app:brauer}; the blow-down sequences and the elementary
transformation below express their resolutions as successive point blow-ups.
Unlike the exceptional-lattice computations in Appendix~\ref{app:brauer},
these presentations specify the ordered centers needed for Orlov's formula.

After fixing uniform notation, we treat the six bases in order.
\Cref{app:explicit-base-one,app:explicit-base-two} begin with \(\F_3\)
and distinguish the incidence conditions that create a second quotient
point. The three- and six-point cases are checked by toric contractions;
the four-point case requires an elementary transformation, and the
five-point case uses the pentagon construction. Each case ends with its
ordered exceptional collection on the resolution, and
\cref{cor:explicit-cascade-collections} extends these collections along the
successive blow-ups of a chosen marked cascade presentation of a surface
of Fano index one.

The following guide records the centers and geometric steps for the resolutions
in Table~\ref{tab:rigid-base-invariants}. All modifications are ordinary
blow-ups at smooth points; ``infinitely near'' describes the location of a
center relative to the starting surface \(Y_0\) listed there.
\begin{center}
\small
\renewcommand{\arraystretch}{1.2}
\begin{tabular}{@{}c p{0.46\linewidth} p{0.42\linewidth}@{}}
\(k\)&Centers&Key geometric step\\ \hline
1&None&Use the negative section.\\
2&One point and two infinitely near points&Create a second \((-3)\)-curve.\\
3&One torus-fixed point and three infinitely near torus-fixed points&Reverse
four boundary contractions.\\
4&Five distinct points and two infinitely near points&Perform an elementary
transformation.\\
5&Four plane points, then five vertices (four infinitely near)&Separate the
pentagon curves.\\
6&Four torus-fixed points and four infinitely near torus-fixed points&Contract to
a hexagon, then a square.
\end{tabular}
\end{center}
The individual cases below specify the centers and their order.

We use the following notation throughout. For a sequence of ordinary
point blow-ups
\[
Y_m\xrightarrow{\beta_m}Y_{m-1}\longrightarrow\cdots
\xrightarrow{\beta_1}Y_0,
\qquad \sigma=\beta_1\circ\cdots\circ\beta_m,
\]
let \(D_i\subset Y_i\) be the exceptional curve created by \(\beta_i\),
and set
\begin{equation}\label{eq:explicit-orlov-objects}
G_i=L\beta_m^*\cdots L\beta_{i+1}^*\OO_{D_i}(-1).
\end{equation}
Thus \(G_m=\OO_{D_m}(-1)\). If
\(\langle L_1,\ldots,L_t\rangle\) is a full exceptional collection of
line bundles on \(Y_0\), \cref{thm:orlov-blow-up} gives the ordered
full exceptional collection
\begin{equation}\label{eq:explicit-orlov-order}
\left\langle G_m,\ldots,G_1,
\sigma^*L_1,\ldots,\sigma^*L_t\right\rangle.
\end{equation}
In particular, a blow-up center may be infinitely near to an earlier point;
the objects in
\eqref{eq:explicit-orlov-objects} are not replaced by sheaves on the
final strict transforms. The collections on \(Y_0\), recalled in
\cref{sec:orlov}, are
\[
\begin{aligned}
\mathscr B_{\Pp^2}
 &=\langle\OO,\OO(H),\OO(2H)\rangle,\\
\mathscr B_{\F_n}
 &=\langle\OO,\OO(f),\OO(C_0+nf),\OO(C_0+(n+1)f)\rangle.
\end{aligned}
\]
Here \(H\) is a line class, and \(C_0,f\) are the section and fiber
classes with \(C_0^2=-n\). We denote the collection on
\(\widetilde Z_k\) by \(\mathscr A_k\), the canonical stack of
\(Z_k\) by \(\mathcal Z_k\), and its Ishii--Ueda embedding by
\(\Phi_k\). In each case, \cref{cor:mixed-basket-length} turns
\(\mathscr A_k\) into a full exceptional collection on \(\mathcal Z_k\) by applying \(\Phi_k\)
and prepending the \(k\) local objects. The lengths are recorded
in Table~\ref{tab:rigid-base-invariants}, so we list only the collections
on the resolutions below.

\subsection{The one-point base}
\label{app:explicit-base-one}

The minimal resolution of \(Z_1=\Pp(1,1,3)\) is \(\F_3\), with
exceptional curve \(C_0\). No point blow-ups are needed, and we take
\[
\mathscr A_1=
\left\langle\OO,\OO(f),\OO(C_0+3f),\OO(C_0+4f)\right\rangle.
\]

\subsection{The two-point base}
\label{app:explicit-base-two}

The directed minimal model program of \cite[Section~6]{Corti2016} contains a two-point branch
ending at \(\Pp(1,1,3)\):
\[
Z_2=X_{2,17/3}\longrightarrow T\longrightarrow\Pp(1,1,3),
\qquad \Sing(T)=\left\{\tfrac13(1,1),A_1\right\}.
\]
For the exceptional collection we need a marked blow-up presentation of the
minimal resolution, which we now construct directly.

\begin{proposition}\label{prop:explicit-base-two}
Choose \(q_1\in\F_3\setminus C_0\), and let \(f_0\) be its ruling
fiber. Blow up \(q_1\), with exceptional curve \(D_1\), and then choose
\(q_2\in D_1\) away from the strict transform of \(f_0\). Blow up
\(q_2\), with exceptional curve \(D_2\). Finally, choose \(q_3\) on
the strict transform of \(D_1\), away from both \(D_2\) and the strict
transform of \(f_0\), and blow it up. Then
\[
\widetilde Z_2\simeq
Y_3=\Bl_{q_3}\Bl_{q_2}\Bl_{q_1}\F_3.
\]
The two curves contracted to the quotient singularities are the untouched
\(C_0\) and the final strict transform \(A\) of \(D_1\).
\end{proposition}

\begin{proof}
Write \(Y_i\) for the surface after the first \(i\) blow-ups and
\(T_i\) for the total transform of the class of \(D_i\) on \(Y_3\).
Then \(T_iT_j=-\delta_{ij}\), and pullbacks from \(\F_3\) are
orthogonal to all \(T_i\). On \(Y_2\) the strict transform
\(A_2\) of \(D_1\) has square \(-2\), since the second blow-up is centered
at \(q_2\in D_1\). The last center lies on
\(A_2\), but on neither of its two neighboring curves. Consequently
\[
A=T_1-T_2-T_3,\qquad A^2=-3,\qquad C_0A=0.
\]
The strict transform \(F'=f-T_1\) of \(f_0\) has square \(-1\);
it meets \(C_0\) and \(A\) once each. The curves \(D_2,D_3\)
remain \((-1)\)-curves, each meeting \(A\) once and disjoint from
\(C_0\) and \(F'\). The dual graph records these transverse
intersections (each edge has intersection multiplicity one):
\[
\begin{array}{ccccc}
&&&&D_2\,(-1)\\[-2pt]
&&&&|\\[-2pt]
C_0\,(-3)&\text{---}&F'\,(-1)&\text{---}&A\,(-3)\\[-2pt]
&&&&|\\[-2pt]
&&&&D_3\,(-1)
\end{array}
\]

To check that contracting \(C_0\) and \(A\) gives the required del
Pezzo surface, consider the rational divisor
\[
P=-K_{Y_3}-\tfrac13(C_0+A)
 =\tfrac53 C_0+5f-\tfrac43T_1-\tfrac23T_2-\tfrac23T_3.
\]
The intersection form in the orthogonal total-transform basis gives
\[
P^2=\tfrac{17}{3},\qquad
PC_0=PA=0,\qquad PF'=\tfrac13,\qquad PD_2=PD_3=\tfrac23.
\]
Every other irreducible curve is the strict transform of a fiber not
through \(q_1\), or of a curve \(B\sim aC_0+bf\) with
\(a\geq1\): an irreducible curve on \(\F_3\) with \(a=0\) is a
fiber, and curves exceptional over \(\F_3\) were already listed.
Since \(B\ne C_0\), nonnegative intersection of
distinct irreducible curves gives \(BC_0=b-3a\geq0\), hence
\(b\geq3a\). Let \(m_1,m_2,m_3\) be its multiplicities at the
successive centers.
Intersection with \(f_0\) gives \(m_1\leq a\), while nonnegative
intersection of its final strict transform with
\(A=T_1-T_2-T_3\) gives
\(\widetilde B A=m_1-m_2-m_3\geq0\), hence
\(m_2+m_3\leq m_1\). Consequently,
\[
P\cdot\widetilde B
 =\tfrac53b-\tfrac43m_1-\tfrac23(m_2+m_3)
 \geq\tfrac53b-2a\geq3a>0.
\]
On a fiber not through \(q_1\), the intersection is \(5/3\).
Thus \(P\) is nef and big, and its only null curves are \(C_0,A\).
The pair \((Y_3,(C_0+A)/3)\) is klt, so the basepoint-free theorem
\cite[Section~3.2]{Kollar1998} makes a sufficiently divisible multiple
of \(P\) semiample. Its associated contraction has precisely these
two disjoint \((-3)\)-curves as exceptional locus and gives a normal
surface \(Z\) with two \(\frac13(1,1)\)-points. The discrepancy
formula identifies \(P\) with the pullback of the ample divisor
\(-K_Z\).

There is also no ambiguity about the Fano index. The divisor-class group
of \(Z\) is the quotient of \(\Pic(Y_3)\) by \(C_0,A\), so it is
free with basis \(f,T_2,T_3\). In this quotient
\[
-K_Z=5f-2T_2-2T_3
\]
is primitive. The surface therefore has Fano index one, and
\cite[Theorem~4 and Remark~7]{Corti2016} identifies it with \(Z_2\).
Since its exceptional curves have square \(-3\), \(Y_3\) is its
minimal resolution.
\end{proof}

With \(\sigma_2:Y_3\to\F_3\) the composite, the seven-object collection is
\[
\mathscr A_2=\left\langle G_3,G_2,G_1,
\sigma_2^*\OO,\sigma_2^*\OO(f),
\sigma_2^*\OO(C_0+3f),\sigma_2^*\OO(C_0+4f)\right\rangle.
\]

\subsection{The three-point base}
\label{app:explicit-base-three}

For the three-point base, use the fan and boundary notation of
\cref{app:brauer-three}. Its minimal resolution has boundary
\(V_1,E_1,V_2,E_2,V_3,E_3,V_4,V_5\), with self-intersections
\[
(0,-3,-1,-3,-1,-3,-1,0).
\]
To obtain a marked blow-up presentation over \(\F_3\), we contract, in this
order, the successive images of \(V_4,V_3,E_3,E_2\). The following table lists
the neighboring rays at the actual stage of each contraction:
\[
\begin{array}{c|c|c}
\text{curve}&\text{neighboring rays}&\text{ray relation}\\ \hline
V_4&u_{34},v_5&v_4=u_{34}+v_5\\
V_3&u_{23},u_{34}&v_3=u_{23}+u_{34}\\
E_3&u_{23},v_5&u_{34}=u_{23}+v_5\\
E_2&v_2,v_5&u_{23}=v_2+v_5
\end{array}
\]
In every row the determinant of the two neighbors is one. The ray-sum
criterion therefore verifies that the curve being removed is a
\((-1)\)-curve and that its contraction is a smooth toric blow-down
\cite[Theorem~10.4.4]{CoxLittleSchenck2011}. The remaining fan has rays
\(v_1,u_{12},v_2,v_5\), with boundary squares \((0,-3,0,3)\);
it is the fan of \(\F_3\), with negative section corresponding to
\(u_{12}\).

Reversing these four contractions prescribes the centers: insert
\(u_{23}\) in the cone \(\langle v_2,v_5\rangle\), then
\(u_{34}\) in \(\langle u_{23},v_5\rangle\), then \(v_3\) in
\(\langle u_{23},u_{34}\rangle\), and finally \(v_4\) in
\(\langle u_{34},v_5\rangle\). Each center is the intersection of
the two indicated boundary curves. This gives
\(\sigma_3:\widetilde Z_3\to\F_3\) and the full exceptional collection
\[
\mathscr A_3=\left\langle G_4,G_3,G_2,G_1,
\sigma_3^*\OO,\sigma_3^*\OO(f),
\sigma_3^*\OO(C_0+3f),\sigma_3^*\OO(C_0+4f)\right\rangle.
\]
The creation stages for \(G_1,G_2,G_3,G_4\) correspond respectively
to the rays \(u_{23},u_{34},v_3,v_4\). In particular, \(G_1\) and \(G_2\)
are understood as derived pullbacks from their creation stages, not as
sheaves on the final \((-3)\)-curves.

\subsection{The four-point base}
\label{app:explicit-base-four}

The surface obtained from the quadrilateral in
\cref{app:brauer-unobstructed} lies in \(X_{4,4/3}\), not in the
rigid base family \(Z_4=X_{4,7/3}\). The following proposition contracts
an additional \((-1)\)-curve and performs an elementary transformation of
\(\F_1\) to give a marked blow-up presentation of \(\widetilde Z_4\)
over \(\F_2\).

\begin{proposition}\label{prop:explicit-base-four}
Let \(S=\Bl_{p_1,\ldots,p_5}\Pp^2\), with the points general, and
let \(H,Q_1,\ldots,Q_5\) be its standard divisor classes. Set
\[
\begin{aligned}
C_1&=Q_1,& C_2&=H-Q_1-Q_2,\\
C_3&=Q_2,& C_4&=2H-Q_1-Q_2-Q_3-Q_4-Q_5.
\end{aligned}
\]
Blow up the vertices \(r_i=C_i\cap C_{i+1}\), cyclically, with
exceptional curves \(R_i\), obtaining \(Y\). The strict transform
of the line tangent at \(p_1\) to the conic through the five points
is a \((-1)\)-curve
\[
D=H-Q_1-R_4.
\]
Its contraction \(Y\to Y'\) identifies \(Y'\) with
\(\widetilde Z_4\). Moreover, an elementary transformation at \(r_4\)
gives an explicit seven-point blow-up presentation
\(\sigma_4:\widetilde Z_4\to\F_2\). Its centers, in order, are
the transported points \(p_2,p_3,p_4,p_5,r_1,r_2,r_3\): the first
five are distinct on \(\F_2\), and the last two are infinitely near
\(p_2\).
\end{proposition}

\begin{proof}
All divisor classes on \(Y\) in the calculations below are total
pullbacks, except where strict transforms are explicitly specified.
The four \((-1)\)-curves \(C_i\) form a cycle with distinct vertices,
so their strict transforms
\[
E_i=C_i-R_{i-1}-R_i
\]
satisfy \(E_i^2=-3\) and \(E_iE_j=0\) for \(i\ne j\).
The contraction \(Y\to X_{4,4/3}\) is the \(4\)-gon construction
of \cite[Section~3.5]{Corti2016}.

The conic through five general points is smooth. Its tangent line at
\(p_1\) contains none of \(p_2,\ldots,p_5\), by B\'ezout, and
meets \(Q_1\) at \(r_4\) after the first blow-up. This point is
different from \(r_1\), the direction of the secant through \(p_2\).
Thus \(D\) is the strict transform of an irreducible smooth
rational curve. In the orthogonal
blow-up basis we have
\[
D^2=-1,\qquad (-K_Y)D=1,\qquad DE_i=0\quad(1\leq i\leq4).
\]
These distinct irreducible curves are therefore disjoint. The image of
\(D\) on \(X_{4,4/3}\) lies in the smooth locus and is a floating
\((-1)\)-curve. Its contraction is a log del Pezzo surface with four
unchanged \(\frac13(1,1)\)-points and degree \(7/3\), hence is
\(Z_4\) by the classification in \cref{sec:cascade}. Contracting
\(D\) on \(Y\) leaves the four \(E_i\) unchanged and gives its
minimal resolution \(Y'\).

To identify the starting ruled surface, first blow up only \(p_1\).
This gives \(\F_1\), with negative section \(Q_1\) and ruling
class \(H-Q_1\). The tangent direction of the conic specifies
\(r_4\) on \(Q_1\). Put
\(W=\Bl_{r_4}\F_1\). The strict transform of that ruling fiber
is \(D\). Contracting it is the elementary transformation of the
ruled surface at a point on its negative section. On \(\F_1\), we have
\(Q_1^2=-1\). Since \(r_4\) lies on \(Q_1\), its strict transform on
\(W\) has square \(-2\). It is disjoint from \(D\), so contracting \(D\)
does not change this self-intersection. The descended ruling therefore has
a section of square \(-2\), and the classification of ruled surfaces
identifies the result with \(\F_2\)
\cite[Chapter~V, Section~2]{Hartshorne1977}.

It remains to account for the marked centers. On \(W\), perform the blow-ups at \(p_2,p_3,p_4,p_5\),
then at \(r_1,r_2,r_3\). All these centers, including the two
directions \(r_2,r_3\) on the exceptional curve over \(p_2\),
lie away from the strict transform of \(D\): the tangent line
avoids the four points, and \(r_1\ne r_4\). The isomorphism
\[
W\setminus D\ \simeq\ \F_2\setminus\{t\},
\]
where \(t\) is the image of \(D\), transports these centers and
the points infinitely near to \(p_2\) to the corresponding blow-up of
\(\F_2\). Blow up the images of \(p_2,p_3,p_4,p_5\), followed by the
transported \(r_1,r_2,r_3\). The transported \(r_1\) lies on the negative
section of \(\F_2\). Thus the transported sequence gives the asserted
morphism \(\sigma_4:\widetilde Z_4\to\F_2\).
\end{proof}

Using precisely this order of seven centers in
\eqref{eq:explicit-orlov-objects}, we obtain
\[
\mathscr A_4=\left\langle G_7,\ldots,G_1,
\sigma_4^*\OO,\sigma_4^*\OO(f),
\sigma_4^*\OO(C_0+2f),\sigma_4^*\OO(C_0+3f)\right\rangle.
\]
Here \(C_0\) is the negative section on the new \(\F_2\).

\subsection{The five-point base}
\label{app:explicit-base-five}

For the five-point base, the four plane blow-ups followed by the five
vertex blow-ups of the pentagon on the degree-five del Pezzo surface give
a marked blow-up presentation of \(\widetilde Z_5\) over the plane and
hence its Orlov collection.
On the degree-five del Pezzo surface
\(S_5=\Bl_{p_1,\ldots,p_4}\Pp^2\), take the pentagon from
\cref{app:brauer-unobstructed}:
\[
\begin{aligned}
C_1&=Q_1,& C_2&=H-Q_1-Q_2,& C_3&=Q_2,\\
C_4&=H-Q_2-Q_3,& C_5&=H-Q_1-Q_4.
\end{aligned}
\]
Each curve has square \(-1\), adjacent curves intersect once, and
non-adjacent curves are disjoint. Blow up the five distinct vertices
\(r_i=C_i\cap C_{i+1}\), with exceptional curves \(R_i\).
Since exactly two vertices lie on each \(C_i\), the strict transforms
\[
E_i=C_i-R_{i-1}-R_i
\]
have square \(-1-2=-3\). Blowing up the unique intersection of any
adjacent pair separates it, so \(E_iE_j=0\) for \(i\ne j\).
By the \(5\)-gon construction of \cite[Section~3.5]{Corti2016},
their contraction belongs to \(X_{5,5/3}\). Uniqueness of this base
identifies the contraction with \(Z_5\), and consequently
\[
\widetilde Z_5\simeq
\Bl_{r_1,\ldots,r_5}\Bl_{p_1,\ldots,p_4}\Pp^2.
\]
Here the vertex blow-ups are on \(S_5\), not at nine distinct points
of the plane: \(r_1,r_5\) lie on \(Q_1\) over \(p_1\),
\(r_2,r_3\) lie on \(Q_2\) over \(p_2\), and \(r_4\) maps to a
fifth distinct plane point. After each vertex blow-up, the remaining
vertices are taken on the strict transforms of the indicated curves.
Order the centers as
\(p_1,\ldots,p_4,r_1,\ldots,r_5\), and write \(\sigma_5\)
for their composite. The twelve-object Orlov collection is
\[
\mathscr A_5=\left\langle G_9,\ldots,G_1,
\sigma_5^*\OO,\sigma_5^*\OO(H),\sigma_5^*\OO(2H)\right\rangle.
\]

\subsection{The six-point base}
\label{app:explicit-base-six}

For the six-point base, we contract six disjoint toric boundary curves of
\(\widetilde Z_6\) to obtain the degree-six del Pezzo surface, then two
opposite boundary curves to obtain \(\F_0\). Reversing these contractions
yields the twelve-object exceptional collection on the resolution,
despite the Brauer obstruction on the singular surface.
Use the fan and notation of \cref{app:brauer-six}. The resolved boundary is
\(V_1,E_1,\ldots,V_6,E_6\), with \(V_i^2=-1\) and \(E_i^2=-3\).
We first contract
\(V_1,\ldots,V_6\), in that order. At every stage, the neighbors
of the ray \(v_i\) being removed are still \(u_{i-1},u_i\), and
\[
u_{i-1}+u_i=v_i,\qquad \det(u_{i-1},u_i)=1.
\]
Thus every contraction is a smooth toric blow-down. Equivalently, the
\(V_i\) are pairwise disjoint, so their self-intersections remain
\(-1\) until they are contracted. After any initial set \(I\) of
these contractions, the self-intersection of the image of \(E_j\) is
\[
-3+\#\bigl(I\cap\{j,j+1\}\bigr),
\]
with cyclic indices. After all six, the fan consists of the \(u_i\)
and every boundary curve has square \(-1\). This is the smooth toric
del Pezzo surface of degree six: its anticanonical divisor intersects
each invariant curve positively, and its square is six.

Next contract the opposite boundary curves with rays \(u_1\) and
\(u_4\). They are disjoint, and their ray relations are
\[
u_1=u_6+u_2,\qquad u_4=u_3+u_5,
\qquad \det(u_6,u_2)=\det(u_3,u_5)=1.
\]
Both are therefore \((-1)\)-curves at the moment of contraction.
The remaining rays are \(u_2,u_3,-u_2,-u_3\), and
\(\det(u_2,u_3)=1\). In this lattice basis the fan is the square
fan of \(\F_0=\Pp^1\times\Pp^1\).

Reversing the sequence prescribes eight toric point blow-ups: first
insert \(u_4\) between \(u_3,u_5\), then \(u_1\) between
\(u_6,u_2\), and then insert \(v_6,v_5,\ldots,v_1\), each
between \(u_{i-1},u_i\). The insertions of \(u_4,u_1,v_6,v_3\)
blow up the four original fixed points of \(\F_0\); the remaining
four blow-ups are at points infinitely near to those producing \(u_4\)
and \(u_1\). The insertions of \(v_5,v_4\) have centers on the successive
strict transforms of the curve with ray \(u_4\), and those of \(v_2,v_1\)
on the successive strict transforms of the curve with ray \(u_1\).
Let
\(\sigma_6:\widetilde Z_6\to\F_0\) be their composite. With
\(f_1,f_2\) the two ruling classes on \(\F_0\), the resulting
twelve-object collection is
\[
\mathscr A_6=\left\langle G_8,\ldots,G_1,
\sigma_6^*\OO,\sigma_6^*\OO(f_1),\sigma_6^*\OO(f_2),
\sigma_6^*\OO(f_1+f_2)\right\rangle.
\]

The marked blow-up presentations in
\cref{app:explicit-base-one,app:explicit-base-two,app:explicit-base-three,app:explicit-base-four,app:explicit-base-five,app:explicit-base-six}
establish \cref{prop:explicit-rigid-bases}. The following corollary extends
the resulting exceptional collections along the blow-ups at distinct smooth
points of \(Z_k\) described in \cref{sec:cascade}, proving
\cref{thm:explicit-cascades-main}. We use \(A^{(k)}_j\) for the entries
of \(\mathscr A_k\) on \(\widetilde Z_k\), reserving \(G_i\) for
blow-up objects.

\begin{corollary}\label{cor:explicit-cascade-collections}
Let \(X\) belong to a Fano index one family \(X_{k,d}\), and put
\(s=K_{Z_k}^2-K_X^2\). Choose a marked cascade presentation by
\(s\) distinct smooth points of \(Z_k\), and let
\(\gamma:\widetilde X\to\widetilde Z_k\) be the induced map of
minimal resolutions. Write the exceptional collection on \(\widetilde Z_k\) as
\[
\mathscr A_k=\langle A^{(k)}_1,\ldots,A^{(k)}_{N_k}\rangle,
\]
with \(N_k\) as in Table~\ref{tab:rigid-base-invariants}, and let \(D_1,\ldots,D_s\) be the disjoint exceptional curves of
\(\gamma\). Set \(H_i=\OO_{D_i}(-1)\); the distinctness of the centers follows
from \eqref{eq:cascade-distinct}. The derived-pullback convention
\eqref{eq:explicit-orlov-objects} is needed only for the collections on
\(\widetilde Z_k\), whose marked blow-up presentations may use points
infinitely near to earlier points as blow-up centers. Then
\[
\Db(\coh\widetilde X)=
\left\langle H_s,\ldots,H_1,
L\gamma^*A^{(k)}_1,\ldots,L\gamma^*A^{(k)}_{N_k}\right\rangle
\]
is a full exceptional collection. If \(\XX\) is the canonical stack
and \(\Phi\) its Ishii--Ueda embedding, the ordered collection
\[
\Db(\coh\XX)=\left\langle
\mathcal E_{p_1},\ldots,\mathcal E_{p_k},
\Phi(H_s),\ldots,\Phi(H_1),
\Phi\bigl(L\gamma^*\mathscr A_k\bigr)
\right\rangle
\]
is full exceptional, of length
\(N_k+s+k=12-K_X^2+4k/3\), the specialization of
\cref{cor:mixed-basket-length} with \(n_i=3\) and \(r=k\).
\end{corollary}

\begin{remark}[Marked presentations and dominance]
\label{rem:presentations-versus-dominance}
Every individual member admits the blow-up presentation of
\cite[Corollary~8(1)--(6)]{Corti2016}, with distinct smooth centers by
\eqref{eq:cascade-distinct}. Applying Orlov's formula and the Ishii--Ueda
embedding gives all the collections in \cref{thm:explicit-cascades-main}.
Appendix~\ref{app:general-X10} compares the parameter space of ordered
eight-point configurations on \(\Pp(1,1,3)\) with the parameter space of
degree-\(10\) hypersurfaces in \(\Pp(1,2,3,5)\). Its automorphism and
fiber calculations show that general hypersurfaces admit general centers,
and general centers give general hypersurfaces.

For other families, an analogous comparison with a specified equation
parameter space or an alternative marked blow-up presentation requires its own
classifying map and fiber analysis.
The infinitely near points used as blow-up centers in
Appendix~\ref{app:explicit-bases} occur only in the fixed marked blow-up
presentations of \(\widetilde Z_k\), not in the additional blow-ups at
distinct smooth points of \(Z_k\).
\end{remark}


% ------------------------------------------------------------

\section*{Acknowledgements}
I am grateful to my advisor, Sergey Galkin, for his guidance, support, and many valuable discussions. I thank Nicolau Saldanha and Carlos Tomei for supporting my travel to the RGAS School 2026, where part of this work was further developed. I also thank Alexander Kuznetsov for a helpful conversation during the school and for pointing me to the work of Karmazyn--Kuznetsov--Shinder, which motivated the comparison with the singular surface in Section~4. I am grateful to Alexey Elagin for useful discussions at IMPA about derived categories of singular varieties and for emphasizing the distinction between the categorical behavior of singular surfaces and that of their canonical stacks. Finally, I thank Alessio Corti for helpful correspondence on the Corti--Heuberger cascades, and in particular for pointing out the relevance of the Reid--Suzuki construction to \(X_{10}\).

\section*{Funding}
The author was supported by the Funda\c{c}\~ao Carlos Chagas Filho de Amparo \`a Pesquisa do Estado do Rio de Janeiro (FAPERJ) through the Nota 10 Excellence Scholarship.

\section*{Declaration of generative AI and AI-assisted technologies in the writing process}
During the preparation of this work, the author used ChatGPT (OpenAI) to assist with language editing, organization, and manuscript formatting. After using this tool, the author reviewed and edited the content as needed and takes full responsibility for the content of the publication.

\printbibliography

\end{document}
